% 由逐条人工译稿生成；保留原索引全部条目与层级，页码指向译本。
\clearpage
\begingroup
\makeatletter
\@nameuse{imki@set@main}
\def\indexname{索引}
\def\imki@maybeaddtotoc{\phantomsection\addcontentsline{toc}{chapter}{索引}}
\makeatother
\setlength{\multicolsep}{8pt}
\begin{theindex}
\raggedcolumns
\renewcommand{\indexspace}{\par\vskip9.5pt plus5pt minus3pt\relax}
\item \leavevmode\hypertarget{source:476}{}\hypertarget{idx:0038}{$\mathcal A$ 超几何方程}，\textit{参见} \hyperlink{idx:0055}{$\mathcal A$ 方程组}
\indexspace
\item \leavevmode\hypertarget{idx:0055}{$\mathcal A$ 方程组},\enspace\IndexPageList[0055]{113}，\textit{另见} \hyperlink{idx:0449}{GKZ 方程组}；\hyperlink{idx:0615}{超几何方程}
\subitem \leavevmode\hypertarget{idx:0056}{定义},\enspace\IndexPageList[0056]{114}
\subitem \leavevmode\hypertarget{idx:0057}{例题},\enspace\IndexPageList[0057]{116,117,118,119,148,185}
\subitem \leavevmode\hypertarget{idx:0058}{在最大幂幺边界点处的全纯解},\enspace\IndexPageList[0058]{183}
\subitem \leavevmode\hypertarget{idx:0059}{Calabi--Yau 环面超曲面的方程组}，\textit{参见} \hyperlink{idx:0099}{Calabi--Yau，Calabi--Yau 环面超曲面}
\subitem \leavevmode\hypertarget{idx:0060}{光滑 Fano 环面簇的方程组},\enspace\IndexPageList[0060]{413,414}
\subitem \leavevmode\hypertarget{idx:0061}{解},\enspace\IndexPageList[0061]{121,123,187}
\indexspace
\item \leavevmode\hypertarget{idx:0002}{A 模型},\enspace\IndexPageList[0002]{30,445,446}
\subitem \leavevmode\hypertarget{idx:0003}{联络},\enspace\IndexPageList[0003]{105,134,281,282,288,289,290,294}
\subsubitem \leavevmode\hypertarget{idx:0004}{渐近性态},\enspace\IndexPageList[0004]{268}
\subsubitem \leavevmode\hypertarget{idx:0005}{典范延拓},\enspace\IndexPageList[0005]{272,274,277,339}
\subsubitem \leavevmode\hypertarget{idx:0006}{收敛性假设},\enspace\IndexPageList[0006]{268,272,284,339}
\subsubitem \leavevmode\hypertarget{idx:0007}{定义},\enspace\IndexPageList[0007]{267}
\subsubitem \leavevmode\hypertarget{idx:0008}{平坦截面},\enspace\IndexPageList[0008]{271,272,338,339,340,341}
\subsubitem \leavevmode\hypertarget{idx:0009}{平坦性},\enspace\IndexPageList[0009]{267}
\subsubitem \leavevmode\hypertarget{idx:0010}{具有最大幂幺单值性},\enspace\IndexPageList[0010]{273}
\subsubitem \leavevmode\hypertarget{idx:0011}{单值性},\enspace\IndexPageList[0011]{269,271,272,339}
\subsubitem \leavevmode\hypertarget{idx:0012}{单值性权滤过},\enspace\IndexPageList[0012]{272,273,276}，\textit{另见} \hyperlink{idx:0907}{单值性权滤过}
\subsubitem \leavevmode\hypertarget{idx:0013}{Calabi--Yau 三维簇的联络},\enspace\IndexPageList[0013]{271,278,279,280,341,351,352}
\subsubitem \leavevmode\hypertarget{idx:0014}{对应的 Picard--Fuchs 算子}，\textit{参见} \hyperlink{idx:1062}{Picard--Fuchs 算子}
\subsubitem \leavevmode\hypertarget{idx:0015}{基 $T_0,\ldots,T_m$ 的作用},\enspace\IndexPageList[0015]{274}
\subsubitem \leavevmode\hypertarget{idx:0016}{满足整性猜想},\enspace\IndexPageList[0016]{273,281}
\subitem \leavevmode\hypertarget{idx:0017}{关联函数}，\textit{参见} \hyperlink{idx:0184}{关联函数，A 模型关联函数}
\subitem \leavevmode\hypertarget{idx:0018}{是 $(1+1)$ 维拓扑量子场论},\enspace\IndexPageList[0018]{457}
\indexspace
\item \leavevmode\hypertarget{idx:0037}{$\mathcal A$ 判别式},\enspace\IndexPageList[0037]{145,148}，\textit{另见} \hyperlink{idx:0214}{判别式轨迹}；\hyperlink{idx:1093}{主 $\mathcal A$ 行列式}
\indexspace
\item \leavevmode\hypertarget{idx:0019}{A 型 Hodge 结构变分},\enspace\IndexPageList[0019]{48,130,245,276,277}，\textit{另见} \hyperlink{idx:0003}{A 模型，联络}
\subitem \leavevmode\hypertarget{idx:0020}{定义},\enspace\IndexPageList[0020]{274,278}
\subitem \leavevmode\hypertarget{idx:0021}{Griffiths 横截性},\enspace\IndexPageList[0021]{275,277}
\subitem \leavevmode\hypertarget{idx:0022}{Hodge 滤过},\enspace\IndexPageList[0022]{274}
\subitem \leavevmode\hypertarget{idx:0023}{整数格},\enspace\IndexPageList[0023]{272,274,277,283}
\subitem \leavevmode\hypertarget{idx:0024}{中间 A 型 Hodge 结构变分},\enspace\IndexPageList[0024]{281}
\subsubitem \leavevmode\hypertarget{idx:0025}{定义},\enspace\IndexPageList[0025]{278}
\subsubitem \leavevmode\hypertarget{idx:0026}{具有最大幂幺单值性},\enspace\IndexPageList[0026]{278}
\subsubitem \leavevmode\hypertarget{idx:0027}{混合 Hodge 结构为 Hodge--Tate 型},\enspace\IndexPageList[0027]{278,287}
\subsubitem \leavevmode\hypertarget{idx:0028}{实结构},\enspace\IndexPageList[0028]{278,287}
\subsubitem \leavevmode\hypertarget{idx:0029}{环面情形},\enspace\IndexPageList[0029]{294}
\subitem \leavevmode\hypertarget{idx:0030}{混合 Hodge 结构},\enspace\IndexPageList[0030]{276,278}，\textit{另见} \hyperlink{idx:0012}{A 模型，联络，单值性权滤过}
\subitem \leavevmode\hypertarget{idx:0031}{极化},\enspace\IndexPageList[0031]{275,277,281}
\indexspace
\item \leavevmode\hypertarget{idx:0001}{$A_\infty$ 范畴},\enspace\IndexPageList[0001]{426}
\indexspace
\item \leavevmode\hypertarget{idx:0032}{Abel--Jacobi 映射},\enspace\IndexPageList[0032]{90,450}
\indexspace
\item \leavevmode\hypertarget{idx:0054}{Arnold 奇异对偶性},\enspace\IndexPageList[0054]{424}
\indexspace
\item \leavevmode\hypertarget{idx:0062}{B 模型},\enspace\IndexPageList[0062]{31,445,446}
\subitem \leavevmode\hypertarget{idx:0063}{关联函数}，\textit{参见} \hyperlink{idx:0189}{关联函数，B 模型关联函数}
\subitem \leavevmode\hypertarget{idx:0064}{势函数},\enspace\IndexPageList[0064]{131,132,284,286}
\subsubitem \leavevmode\hypertarget{idx:0065}{用周期表示},\enspace\IndexPageList[0065]{412,422}
\subsubitem \leavevmode\hypertarget{idx:0066}{五次镜像的势函数}，\textit{参见} \hyperlink{idx:1168}{五次镜像，B 模型势函数}
\subitem \leavevmode\hypertarget{idx:0067}{环面情形},\enspace\IndexPageList[0067]{292,411}
\indexspace
\item \leavevmode\hypertarget{source:477}{}\hypertarget{idx:0068}{Batyrev 镜像}，\textit{参见} \hyperlink{idx:0799}{镜像构造，Batyrev 构造}
\indexspace
\item \leavevmode\hypertarget{idx:0069}{Batyrev 量子环},\enspace\IndexPageList[0069]{246,248,412}
\subitem \leavevmode\hypertarget{idx:0070}{定义},\enspace\IndexPageList[0070]{247}
\subitem \leavevmode\hypertarget{idx:0071}{$-K_X$ 为 nef 的情形},\enspace\IndexPageList[0071]{416}
\subitem \leavevmode\hypertarget{idx:0072}{$\mathbb F_2$ 的情形},\enspace\IndexPageList[0072]{416,417,418}
\subitem \leavevmode\hypertarget{idx:0073}{$\mathbb P^n$ 的情形},\enspace\IndexPageList[0073]{247}
\subitem \leavevmode\hypertarget{idx:0074}{Fano 环面簇的情形},\enspace\IndexPageList[0074]{415}
\indexspace
\item \leavevmode\hypertarget{idx:0085}{BRST 上同调},\enspace\IndexPageList[0085]{325,445,446}
\indexspace
\item \leavevmode\hypertarget{idx:0086}{BRST 算子},\enspace\IndexPageList[0086]{445,446}
\indexspace
\item \leavevmode\hypertarget{idx:0087}{Calabi--Yau}
\subitem \leavevmode\hypertarget{idx:0088}{自同构群},\enspace\IndexPageList[0088]{35,150,152,153,154,161}
\subitem \leavevmode\hypertarget{idx:0089}{由 K3 曲面构造},\enspace\IndexPageList[0089]{89}
\subitem \leavevmode\hypertarget{idx:0090}{广义 Calabi--Yau},\enspace\IndexPageList[0090]{88}
\subitem \leavevmode\hypertarget{idx:0091}{Calabi--Yau 流形},\enspace\IndexPageList[0091]{24,33,208,241,245,266,269,287,288,289,290,338,340,382,404,444,450}
\subsubitem \leavevmode\hypertarget{idx:0092}{具有 K3 或椭圆纤维化},\enspace\IndexPageList[0092]{427}
\subitem \leavevmode\hypertarget{idx:0093}{极小 Calabi--Yau},\enspace\IndexPageList[0093]{290}
\subitem \leavevmode\hypertarget{idx:0094}{模}，\textit{参见} \hyperlink{idx:0870}{模}
\subitem \leavevmode\hypertarget{idx:0095}{Calabi--Yau 轨形},\enspace\IndexPageList[0095]{35,40,78,431,433}
\subitem \leavevmode\hypertarget{idx:0096}{刚性 Calabi--Yau},\enspace\IndexPageList[0096]{88,284,290}
\subitem \leavevmode\hypertarget{idx:0097}{Calabi--Yau 三维簇},\enspace\IndexPageList[0097]{26,27,29,30,32,33,76,97,102,125,126,127,128,129,130,131,132,133,134,135,153,155,156,157,166,174,224,225,226,227,228,229,230,231,232,233,248,249,250,261,270,271,274,278,279,280,281,282,283,284,285,290,305,311,312,313,314,315,316,331,341,348,349,350,351,352,375,389,411,412,428,446,448,450}，\textit{另见} \hyperlink{idx:0184}{关联函数，A 模型关联函数}；\hyperlink{idx:0870}{模}；\hyperlink{idx:1049}{Picard--Fuchs 方程}；\hyperlink{idx:1444}{Yukawa 耦合}
\subitem \leavevmode\hypertarget{idx:0098}{Calabi--Yau 环面完全交}，\textit{参见} \hyperlink{idx:1348}{环面完全交，Calabi--Yau 环面完全交}
\subitem \leavevmode\hypertarget{idx:0099}{Calabi--Yau 环面超曲面},\enspace\IndexPageList[0099]{76,79,83,139,140,141,142,143,144,145,146,147,148,149,150,157,158,159,160,161,162,163,164,165,166,178,179,180,181,182,183,184,185,186,187,188,189,290,293,406,453}
\subsubitem \leavevmode\hypertarget{idx:0100}{$\mathcal A$ 方程组},\enspace\IndexPageList[0100]{114,115,116}
\subsubitem \leavevmode\hypertarget{idx:0101}{复模}，\textit{参见} \hyperlink{idx:0871}{模，复模}
\subsubitem \leavevmode\hypertarget{idx:0102}{Kähler 模}，\textit{参见} \hyperlink{idx:0882}{模，Kähler 模}
\subsubitem \leavevmode\hypertarget{idx:0103}{非退化情形},\enspace\IndexPageList[0103]{145}
\subsubitem \leavevmode\hypertarget{idx:0104}{Picard--Fuchs 方程},\enspace\IndexPageList[0104]{114,115,406}
\subitem \leavevmode\hypertarget{idx:0105}{Calabi--Yau 簇},\enspace\IndexPageList[0105]{78,83,86,430}
\subsubitem \leavevmode\hypertarget{idx:0106}{定义},\enspace\IndexPageList[0106]{34}
\subsubitem \leavevmode\hypertarget{idx:0107}{极小 Calabi--Yau 簇},\enspace\IndexPageList[0107]{34,78,86,139,430,432,433}
\subitem \leavevmode\hypertarget{idx:0108}{加权射影 Calabi--Yau 超曲面},\enspace\IndexPageList[0108]{93,230,428}
\subsubitem \leavevmode\hypertarget{idx:0109}{7555 个超曲面的列表},\enspace\IndexPageList[0109]{83,95,168}
\indexspace
\item \leavevmode\hypertarget{idx:0110}{Calabi--Yau 锥},\enspace\IndexPageList[0110]{162}
\subitem \leavevmode\hypertarget{idx:0111}{例题},\enspace\IndexPageList[0111]{163}
\indexspace
\item \leavevmode\hypertarget{idx:0115}{Cauchy--Riemann 方程},\enspace\IndexPageList[0115]{208,210,214}
\subitem \leavevmode\hypertarget{idx:0116}{非齐次方程},\enspace\IndexPageList[0116]{210}
\subitem \leavevmode\hypertarget{idx:0117}{扰动方程},\enspace\IndexPageList[0117]{210}
\indexspace
\item \leavevmode\hypertarget{idx:0118}{Cayley 技巧},\enspace\IndexPageList[0118]{87}
\indexspace
\item \leavevmode\hypertarget{idx:0122}{Chow 商},\enspace\IndexPageList[0122]{140,142,143,171}
\indexspace
\item \leavevmode\hypertarget{idx:0124}{Clemens 猜想},\enspace\IndexPageList[0124]{38,39,225,229,237,423}
\subitem \leavevmode\hypertarget{idx:0125}{陈述},\enspace\IndexPageList[0125]{225}
\indexspace
\item \leavevmode\hypertarget{source:478}{}\hypertarget{idx:0199}{$\operatorname{cpl}(\Sigma)$},\enspace\IndexPageList[0199]{64,65,66,67,68,158,159,160,162,184}
\subitem \leavevmode\hypertarget{idx:0200}{在 $A^+(\Sigma)$ 中的定义},\enspace\IndexPageList[0200]{66}
\subitem \leavevmode\hypertarget{idx:0201}{在 $A^+_{n-1}(X_\Sigma)\otimes\mathbb R$ 中的定义},\enspace\IndexPageList[0201]{62}
\subitem \leavevmode\hypertarget{idx:0202}{$\Sigma$ 凸性},\enspace\IndexPageList[0202]{66}
\indexspace
\item \leavevmode\hypertarget{idx:0207}{D 膜},\enspace\IndexPageList[0207]{425,427}
\indexspace
\item \leavevmode\hypertarget{idx:0226}{Dubrovin 联络},\enspace\IndexPageList[0226]{262,334,335}
\subitem \leavevmode\hypertarget{idx:0227}{定义},\enspace\IndexPageList[0227]{262}
\subitem \leavevmode\hypertarget{idx:0228}{Euler 向量场},\enspace\IndexPageList[0228]{264}
\subitem \leavevmode\hypertarget{idx:0229}{平坦性},\enspace\IndexPageList[0229]{263,267}
\subitem \leavevmode\hypertarget{idx:0230}{形式联络},\enspace\IndexPageList[0230]{265}
\subitem \leavevmode\hypertarget{idx:0231}{齐次性},\enspace\IndexPageList[0231]{264}
\subitem \leavevmode\hypertarget{idx:0232}{势函数},\enspace\IndexPageList[0232]{262}
\subitem \leavevmode\hypertarget{idx:0233}{挠率},\enspace\IndexPageList[0233]{263}
\subitem \leavevmode\hypertarget{idx:0234}{WDVV 方程},\enspace\IndexPageList[0234]{263,264}
\indexspace
\item \leavevmode\hypertarget{idx:0235}{Dubrovin 形式体系},\enspace\IndexPageList[0235]{257,262,263,264,265,267,327}，\textit{另见} \hyperlink{idx:0126}{上同调场论}
\subitem \leavevmode\hypertarget{idx:0236}{二元乘积},\enspace\IndexPageList[0236]{262}
\subitem \leavevmode\hypertarget{idx:0237}{联络}，\textit{参见} \hyperlink{idx:0226}{Dubrovin 联络}
\indexspace
\item \leavevmode\hypertarget{source:479}{}\hypertarget{idx:0271}{Euler $\Gamma$ 函数},\enspace\IndexPageList[0271]{185,186,188}
\indexspace
\item \leavevmode\hypertarget{idx:0273}{Euler 常数},\enspace\IndexPageList[0273]{185}
\indexspace
\item \leavevmode\hypertarget{idx:0272}{Euler 积分},\enspace\IndexPageList[0272]{183}，\textit{另见} \hyperlink{idx:0449}{GKZ 方程组}
\indexspace
\item \leavevmode\hypertarget{idx:0257}{Euler 类},\enspace\IndexPageList[0257]{196,312,374,381,383}
\subitem \leavevmode\hypertarget{idx:0258}{等变 Euler 类},\enspace\IndexPageList[0258]{299,315,359,361,364,365,367,400}
\subsubitem \leavevmode\hypertarget{idx:0259}{$N_\Gamma$ 的 Euler 类},\enspace\IndexPageList[0259]{308,310,311,312,314,316,393}
\subitem \leavevmode\hypertarget{idx:0260}{零丛的 Euler 类},\enspace\IndexPageList[0260]{381,388,395}
\subitem \leavevmode\hypertarget{idx:0261}{代数叠与粗模空间之间的比较},\enspace\IndexPageList[0261]{204,207,325,355,360}
\indexspace
\item \leavevmode\hypertarget{idx:0256}{Euler 示性数},\enspace\IndexPageList[0256]{333}
\indexspace
\item \leavevmode\hypertarget{idx:0262}{Euler 数据},\enspace\IndexPageList[0262]{50,51,368,371,372,373,403}
\subitem \leavevmode\hypertarget{idx:0263}{定义},\enspace\IndexPageList[0263]{362}
\subitem \leavevmode\hypertarget{idx:0264}{例题}，\textit{参见} \hyperlink{idx:1043}{$\widehat P$}；\hyperlink{idx:1101}{$\widehat Q$}
\subitem \leavevmode\hypertarget{idx:0265}{超几何函数},\enspace\IndexPageList[0265]{369}，\textit{参见} \hyperlink{idx:0584}{$HG$，Euler 数据的 $HG$}
\subitem \leavevmode\hypertarget{idx:0266}{相关联的 Euler 数据},\enspace\IndexPageList[0266]{366,368,376,378}
\subsubitem \leavevmode\hypertarget{idx:0267}{定义},\enspace\IndexPageList[0267]{365}
\subsubitem \leavevmode\hypertarget{idx:0268}{镜像变换},\enspace\IndexPageList[0268]{371,375,376,404}，\textit{另见} \hyperlink{idx:0581}{$HG$，镜像变换}
\subitem \leavevmode\hypertarget{idx:0269}{多重覆盖解释},\enspace\IndexPageList[0269]{366}
\subitem \leavevmode\hypertarget{idx:0270}{唯一性定理},\enspace\IndexPageList[0270]{367,371,376,378,403}
\indexspace
\item \leavevmode\hypertarget{source:487}{}\hypertarget{idx:0966}{$\Omega$--Euler 数据}，\textit{参见} \hyperlink{idx:0262}{Euler 数据}
\indexspace
\item \leavevmode\hypertarget{idx:0274}{Euler--Lagrange 方程},\enspace\IndexPageList[0274]{435,436,437,441,444,445,447}
\subitem \leavevmode\hypertarget{idx:0275}{电磁场的方程},\enspace\IndexPageList[0275]{440}
\indexspace
\item \leavevmode\hypertarget{idx:0295}{F 理论},\enspace\IndexPageList[0295]{427}
\indexspace
\item \leavevmode\hypertarget{idx:0308}{Fano 簇},\enspace\IndexPageList[0308]{69,241,242,246,251,428}
\indexspace
\item \leavevmode\hypertarget{idx:0315}{Feynman 积分},\enspace\IndexPageList[0315]{25,30}，\textit{另见} \hyperlink{idx:1037}{路径积分}
\indexspace
\item \leavevmode\hypertarget{idx:0314}{Feynman 图},\enspace\IndexPageList[0314]{219}
\indexspace
\item \leavevmode\hypertarget{idx:0336}{Fourier 变换},\enspace\IndexPageList[0336]{437,441}
\indexspace
\item \leavevmode\hypertarget{idx:0337}{Frobenius 代数},\enspace\IndexPageList[0337]{244,257,263,453,458}
\subitem \leavevmode\hypertarget{idx:0338}{定义},\enspace\IndexPageList[0338]{455}
\subitem \leavevmode\hypertarget{idx:0339}{给出 $(1+1)$ 维拓扑量子场论},\enspace\IndexPageList[0339]{455,457}
\subitem \leavevmode\hypertarget{idx:0340}{$n$ 点函数},\enspace\IndexPageList[0340]{455,457}
\subitem \leavevmode\hypertarget{source:480}{}\hypertarget{idx:0341}{三点函数},\enspace\IndexPageList[0341]{455}
\subitem \leavevmode\hypertarget{idx:0342}{迹映射},\enspace\IndexPageList[0342]{455}
\indexspace
\item \leavevmode\hypertarget{idx:0344}{Frobenius 方法},\enspace\IndexPageList[0344]{49,100,184,185,186,356}
\subitem \leavevmode\hypertarget{idx:0345}{例题},\enspace\IndexPageList[0345]{185,188}
\subitem \leavevmode\hypertarget{idx:0346}{在 $\mathbb P^4$ 的上同调环中},\enspace\IndexPageList[0346]{187}
\indexspace
\item \leavevmode\hypertarget{idx:0343}{Frobenius 流形},\enspace\IndexPageList[0343]{264}
\indexspace
\item \leavevmode\hypertarget{idx:0348}{Gale 变换},\enspace\IndexPageList[0348]{67}
\indexspace
\item \leavevmode\hypertarget{idx:0365}{Gauss--Manin 联络},\enspace\IndexPageList[0365]{31,48,98,105,174,273,282,288,289,290}
\subitem \leavevmode\hypertarget{idx:0366}{联络矩阵}，\textit{参见} \hyperlink{idx:0174}{联络矩阵}
\subitem \leavevmode\hypertarget{idx:0367}{定义},\enspace\IndexPageList[0367]{97}
\subitem \leavevmode\hypertarget{idx:0368}{平坦整截面},\enspace\IndexPageList[0368]{105,129,134,173,174,410}
\subitem \leavevmode\hypertarget{idx:0369}{Calabi--Yau 三维簇的 Gauss--Manin 联络},\enspace\IndexPageList[0369]{129,131,132,274,280,281}
\subitem \leavevmode\hypertarget{idx:0370}{与归一化 Yukawa 耦合的关系},\enspace\IndexPageList[0370]{128}
\indexspace
\item \leavevmode\hypertarget{idx:0381}{GIT 商},\enspace\IndexPageList[0381]{140,143,145}
\indexspace
\item \leavevmode\hypertarget{idx:0382}{Givental $I$ 函数}
\subitem \leavevmode\hypertarget{idx:0383}{$I_{\mathcal V}$},\enspace\IndexPageList[0383]{49,50,123,381,402,403}
\subsubitem \leavevmode\hypertarget{idx:0384}{定义},\enspace\IndexPageList[0384]{379}
\subsubitem \leavevmode\hypertarget{idx:0385}{例题},\enspace\IndexPageList[0385]{387}
\subsubitem \leavevmode\hypertarget{idx:0386}{凸簇的 $I_{\mathcal V}$},\enspace\IndexPageList[0386]{404}
\subsubitem \leavevmode\hypertarget{idx:0387}{环面簇的 $I_{\mathcal V}$},\enspace\IndexPageList[0387]{404,406}
\subsubitem \leavevmode\hypertarget{idx:0388}{$\mathbb P(1,1,2,2,2)$ 中超曲面的 $I_{\mathcal V}$},\enspace\IndexPageList[0388]{407}
\subsubitem \leavevmode\hypertarget{idx:0389}{五次三维簇的 $I_{\mathcal V}$},\enspace\IndexPageList[0389]{49,121,124,187,384,385}
\subsubitem \leavevmode\hypertarget{idx:0390}{满足镜像的 Picard--Fuchs 方程},\enspace\IndexPageList[0390]{124,144,151,406,408,422}
\subitem \leavevmode\hypertarget{idx:0391}{$\overline I_{\mathcal V}$},\enspace\IndexPageList[0391]{388,389}
\subitem \leavevmode\hypertarget{idx:0392}{$\widetilde I$},\enspace\IndexPageList[0392]{121,122,123,124,151,187,382}
\subsubitem \leavevmode\hypertarget{idx:0393}{满足 $\mathcal A$ 方程组},\enspace\IndexPageList[0393]{121,414}
\subitem \leavevmode\hypertarget{idx:0394}{$I_{\mathbb P^n}$ 的等变版本},\enspace\IndexPageList[0394]{394}
\subitem \leavevmode\hypertarget{idx:0395}{$I_{\mathcal V}$ 的等变版本},\enspace\IndexPageList[0395]{402,405}
\subitem \leavevmode\hypertarget{idx:0396}{$\mathbb F_2$ 的 $I$ 函数},\enspace\IndexPageList[0396]{416,417}
\subitem \leavevmode\hypertarget{idx:0397}{$\mathbb P^1$ 的 $I$ 函数},\enspace\IndexPageList[0397]{380}
\subitem \leavevmode\hypertarget{idx:0398}{$\mathbb P^n$ 的 $I$ 函数},\enspace\IndexPageList[0398]{380,388,394}
\subitem \leavevmode\hypertarget{idx:0399}{环面簇的 $I$ 函数},\enspace\IndexPageList[0399]{412}
\subsubitem \leavevmode\hypertarget{idx:0400}{当 $-K_X$ 为 nef 时},\enspace\IndexPageList[0400]{416}
\indexspace
\item \leavevmode\hypertarget{idx:0401}{Givental $J$ 函数},\enspace\IndexPageList[0401]{48,343,344,345,346,347,348,381,382}
\subitem \leavevmode\hypertarget{idx:0402}{$J_{\mathcal V}$},\enspace\IndexPageList[0402]{49,381,382,402,403}
\subsubitem \leavevmode\hypertarget{idx:0403}{定义},\enspace\IndexPageList[0403]{381}
\subsubitem \leavevmode\hypertarget{idx:0404}{例题},\enspace\IndexPageList[0404]{387}
\subsubitem \leavevmode\hypertarget{idx:0405}{凸簇的 $J_{\mathcal V}$},\enspace\IndexPageList[0405]{388}
\subsubitem \leavevmode\hypertarget{idx:0406}{环面簇的 $J_{\mathcal V}$},\enspace\IndexPageList[0406]{404,405,406}
\subsubitem \leavevmode\hypertarget{idx:0407}{$\mathbb P(1,1,2,2,2)$ 中超曲面的 $J_{\mathcal V}$},\enspace\IndexPageList[0407]{407,411}
\subsubitem \leavevmode\hypertarget{idx:0408}{五次三维簇的 $J_{\mathcal V}$},\enspace\IndexPageList[0408]{49,385}
\subsubitem \leavevmode\hypertarget{idx:0409}{给出修正量子上同调中的关系},\enspace\IndexPageList[0409]{386}
\subsubitem \leavevmode\hypertarget{idx:0410}{给出量子上同调中的关系},\enspace\IndexPageList[0410]{387}
\subsubitem \leavevmode\hypertarget{idx:0411}{符号记法},\enspace\IndexPageList[0411]{381}
\subitem \leavevmode\hypertarget{idx:0412}{定义},\enspace\IndexPageList[0412]{343}
\subitem \leavevmode\hypertarget{idx:0413}{$J_{\mathbb P^n}$ 的等变版本},\enspace\IndexPageList[0413]{389}
\subsubitem \leavevmode\hypertarget{idx:0414}{$S$ 的定义},\enspace\IndexPageList[0414]{389}
\subsubitem \leavevmode\hypertarget{idx:0415}{$Z_i$ 的定义},\enspace\IndexPageList[0415]{390}
\subitem \leavevmode\hypertarget{idx:0416}{$J_{\mathcal V}$ 的等变版本},\enspace\IndexPageList[0416]{395,405}
\subsubitem \leavevmode\hypertarget{idx:0417}{$\widehat S_{\mathcal V}$ 的定义},\enspace\IndexPageList[0417]{395,396}
\subsubitem \leavevmode\hypertarget{idx:0418}{$S_{\mathcal V}$ 的定义},\enspace\IndexPageList[0418]{395,396}
\subsubitem \leavevmode\hypertarget{idx:0419}{$Z_{i,\mathcal V}$ 的定义},\enspace\IndexPageList[0419]{396}
\subitem \leavevmode\hypertarget{idx:0420}{$\mathbb F_2$ 的 $J$ 函数},\enspace\IndexPageList[0420]{416,417}
\subitem \leavevmode\hypertarget{idx:0421}{$\mathbb P^1$ 的 $J$ 函数},\enspace\IndexPageList[0421]{347,380}
\subitem \leavevmode\hypertarget{idx:0422}{$\mathbb P^n$ 的 $J$ 函数},\enspace\IndexPageList[0422]{347,348,369,380,388}
\subitem \leavevmode\hypertarget{idx:0423}{Calabi--Yau 三维簇的 $J$ 函数},\enspace\IndexPageList[0423]{348,349,350,351,352}
\subsubitem \leavevmode\hypertarget{idx:0424}{公式},\enspace\IndexPageList[0424]{348,349}
\subsubitem \leavevmode\hypertarget{idx:0425}{用 Gromov--Witten 势表示},\enspace\IndexPageList[0425]{349}
\subsubitem \leavevmode\hypertarget{idx:0426}{二阶偏导数},\enspace\IndexPageList[0426]{349,386}
\subitem \leavevmode\hypertarget{idx:0427}{环面簇的 $J$ 函数},\enspace\IndexPageList[0427]{412}
\subsubitem \leavevmode\hypertarget{idx:0428}{当 $-K_X$ 为 nef 时},\enspace\IndexPageList[0428]{416}
\subitem \leavevmode\hypertarget{idx:0429}{$\mathbb P(1,1,2,2,2)$ 中超曲面的 $J$ 函数},\enspace\IndexPageList[0429]{409,411}
\subitem \leavevmode\hypertarget{idx:0430}{五次三维簇的 $J$ 函数},\enspace\IndexPageList[0430]{48,385}
\subitem \leavevmode\hypertarget{idx:0431}{公式},\enspace\IndexPageList[0431]{345,379}
\subitem \leavevmode\hypertarget{idx:0432}{给出量子上同调中的关系},\enspace\IndexPageList[0432]{344}
\subitem \leavevmode\hypertarget{idx:0433}{本原部分},\enspace\IndexPageList[0433]{384}
\subitem \leavevmode\hypertarget{idx:0434}{$J_X$ 与 $J_{\mathcal V}$ 的关系},\enspace\IndexPageList[0434]{382,383,405,409,423}
\subitem \leavevmode\hypertarget{idx:0435}{符号记法},\enspace\IndexPageList[0435]{345,379,383}
\subitem \leavevmode\hypertarget{idx:0436}{环面部分},\enspace\IndexPageList[0436]{405}
\indexspace
\item \leavevmode\hypertarget{idx:0437}{Givental 联络},\enspace\IndexPageList[0437]{339}
\subitem \leavevmode\hypertarget{idx:0438}{定义},\enspace\IndexPageList[0438]{334}
\subitem \leavevmode\hypertarget{idx:0439}{对偶联络},\enspace\IndexPageList[0439]{344,352}
\subsubitem \leavevmode\hypertarget{source:481}{}\hypertarget{idx:0440}{平坦截面},\enspace\IndexPageList[0440]{334,335,336,342}
\IndexSubSubSubItem \leavevmode\hypertarget{idx:0441}{限制},\enspace\IndexPageList[0441]{337,345}
\IndexSubSubSubItem \leavevmode\hypertarget{idx:0442}{符号记法},\enspace\IndexPageList[0442]{335}
\subitem \leavevmode\hypertarget{idx:0443}{与 Dubrovin 联络的关系},\enspace\IndexPageList[0443]{268}
\subitem \leavevmode\hypertarget{idx:0444}{在 $H^0(X,\mathbb C)\oplus H^2(X,\mathbb C)$ 上的限制},\enspace\IndexPageList[0444]{337,343}
\subitem \leavevmode\hypertarget{idx:0445}{在 $H^2(V,\mathbb C)$ 上的限制},\enspace\IndexPageList[0445]{338}
\indexspace
\item \leavevmode\hypertarget{idx:0449}{GKZ 方程组},\enspace\IndexPageList[0449]{118,119,123,124,145,146,148,181,182,183,184,185,186,187,188,189,356,406}，\textit{另见} \hyperlink{idx:0055}{$\mathcal A$ 方程组}；\hyperlink{idx:0615}{超几何方程}
\indexspace
\item \leavevmode\hypertarget{idx:0446}{GKZ 分解},\enspace\IndexPageList[0446]{66,67,83,143,144,162,166,178,184,248,453}
\subitem \leavevmode\hypertarget{idx:0447}{扩大 GKZ 分解},\enspace\IndexPageList[0447]{66,453}
\subitem \leavevmode\hypertarget{idx:0448}{例题},\enspace\IndexPageList[0448]{73,74,75,147,150,163}
\indexspace
\item \leavevmode\hypertarget{idx:0450}{GLSM}，\textit{参见} \hyperlink{idx:0354}{规范线性 $\sigma$ 模型}
\indexspace
\item \leavevmode\hypertarget{idx:0452}{Gr\"obner 基方法},\enspace\IndexPageList[0452]{107,109,110,111,113,248}
\indexspace
\item \leavevmode\hypertarget{idx:0453}{Gr\"obner 扇},\enspace\IndexPageList[0453]{144,248}，\textit{另见} \hyperlink{idx:1239}{次级扇}
\indexspace
\item \leavevmode\hypertarget{idx:0454}{Grassmann 簇},\enspace\IndexPageList[0454]{190,196,242,250,251,304,316,428}
\subitem \leavevmode\hypertarget{idx:0455}{其中的 Calabi--Yau 完全交},\enspace\IndexPageList[0455]{428}
\subsubitem \leavevmode\hypertarget{idx:0456}{三维簇},\enspace\IndexPageList[0456]{389}
\indexspace
\item \leavevmode\hypertarget{idx:0483}{Griffiths 横截性},\enspace\IndexPageList[0483]{97,99,125,126,134,275,277}
\subitem \leavevmode\hypertarget{idx:0484}{射影超曲面的 Griffiths 横截性},\enspace\IndexPageList[0484]{106}
\subitem \leavevmode\hypertarget{idx:0485}{$A$ 型 Hodge 结构变分的 Griffiths 横截性}，\textit{参见} \hyperlink{idx:0021}{A 型 Hodge 结构变分，Griffiths 横截性}
\indexspace
\item \leavevmode\hypertarget{idx:0486}{Griffiths--Dwork 方法},\enspace\IndexPageList[0486]{106,107,108,109,110,117,119}，\textit{另见} \hyperlink{idx:1049}{Picard--Fuchs 方程}
\subitem \leavevmode\hypertarget{idx:0487}{方法说明},\enspace\IndexPageList[0487]{107}
\subitem \leavevmode\hypertarget{idx:0488}{例题},\enspace\IndexPageList[0488]{110,111,112,113}
\subitem \leavevmode\hypertarget{idx:0489}{用于射影超曲面},\enspace\IndexPageList[0489]{106,107,108}
\subitem \leavevmode\hypertarget{idx:0490}{用于环面超曲面},\enspace\IndexPageList[0490]{108,109,110}
\subitem \leavevmode\hypertarget{idx:0491}{局限},\enspace\IndexPageList[0491]{109}
\subitem \leavevmode\hypertarget{idx:0492}{用于加权射影超曲面},\enspace\IndexPageList[0492]{108}
\indexspace
\item \leavevmode\hypertarget{idx:0493}{Gromov--Uhlenbeck 紧化},\enspace\IndexPageList[0493]{210,232}
\indexspace
\item \leavevmode\hypertarget{idx:0515}{Gromov--Witten 不变量},\enspace\IndexPageList[0515]{30,32,38,39,47,48,51,288,289,296,304,321,322,323,324,330}
\subitem \leavevmode\hypertarget{idx:0516}{代数定义},\enspace\IndexPageList[0516]{191,192,193,194,195,196,197,198,199,200,201,202,203,204,205,206,207}，\textit{另见} \hyperlink{idx:1406}{虚拟基本类}
\subsubitem \leavevmode\hypertarget{source:482}{}\hypertarget{idx:0517}{一般情形},\enspace\IndexPageList[0517]{206}
\subsubitem \leavevmode\hypertarget{idx:0518}{最简单情形},\enspace\IndexPageList[0518]{194}
\subsubitem \leavevmode\hypertarget{idx:0519}{直观思想},\enspace\IndexPageList[0519]{191}
\subsubitem \leavevmode\hypertarget{idx:0520}{基本类的作用},\enspace\IndexPageList[0520]{195}
\subitem \leavevmode\hypertarget{idx:0521}{在物理中的出现},\enspace\IndexPageList[0521]{445}
\subitem \leavevmode\hypertarget{idx:0522}{公理},\enspace\IndexPageList[0522]{215,216}，\textit{另见} \hyperlink{idx:0495}{Gromov--Witten 类，公理}
\subitem \leavevmode\hypertarget{idx:0523}{计算},\enspace\IndexPageList[0523]{219,220,221,222,223,224,225,226,227,228,229,230,231,232,233,234,235,236,237,238}
\subsubitem \leavevmode\hypertarget{idx:0524}{$\mathbb P^2$ 的椭圆不变量},\enspace\IndexPageList[0524]{233,234,235,333}
\subsubitem \leavevmode\hypertarget{idx:0525}{$\mathbb P^3$ 的椭圆不变量},\enspace\IndexPageList[0525]{236}
\subsubitem \leavevmode\hypertarget{idx:0526}{五次三维簇的椭圆不变量},\enspace\IndexPageList[0526]{236,237}
\subsubitem \leavevmode\hypertarget{idx:0527}{$\mathbb P^2$ 的亏格 $0$ 不变量},\enspace\IndexPageList[0527]{220,221,222,223,224,259,260,333}
\subsubitem \leavevmode\hypertarget{idx:0528}{Calabi--Yau 三维簇的亏格 $0$ 不变量},\enspace\IndexPageList[0528]{224,225,226,227,228,229,230,231,232,233,331,348}，\textit{另见} \hyperlink{idx:0628}{瞬子数}；\hyperlink{idx:0915}{多重覆盖}
\subsubitem \leavevmode\hypertarget{idx:0529}{五次三维簇的亏格 $0$ 不变量},\enspace\IndexPageList[0529]{207,224,225,226}，\textit{另见} \hyperlink{idx:1186}{五次三维簇，Gromov--Witten 不变量}
\subsubitem \leavevmode\hypertarget{idx:0530}{$\mathbb P^2$ 的亏格 $2$ 不变量},\enspace\IndexPageList[0530]{334}
\subsubitem \leavevmode\hypertarget{idx:0531}{其他例题},\enspace\IndexPageList[0531]{224,237,428}
\subsubitem \leavevmode\hypertarget{idx:0532}{$\mathbb P^2$ 的不变量的递推关系},\enspace\IndexPageList[0532]{221,223,234,243,255,260}
\subsubitem \leavevmode\hypertarget{idx:0533}{$\mathbb P^3$ 的不变量的递推关系},\enspace\IndexPageList[0533]{236}
\subitem \leavevmode\hypertarget{idx:0534}{枚举意义},\enspace\IndexPageList[0534]{218,219,221,225,234,235,237,285,423}，\textit{另见} \hyperlink{idx:0632}{瞬子数，枚举意义}
\subitem \leavevmode\hypertarget{idx:0535}{代数定义与辛定义的等价性},\enspace\IndexPageList[0535]{218}
\subitem \leavevmode\hypertarget{idx:0536}{等变不变量},\enspace\IndexPageList[0536]{219,321,322,323,342}
\subsubitem \leavevmode\hypertarget{idx:0537}{定义},\enspace\IndexPageList[0537]{321}
\subsubitem \leavevmode\hypertarget{idx:0538}{非等变极限},\enspace\IndexPageList[0538]{321}
\subsubitem \leavevmode\hypertarget{idx:0539}{五次三维簇的等变不变量},\enspace\IndexPageList[0539]{322}
\subsubitem \leavevmode\hypertarget{idx:0540}{等变丛零截面的等变不变量},\enspace\IndexPageList[0540]{322}
\subitem \leavevmode\hypertarget{idx:0541}{辛定义},\enspace\IndexPageList[0541]{207,211,212,213,214,228,231,232}
\subsubitem \leavevmode\hypertarget{idx:0542}{定义的等价性},\enspace\IndexPageList[0542]{214,218}
\subsubitem \leavevmode\hypertarget{idx:0543}{直观思想},\enspace\IndexPageList[0543]{209}
\subsubitem \leavevmode\hypertarget{idx:0544}{Li--Tian 方法},\enspace\IndexPageList[0544]{213}
\subsubitem \leavevmode\hypertarget{idx:0545}{混合不变量与辛不变量的关系},\enspace\IndexPageList[0545]{213}
\subsubitem \leavevmode\hypertarget{idx:0546}{Ruan--Tian 混合不变量},\enspace\IndexPageList[0546]{211,212,233,235}
\subsubitem \leavevmode\hypertarget{idx:0547}{Ruan--Tian 辛不变量},\enspace\IndexPageList[0547]{212,233,235}
\subitem \leavevmode\hypertarget{idx:0548}{树级不变量},\enspace\IndexPageList[0548]{219,253}
\subitem \leavevmode\hypertarget{idx:0549}{用于定义 $A$ 模型的关联函数},\enspace\IndexPageList[0549]{244}
\subitem \leavevmode\hypertarget{idx:0550}{用于定义 Gromov--Witten 势},\enspace\IndexPageList[0550]{253}
\subitem \leavevmode\hypertarget{idx:0551}{用于定义量子上同调},\enspace\IndexPageList[0551]{241}
\indexspace
\item \leavevmode\hypertarget{idx:0494}{Gromov--Witten 类},\enspace\IndexPageList[0494]{223,235,243,255,328,332}
\subitem \leavevmode\hypertarget{idx:0495}{公理},\enspace\IndexPageList[0495]{214,215,216,217,218,327}
\subsubitem \leavevmode\hypertarget{idx:0496}{复合律},\enspace\IndexPageList[0496]{217}
\subsubitem \leavevmode\hypertarget{idx:0497}{形变公理},\enspace\IndexPageList[0497]{217,418}
\subsubitem \leavevmode\hypertarget{idx:0498}{次数公理},\enspace\IndexPageList[0498]{215,221,224,288}
\subsubitem \leavevmode\hypertarget{idx:0499}{除子公理},\enspace\IndexPageList[0499]{216,221,223,224,225,226,235,258,259,262,269,288,330}
\subsubitem \leavevmode\hypertarget{idx:0500}{有效性公理},\enspace\IndexPageList[0500]{215,251}
\subsubitem \leavevmode\hypertarget{idx:0501}{等变性公理},\enspace\IndexPageList[0501]{215,222,224,242,243,244,256,264}
\subsubitem \leavevmode\hypertarget{idx:0502}{基本类公理},\enspace\IndexPageList[0502]{216,221,222,224,235,242,257,258,288}
\subsubitem \leavevmode\hypertarget{idx:0503}{线性公理},\enspace\IndexPageList[0503]{215}
\subsubitem \leavevmode\hypertarget{idx:0504}{动机公理},\enspace\IndexPageList[0504]{217}
\subsubitem \leavevmode\hypertarget{idx:0505}{点映射公理},\enspace\IndexPageList[0505]{216,220,223,224,227,234,242,254,257,258,262,330,333,337}
\subsubitem \leavevmode\hypertarget{idx:0506}{约化公理},\enspace\IndexPageList[0506]{217,235,333}
\subsubitem \leavevmode\hypertarget{idx:0507}{分裂公理},\enspace\IndexPageList[0507]{217,219,221,222,242,243,256,264}
\subitem \leavevmode\hypertarget{idx:0508}{最简单情形下的定义},\enspace\IndexPageList[0508]{194}
\subitem \leavevmode\hypertarget{idx:0509}{定义},\enspace\IndexPageList[0509]{206}，\textit{另见} \hyperlink{idx:0516}{Gromov--Witten 不变量，代数定义}；\hyperlink{idx:1406}{虚拟基本类}
\subitem \leavevmode\hypertarget{idx:0510}{直观定义},\enspace\IndexPageList[0510]{191}
\subitem \leavevmode\hypertarget{idx:0511}{$\mathbb P^2$ 的 Gromov--Witten 类},\enspace\IndexPageList[0511]{221,222}
\subitem \leavevmode\hypertarget{idx:0512}{基本类的作用},\enspace\IndexPageList[0512]{195}
\subitem \leavevmode\hypertarget{idx:0513}{辛 Gromov--Witten 类},\enspace\IndexPageList[0513]{213}
\subitem \leavevmode\hypertarget{idx:0514}{树级 Gromov--Witten 类},\enspace\IndexPageList[0514]{219}
\indexspace
\item \leavevmode\hypertarget{idx:0552}{Gromov--Witten 势},\enspace\IndexPageList[0552]{51,132,253,254,255,256,257,267,270,326,327,412}
\subitem \leavevmode\hypertarget{idx:0553}{经典部分},\enspace\IndexPageList[0553]{258}
\subitem \leavevmode\hypertarget{idx:0554}{定义},\enspace\IndexPageList[0554]{253}
\subitem \leavevmode\hypertarget{idx:0555}{等变势},\enspace\IndexPageList[0555]{342,343}
\subitem \leavevmode\hypertarget{idx:0556}{亏格 $0$ 自由能},\enspace\IndexPageList[0556]{254}
\subitem \leavevmode\hypertarget{idx:0557}{亏格 $g$ 的势},\enspace\IndexPageList[0557]{254}
\subitem \leavevmode\hypertarget{idx:0558}{齐次性},\enspace\IndexPageList[0558]{258}
\subitem \leavevmode\hypertarget{idx:0559}{$\mathbb P^2$ 的势},\enspace\IndexPageList[0559]{259}
\subitem \leavevmode\hypertarget{idx:0560}{Calabi--Yau 三维簇的势},\enspace\IndexPageList[0560]{261,262,279,284,341,349}
\subsubitem \leavevmode\hypertarget{idx:0561}{可用镜像的周期表示},\enspace\IndexPageList[0561]{356,412}
\subsubitem \leavevmode\hypertarget{idx:0562}{环面部分},\enspace\IndexPageList[0562]{412}
\subitem \leavevmode\hypertarget{idx:0563}{椭圆曲线的势},\enspace\IndexPageList[0563]{254,258}
\subitem \leavevmode\hypertarget{idx:0564}{五次三维簇的势},\enspace\IndexPageList[0564]{51,260,261,356}
\subitem \leavevmode\hypertarget{idx:0565}{量子部分},\enspace\IndexPageList[0565]{258}
\subitem \leavevmode\hypertarget{idx:0566}{结构},\enspace\IndexPageList[0566]{257,258,259}
\subitem \leavevmode\hypertarget{idx:0567}{三阶偏导数},\enspace\IndexPageList[0567]{254,257}
\subitem \leavevmode\hypertarget{idx:0568}{WDVV 方程},\enspace\IndexPageList[0568]{255,256,257,259,260,262}
\indexspace
\item \leavevmode\hypertarget{idx:0570}{Gysin 谱序列},\enspace\IndexPageList[0570]{82}
\indexspace
\item \leavevmode\hypertarget{idx:0569}{Gysin 映射},\enspace\IndexPageList[0569]{194}
\indexspace
\item \leavevmode\hypertarget{idx:0571}{Hamilton 方程},\enspace\IndexPageList[0571]{436,437}
\indexspace
\item \leavevmode\hypertarget{idx:0572}{Hamilton 量},\enspace\IndexPageList[0572]{435,436,437,438,453}
\subitem \leavevmode\hypertarget{idx:0573}{电磁理论的 Hamilton 量},\enspace\IndexPageList[0573]{441}
\indexspace
\item \leavevmode\hypertarget{idx:0576}{Hamilton 向量场},\enspace\IndexPageList[0576]{63}
\indexspace
\item \leavevmode\hypertarget{idx:0574}{Hamilton 作用},\enspace\IndexPageList[0574]{63,452}
\subitem \leavevmode\hypertarget{idx:0575}{矩映射}，\textit{参见} \hyperlink{idx:0898}{矩映射}
\indexspace
\item \leavevmode\hypertarget{idx:0578}{Heisenberg 绘景},\enspace\IndexPageList[0578]{436}
\indexspace
\item \leavevmode\hypertarget{idx:0579}{$HG$}
\subitem \leavevmode\hypertarget{idx:0580}{定义},\enspace\IndexPageList[0580]{369}
\subitem \leavevmode\hypertarget{idx:0581}{镜像变换},\enspace\IndexPageList[0581]{371,372,376,377}，\textit{另见} \hyperlink{idx:0268}{Euler 数据，相关联的 Euler 数据，镜像变换}
\subitem \leavevmode\hypertarget{idx:0582}{$\widehat P$ 的 $HG$},\enspace\IndexPageList[0582]{51,124,187,370,375,376,377,380}
\subitem \leavevmode\hypertarget{idx:0583}{$\widehat Q$ 的 $HG$},\enspace\IndexPageList[0583]{51,370,373,375,376,377,380}
\subitem \leavevmode\hypertarget{idx:0584}{Euler 数据的 $HG$},\enspace\IndexPageList[0584]{369}
\indexspace
\item \leavevmode\hypertarget{idx:0585}{Hilbert 概形},\enspace\IndexPageList[0585]{231,305}
\indexspace
\item \leavevmode\hypertarget{idx:0586}{Hirzebruch--Riemann--Roch 定理},\enspace\IndexPageList[0586]{198}
\indexspace
\item \leavevmode\hypertarget{idx:0587}{Hochschild 上同调},\enspace\IndexPageList[0587]{426}
\indexspace
\item \leavevmode\hypertarget{idx:0593}{Hodge 结构}
\subitem \leavevmode\hypertarget{idx:0594}{退化},\enspace\IndexPageList[0594]{98}
\subitem \leavevmode\hypertarget{idx:0595}{Hodge 滤过},\enspace\IndexPageList[0595]{96,128,176}
\subsubitem \leavevmode\hypertarget{idx:0596}{极限 Hodge 滤过},\enspace\IndexPageList[0596]{99,134}
\subitem \leavevmode\hypertarget{idx:0597}{整数格},\enspace\IndexPageList[0597]{96,99,283,291}
\subitem \leavevmode\hypertarget{idx:0598}{混合 Hodge 结构},\enspace\IndexPageList[0598]{99,276}
\subsubitem \leavevmode\hypertarget{idx:0599}{Hodge--Tate 型},\enspace\IndexPageList[0599]{103,128,278,287,289,291}
\subsubitem \leavevmode\hypertarget{idx:0600}{极限混合 Hodge 结构},\enspace\IndexPageList[0600]{98,128,278,289}，\textit{另见} \hyperlink{idx:0907}{单值性权滤过}
\subitem \leavevmode\hypertarget{idx:0601}{轨形的 Hodge 结构},\enspace\IndexPageList[0601]{432}
\subitem \leavevmode\hypertarget{source:483}{}\hypertarget{idx:0602}{权为 $k$ 的 Hodge 结构},\enspace\IndexPageList[0602]{96}
\subitem \leavevmode\hypertarget{idx:0603}{极化 Hodge 结构},\enspace\IndexPageList[0603]{97,130,276,281,290}
\subitem \leavevmode\hypertarget{idx:0604}{纯 Hodge 结构},\enspace\IndexPageList[0604]{82}
\subitem \leavevmode\hypertarget{idx:0605}{Hodge 结构变分},\enspace\IndexPageList[0605]{48,97,130,132,175,181,276,277,281,290}
\subsubitem \leavevmode\hypertarget{idx:0606}{$\mathbb C$ 变分},\enspace\IndexPageList[0606]{291,294,295}
\subsubitem \leavevmode\hypertarget{idx:0607}{来自几何的变分},\enspace\IndexPageList[0607]{285}
\subsubitem \leavevmode\hypertarget{idx:0608}{形式退化的变分},\enspace\IndexPageList[0608]{284}
\subsubitem \leavevmode\hypertarget{idx:0609}{实变分},\enspace\IndexPageList[0609]{278,287}
\indexspace
\item \leavevmode\hypertarget{idx:0588}{Hodge 菱形图},\enspace\IndexPageList[0588]{27,28}
\indexspace
\item \leavevmode\hypertarget{idx:0589}{Hodge 数},\enspace\IndexPageList[0589]{27,29,89,250}
\subitem \leavevmode\hypertarget{idx:0590}{Calabi--Yau 环面完全交的 Hodge 数},\enspace\IndexPageList[0590]{86}
\subitem \leavevmode\hypertarget{idx:0591}{Calabi--Yau 环面超曲面的 Hodge 数},\enspace\IndexPageList[0591]{79,80,81,82,83,142,168,169}
\subitem \leavevmode\hypertarget{idx:0592}{弦论 Hodge 数},\enspace\IndexPageList[0592]{83,86,296,425}
\indexspace
\item \leavevmode\hypertarget{idx:0610}{Hodge--Riemann 双线性关系},\enspace\IndexPageList[0610]{97,277,291}
\indexspace
\item \leavevmode\hypertarget{idx:0611}{Hodge--Tate 型}，\textit{参见} \hyperlink{idx:0599}{Hodge 结构，混合 Hodge 结构，Hodge--Tate 型}
\indexspace
\item \leavevmode\hypertarget{idx:0614}{Hopf 纤维化},\enspace\IndexPageList[0614]{72}
\indexspace
\item \leavevmode\hypertarget{idx:0641}{$I_{\mathcal V}$}，\textit{参见} \hyperlink{idx:0383}{Givental $I$ 函数，$I_{\mathcal V}$}
\indexspace
\item \leavevmode\hypertarget{idx:0644}{$J$ 全纯曲线},\enspace\IndexPageList[0644]{232}
\subitem \leavevmode\hypertarget{idx:0645}{定义},\enspace\IndexPageList[0645]{231}
\indexspace
\item \leavevmode\hypertarget{idx:0646}{$J$ 全纯映射},\enspace\IndexPageList[0646]{208}
\subitem \leavevmode\hypertarget{idx:0647}{模空间}
\subsubitem \leavevmode\hypertarget{idx:0648}{维数},\enspace\IndexPageList[0648]{209,210}
\subsubitem \leavevmode\hypertarget{idx:0649}{基本类},\enspace\IndexPageList[0649]{211}
\subsubitem \leavevmode\hypertarget{idx:0650}{Li--Tian 方法},\enspace\IndexPageList[0650]{211,213}
\subsubitem \leavevmode\hypertarget{idx:0651}{Ruan 方法},\enspace\IndexPageList[0651]{209}
\subsubitem \leavevmode\hypertarget{idx:0652}{Ruan--Tian 方法},\enspace\IndexPageList[0652]{210}
\subitem \leavevmode\hypertarget{idx:0653}{多重覆盖},\enspace\IndexPageList[0653]{209}
\subitem \leavevmode\hypertarget{idx:0654}{受扰动的 $J$ 全纯映射},\enspace\IndexPageList[0654]{210,213,228}
\subitem \leavevmode\hypertarget{idx:0655}{简单 $J$ 全纯映射},\enspace\IndexPageList[0655]{209,231,233}
\subitem \leavevmode\hypertarget{idx:0656}{稳定 $J$ 全纯映射},\enspace\IndexPageList[0656]{211}
\indexspace
\item \leavevmode\hypertarget{idx:0658}{$J_{\mathcal V}$}，\textit{参见} \hyperlink{idx:0402}{Givental $J$ 函数，$J_{\mathcal V}$}
\indexspace
\item \leavevmode\hypertarget{idx:0643}{Jacobian 环},\enspace\IndexPageList[0643]{450,451}
\indexspace
\item \leavevmode\hypertarget{idx:0642}{Jacobian 理想},\enspace\IndexPageList[0642]{107,109,113}
\indexspace
\item \leavevmode\hypertarget{idx:0657}{Jones 多项式},\enspace\IndexPageList[0657]{458}
\indexspace
\item \leavevmode\hypertarget{idx:0677}{$K_d$},\enspace\IndexPageList[0677]{356,373}，\textit{另见} \hyperlink{idx:0927}{$N_d$，五次三维簇的 Gromov--Witten 不变量}
\subitem \leavevmode\hypertarget{idx:0678}{临界丛的 $K_d$},\enspace\IndexPageList[0678]{375}
\subitem \leavevmode\hypertarget{idx:0679}{多重覆盖情形的 $K_d$},\enspace\IndexPageList[0679]{375,378}
\indexspace
\item \leavevmode\hypertarget{idx:0659}{K3 曲面},\enspace\IndexPageList[0659]{382}
\subitem \leavevmode\hypertarget{idx:0660}{镜像对称},\enspace\IndexPageList[0660]{424}
\subitem \leavevmode\hypertarget{idx:0661}{其上的有理曲线},\enspace\IndexPageList[0661]{423}
\subitem \leavevmode\hypertarget{idx:0662}{其上的向量丛},\enspace\IndexPageList[0662]{427}
\subitem \leavevmode\hypertarget{idx:0663}{带有对合的 K3 曲面},\enspace\IndexPageList[0663]{88,89,90}
\indexspace
\item \leavevmode\hypertarget{idx:0664}{Kähler}
\subitem \leavevmode\hypertarget{idx:0665}{Kähler 锥},\enspace\IndexPageList[0665]{35,143,150,151,153,173,248,250,388}
\subsubitem \leavevmode\hypertarget{idx:0666}{复结构变动时的变化},\enspace\IndexPageList[0666]{157}
\subsubitem \leavevmode\hypertarget{idx:0667}{闭包},\enspace\IndexPageList[0667]{151,155,158,165,176,266,268,290,339,340,404}
\subsubitem \leavevmode\hypertarget{idx:0668}{面},\enspace\IndexPageList[0668]{155,156,157,167,169,181,250}
\subsubitem \leavevmode\hypertarget{idx:0669}{局部多面锥性质},\enspace\IndexPageList[0669]{155}
\subsubitem \leavevmode\hypertarget{idx:0670}{Calabi--Yau 三维簇的 Kähler 锥},\enspace\IndexPageList[0670]{155,167,169}，\textit{另见} \hyperlink{idx:0913}{可动锥}
\subsubitem \leavevmode\hypertarget{idx:0671}{多面锥性质},\enspace\IndexPageList[0671]{62,152}
\subsubitem \leavevmode\hypertarget{idx:0672}{单纯锥性质},\enspace\IndexPageList[0672]{120}
\subsubitem \leavevmode\hypertarget{idx:0673}{环面部分},\enspace\IndexPageList[0673]{158,159,161,162,422}
\subitem \leavevmode\hypertarget{idx:0674}{轨形上的 Kähler 形式},\enspace\IndexPageList[0674]{432}
\subitem \leavevmode\hypertarget{idx:0675}{Kähler 模}，\textit{参见} \hyperlink{idx:0882}{模，Kähler 模}
\subitem \leavevmode\hypertarget{idx:0676}{Kähler 空间}，\textit{参见} \hyperlink{idx:0143}{复化，复化 Kähler 空间}
\indexspace
\item \leavevmode\hypertarget{idx:0680}{Klein--Gordon 方程},\enspace\IndexPageList[0680]{435,437,438,441}
\indexspace
\item \leavevmode\hypertarget{idx:0681}{Kodaira--Spencer 类},\enspace\IndexPageList[0681]{156}
\indexspace
\item \leavevmode\hypertarget{idx:0682}{Kodaira--Spencer 映射},\enspace\IndexPageList[0682]{176}
\indexspace
\item \leavevmode\hypertarget{idx:0683}{Kuranishi 映射},\enspace\IndexPageList[0683]{198,200}
\indexspace
\item \leavevmode\hypertarget{source:484}{}\hypertarget{idx:0705}{$L_d$},\enspace\IndexPageList[0705]{359,398}
\subitem \leavevmode\hypertarget{idx:0706}{不动点轨迹},\enspace\IndexPageList[0706]{399}
\subsubitem \leavevmode\hypertarget{idx:0707}{法丛},\enspace\IndexPageList[0707]{400}
\subitem \leavevmode\hypertarget{idx:0708}{到 $N_d$ 的映射 $\mu$},\enspace\IndexPageList[0708]{398}
\indexspace
\item \leavevmode\hypertarget{idx:0715}{$\mathcal L_i$}，\textit{参见} \hyperlink{idx:0198}{余切线 $\mathcal L_i$}
\indexspace
\item \leavevmode\hypertarget{idx:0686}{Lagrange 范畴},\enspace\IndexPageList[0686]{426}
\indexspace
\item \leavevmode\hypertarget{idx:0685}{Lagrange 量},\enspace\IndexPageList[0685]{24,435,443}
\indexspace
\item \leavevmode\hypertarget{idx:0687}{Lagrange 密度},\enspace\IndexPageList[0687]{434,435,437,438,439}
\subitem \leavevmode\hypertarget{idx:0688}{电磁理论的 Lagrange 密度},\enspace\IndexPageList[0688]{440,441}
\indexspace
\item \leavevmode\hypertarget{idx:0684}{Lagrange 提升},\enspace\IndexPageList[0684]{372}，\textit{另见} \hyperlink{idx:0268}{Euler 数据，相关联的 Euler 数据，镜像变换}
\indexspace
\item \leavevmode\hypertarget{idx:0689}{Lagrange 子流形},\enspace\IndexPageList[0689]{426}
\subitem \leavevmode\hypertarget{idx:0690}{特殊 Lagrange 子流形},\enspace\IndexPageList[0690]{34,425,426}
\indexspace
\item \leavevmode\hypertarget{idx:0691}{Landau--Ginzburg}
\subitem \leavevmode\hypertarget{idx:0692}{模型},\enspace\IndexPageList[0692]{449}
\subitem \leavevmode\hypertarget{idx:0693}{轨形},\enspace\IndexPageList[0693]{29,166,450,452}
\subitem \leavevmode\hypertarget{idx:0694}{势},\enspace\IndexPageList[0694]{451}
\subitem \leavevmode\hypertarget{idx:0695}{理论},\enspace\IndexPageList[0695]{449,450,451,452}
\indexspace
\item \leavevmode\hypertarget{idx:0703}{Laurent 单项式},\enspace\IndexPageList[0703]{58,141,149}
\indexspace
\item \leavevmode\hypertarget{idx:0704}{Laurent 多项式},\enspace\IndexPageList[0704]{58,139,140,145,178}
\indexspace
\item \leavevmode\hypertarget{idx:0710}{Lefschetz 定理},\enspace\IndexPageList[0710]{384}
\indexspace
\item \leavevmode\hypertarget{idx:0709}{Lefschetz 分解},\enspace\IndexPageList[0709]{282,287,384}
\indexspace
\item \leavevmode\hypertarget{idx:0714}{Levi--Civita 联络},\enspace\IndexPageList[0714]{443}
\indexspace
\item \leavevmode\hypertarget{idx:0744}{Lorentz 变换},\enspace\IndexPageList[0744]{438}
\indexspace
\item \leavevmode\hypertarget{idx:0746}{$M(X)$}，\textit{参见} \hyperlink{idx:0911}{Mori 锥}
\indexspace
\item \leavevmode\hypertarget{idx:0725}{$\overline M_{0,n}(\mathbb P^n,d)$ 中的局部化},\enspace\IndexPageList[0725]{305,306,307,308,309,310,311,312,313,314,315,316}
\subitem \leavevmode\hypertarget{idx:0726}{不动点轨迹},\enspace\IndexPageList[0726]{305,363,371,390,391,399,400}
\subsubitem \leavevmode\hypertarget{idx:0727}{描述},\enspace\IndexPageList[0727]{306}
\subsubitem \leavevmode\hypertarget{idx:0728}{例题},\enspace\IndexPageList[0728]{307,308,312}
\subsubitem \leavevmode\hypertarget{idx:0729}{对应的图},\enspace\IndexPageList[0729]{306,312,313,390,391,393}
\subsubitem \leavevmode\hypertarget{idx:0730}{法丛},\enspace\IndexPageList[0730]{308,309,310,312,314,316,393}，\textit{另见} \hyperlink{idx:0259}{Euler 类，等变 Euler 类，$N_\Gamma$ 的 Euler 类}
\subsubitem \leavevmode\hypertarget{idx:0731}{A 型图},\enspace\IndexPageList[0731]{391,392,397}
\subsubitem \leavevmode\hypertarget{idx:0732}{B 型图},\enspace\IndexPageList[0732]{391,397}
\subitem \leavevmode\hypertarget{idx:0733}{其上的环面作用},\enspace\IndexPageList[0733]{305}
\indexspace
\item \leavevmode\hypertarget{idx:0777}{$M_d$},\enspace\IndexPageList[0777]{357}
\subitem \leavevmode\hypertarget{idx:0778}{不动点轨迹},\enspace\IndexPageList[0778]{357,362}
\subsubitem \leavevmode\hypertarget{idx:0779}{描述},\enspace\IndexPageList[0779]{357}
\subsubitem \leavevmode\hypertarget{idx:0780}{图 $\Gamma$},\enspace\IndexPageList[0780]{362}
\subsubitem \leavevmode\hypertarget{idx:0781}{法丛},\enspace\IndexPageList[0781]{364}
\subitem \leavevmode\hypertarget{idx:0782}{到 $N_d$ 的映射 $\varphi$},\enspace\IndexPageList[0782]{359,360,362,363,364,367,398}
\subitem \leavevmode\hypertarget{idx:0783}{其上的环面作用},\enspace\IndexPageList[0783]{357}
\indexspace
\item \leavevmode\hypertarget{idx:0747}{Macaulay},\enspace\IndexPageList[0747]{320}
\indexspace
\item \leavevmode\hypertarget{idx:0775}{Maxwell 方程},\enspace\IndexPageList[0775]{440,441}
\indexspace
\item \leavevmode\hypertarget{idx:0776}{McKay 对应},\enspace\IndexPageList[0776]{425}
\indexspace
\item \leavevmode\hypertarget{idx:0785}{Minkowski 和},\enspace\IndexPageList[0785]{56,86}
\indexspace
\item \leavevmode\hypertarget{source:486}{}\hypertarget{idx:0911}{Mori 锥},\enspace\IndexPageList[0911]{62,121,251}
\subitem \leavevmode\hypertarget{idx:0912}{环面簇的 Mori 锥}，\textit{参见} \hyperlink{idx:1372}{环面簇，Mori 锥}
\indexspace
\item \leavevmode\hypertarget{idx:0925}{$N_d$}
\subitem \leavevmode\hypertarget{idx:0926}{$\mathbb P^2$ 的 Gromov--Witten 不变量},\enspace\IndexPageList[0926]{221}
\subitem \leavevmode\hypertarget{idx:0927}{五次三维簇的 Gromov--Witten 不变量},\enspace\IndexPageList[0927]{224}
\subitem \leavevmode\hypertarget{idx:0928}{射影空间},\enspace\IndexPageList[0928]{357,398}
\subsubitem \leavevmode\hypertarget{idx:0929}{等变上同调},\enspace\IndexPageList[0929]{358}
\subsubitem \leavevmode\hypertarget{idx:0930}{等变超平面类},\enspace\IndexPageList[0930]{358}
\subsubitem \leavevmode\hypertarget{idx:0931}{不动点},\enspace\IndexPageList[0931]{358,363}
\subsubitem \leavevmode\hypertarget{idx:0932}{其上的环面作用},\enspace\IndexPageList[0932]{358}
\subsubitem \leavevmode\hypertarget{idx:0933}{权},\enspace\IndexPageList[0933]{358}
\indexspace
\item \leavevmode\hypertarget{idx:0935}{nef 分拆},\enspace\IndexPageList[0935]{85,86,120,123}
\subitem \leavevmode\hypertarget{idx:0936}{对偶 nef 分拆},\enspace\IndexPageList[0936]{86}
\indexspace
\item \leavevmode\hypertarget{idx:0934}{nef 完全交}，\textit{参见} \hyperlink{idx:1349}{环面完全交，nef 环面完全交}
\indexspace
\item \leavevmode\hypertarget{idx:0941}{Noether 定理},\enspace\IndexPageList[0941]{442}
\indexspace
\item \leavevmode\hypertarget{idx:0955}{Novikov 环},\enspace\IndexPageList[0955]{252}
\indexspace
\item \leavevmode\hypertarget{source:488}{}\hypertarget{idx:1043}{$\widehat P$},\enspace\IndexPageList[1043]{50,51,124,187,366,369,376}
\subitem \leavevmode\hypertarget{idx:1044}{定义},\enspace\IndexPageList[1044]{360}
\subitem \leavevmode\hypertarget{idx:1045}{凸凹丛的 $\widehat P$},\enspace\IndexPageList[1045]{373}
\subitem \leavevmode\hypertarget{idx:1046}{多重覆盖情形的 $\widehat P$},\enspace\IndexPageList[1046]{378}
\subitem \leavevmode\hypertarget{idx:1047}{是 Euler 数据},\enspace\IndexPageList[1047]{362}
\indexspace
\item \leavevmode\hypertarget{idx:0123}{$\mathcal P$ 类}，\textit{参见} \hyperlink{idx:0847}{镜像定理，Givental 镜像定理，$\mathcal P$ 类}
\indexspace
\item \leavevmode\hypertarget{idx:1001}{$\mathbb P(1,1,2,2,2)$ 及其环面爆破},\enspace\IndexPageList[1001]{72,112}
\subitem \leavevmode\hypertarget{idx:1002}{其中的 Calabi--Yau 超曲面},\enspace\IndexPageList[1002]{450}
\subsubitem \leavevmode\hypertarget{idx:1003}{A 模型关联函数},\enspace\IndexPageList[1003]{292,409,411}
\subsubitem \leavevmode\hypertarget{idx:1004}{Givental 函数 $I_{\mathcal V}$},\enspace\IndexPageList[1004]{407}
\subsubitem \leavevmode\hypertarget{idx:1005}{Givental 函数 $J_X$},\enspace\IndexPageList[1005]{409,411}
\subsubitem \leavevmode\hypertarget{idx:1006}{Givental 函数 $J_{\mathcal V}$},\enspace\IndexPageList[1006]{407,411}
\subsubitem \leavevmode\hypertarget{idx:1007}{Gromov--Witten 势},\enspace\IndexPageList[1007]{410,411}
\subsubitem \leavevmode\hypertarget{idx:1008}{Hodge 数},\enspace\IndexPageList[1008]{169}
\subsubitem \leavevmode\hypertarget{idx:1009}{瞬子数},\enspace\IndexPageList[1009]{292,293}
\subsubitem \leavevmode\hypertarget{idx:1010}{Kähler 锥},\enspace\IndexPageList[1010]{169,292,406}
\subsubitem \leavevmode\hypertarget{idx:1011}{Kähler 模},\enspace\IndexPageList[1011]{169,170}
\subsubitem \leavevmode\hypertarget{idx:1012}{镜像映射},\enspace\IndexPageList[1012]{181,182,187,188,189,293,408,411}
\subsubitem \leavevmode\hypertarget{idx:1013}{镜像定理},\enspace\IndexPageList[1013]{406,409,411}
\subsubitem \leavevmode\hypertarget{idx:1014}{模坐标},\enspace\IndexPageList[1014]{292}
\subsubitem \leavevmode\hypertarget{idx:1015}{超共形场论模},\enspace\IndexPageList[1015]{170}
\subsubitem \leavevmode\hypertarget{idx:1016}{环面 Kähler 锥},\enspace\IndexPageList[1016]{171}
\subsubitem \leavevmode\hypertarget{idx:1017}{环面 Kähler 模},\enspace\IndexPageList[1017]{171,292}
\subitem \leavevmode\hypertarget{idx:1018}{GKZ 分解},\enspace\IndexPageList[1018]{73,147,169}
\subitem \leavevmode\hypertarget{idx:1019}{镜像猜想},\enspace\IndexPageList[1019]{292}
\subitem \leavevmode\hypertarget{idx:1020}{其中 Calabi--Yau 超曲面的镜像},\enspace\IndexPageList[1020]{111,118,126,146,168,187}
\subsubitem \leavevmode\hypertarget{idx:1021}{$\mathcal A$ 方程组},\enspace\IndexPageList[1021]{119,408}
\subsubitem \leavevmode\hypertarget{idx:1022}{判别式轨迹},\enspace\IndexPageList[1022]{170}
\subsubitem \leavevmode\hypertarget{idx:1023}{Frobenius 方法},\enspace\IndexPageList[1023]{188}
\subsubitem \leavevmode\hypertarget{idx:1024}{Gauss--Manin 联络},\enspace\IndexPageList[1024]{292,408,409}
\subsubitem \leavevmode\hypertarget{idx:1025}{Hodge 数},\enspace\IndexPageList[1025]{169}
\subsubitem \leavevmode\hypertarget{idx:1026}{齐次坐标环},\enspace\IndexPageList[1026]{112}
\subsubitem \leavevmode\hypertarget{idx:1027}{模坐标},\enspace\IndexPageList[1027]{112,118,171,181,410}
\subsubitem \leavevmode\hypertarget{idx:1028}{单值性},\enspace\IndexPageList[1028]{293,408,410}
\subsubitem \leavevmode\hypertarget{idx:1029}{周期},\enspace\IndexPageList[1029]{187,188,408}
\subsubitem \leavevmode\hypertarget{idx:1030}{Picard--Fuchs 方程},\enspace\IndexPageList[1030]{113,119,126,181,188,408,409}
\subsubitem \leavevmode\hypertarget{idx:1031}{满足整性猜想},\enspace\IndexPageList[1031]{293,408}
\subsubitem \leavevmode\hypertarget{idx:1032}{简化模},\enspace\IndexPageList[1032]{170,171}
\subsubitem \leavevmode\hypertarget{idx:1033}{Yukawa 耦合},\enspace\IndexPageList[1033]{126,127,148,408,411}
\subitem \leavevmode\hypertarget{idx:1034}{多胞形},\enspace\IndexPageList[1034]{111}
\indexspace
\item \leavevmode\hypertarget{idx:1049}{Picard--Fuchs 方程},\enspace\IndexPageList[1049]{42,98,100,103,104,107,125,126}
\subitem \leavevmode\hypertarget{idx:1050}{通过 $\mathcal A$ 方程组计算}，\textit{参见} \hyperlink{idx:0055}{$\mathcal A$ 方程组}
\subitem \leavevmode\hypertarget{idx:1051}{通过 Griffiths--Dwork 方法计算}，\textit{参见} \hyperlink{idx:0486}{Griffiths--Dwork 方法}
\subitem \leavevmode\hypertarget{idx:1052}{例题},\enspace\IndexPageList[1052]{110,111,112,113}
\subitem \leavevmode\hypertarget{idx:1053}{$\mathcal A$ 方程组中缺失的方程},\enspace\IndexPageList[1053]{117,118,124,406,408,422}
\subitem \leavevmode\hypertarget{idx:1054}{归一化 Picard--Fuchs 方程},\enspace\IndexPageList[1054]{50,128,129,130}
\subitem \leavevmode\hypertarget{idx:1055}{Calabi--Yau 三维簇的 Picard--Fuchs 方程},\enspace\IndexPageList[1055]{102,128,350}
\subitem \leavevmode\hypertarget{idx:1056}{Calabi--Yau 环面超曲面的 Picard--Fuchs 方程}，\textit{参见} \hyperlink{idx:0099}{Calabi--Yau，Calabi--Yau 环面超曲面}
\subitem \leavevmode\hypertarget{idx:1057}{Calabi--Yau 环面完全交的镜像的 Picard--Fuchs 方程},\enspace\IndexPageList[1057]{123,124}
\subitem \leavevmode\hypertarget{idx:1058}{五次镜像的 Picard--Fuchs 方程}，\textit{参见} \hyperlink{idx:1175}{五次镜像，Picard--Fuchs 方程}
\subitem \leavevmode\hypertarget{idx:1059}{解},\enspace\IndexPageList[1059]{42,119,124,172,179,183,184,286,291}，\textit{另见} \hyperlink{idx:0344}{Frobenius 方法}
\indexspace
\item \leavevmode\hypertarget{idx:1060}{Picard--Fuchs 理想},\enspace\IndexPageList[1060]{97,100,117}
\subitem \leavevmode\hypertarget{idx:1061}{相应的 $\mathcal D$ 模},\enspace\IndexPageList[1061]{97,98,124}
\indexspace
\item \leavevmode\hypertarget{idx:1062}{Picard--Fuchs 算子},\enspace\IndexPageList[1062]{351,352}
\subitem \leavevmode\hypertarget{idx:1063}{$\hbar$ 齐次化},\enspace\IndexPageList[1063]{351,352}
\subitem \leavevmode\hypertarget{idx:1064}{给出量子上同调中的关系},\enspace\IndexPageList[1064]{352}
\subitem \leavevmode\hypertarget{idx:1065}{主部},\enspace\IndexPageList[1065]{351}
\indexspace
\item \leavevmode\hypertarget{idx:1066}{Picard--Lefschetz 变换}，\textit{参见} \hyperlink{idx:0901}{单值性变换}
\indexspace
\item \leavevmode\hypertarget{idx:1067}{Planck 常数 $\hbar$},\enspace\IndexPageList[1067]{436}
\indexspace
\item \leavevmode\hypertarget{idx:1068}{Poincaré 对偶},\enspace\IndexPageList[1068]{105,404}
\subitem \leavevmode\hypertarget{idx:1069}{轨形的 Poincaré 对偶},\enspace\IndexPageList[1069]{194,195,206,328,431}
\indexspace
\item \leavevmode\hypertarget{idx:1070}{Poisson 括号},\enspace\IndexPageList[1070]{436,437}
\indexspace
\item \leavevmode\hypertarget{idx:1101}{$\widehat Q$},\enspace\IndexPageList[1101]{50,51,363,366,367,370,371,373,374,376}
\subitem \leavevmode\hypertarget{idx:1102}{定义},\enspace\IndexPageList[1102]{360}
\subitem \leavevmode\hypertarget{idx:1103}{凸凹丛的 $\widehat Q$},\enspace\IndexPageList[1103]{372}
\subitem \leavevmode\hypertarget{idx:1104}{多重覆盖情形的 $\widehat Q$},\enspace\IndexPageList[1104]{375,378}
\subitem \leavevmode\hypertarget{idx:1105}{是 Euler 数据},\enspace\IndexPageList[1105]{362}
\indexspace
\item \leavevmode\hypertarget{source:490}{}\hypertarget{idx:1228}{Reid 幻想},\enspace\IndexPageList[1228]{166,428}
\indexspace
\item \leavevmode\hypertarget{idx:1232}{Ricci 曲率张量},\enspace\IndexPageList[1232]{33}
\indexspace
\item \leavevmode\hypertarget{idx:1233}{Riemann--Roch 定理},\enspace\IndexPageList[1233]{312,317,318}
\indexspace
\item \leavevmode\hypertarget{idx:1237}{Satake--Baily--Borel 紧化},\enspace\IndexPageList[1237]{155}
\indexspace
\item \leavevmode\hypertarget{idx:1238}{Schrödinger 绘景},\enspace\IndexPageList[1238]{436}
\indexspace
\item \leavevmode\hypertarget{idx:1243}{Seiberg--Witten 曲线},\enspace\IndexPageList[1243]{429}
\indexspace
\item \leavevmode\hypertarget{idx:1246}{Serre--Grothendieck 对偶},\enspace\IndexPageList[1246]{78}
\indexspace
\item \leavevmode\hypertarget{idx:1247}{Severi 簇},\enspace\IndexPageList[1247]{221}
\indexspace
\item \leavevmode\hypertarget{idx:1284}{Stanley--Reisner 理想},\enspace\IndexPageList[1284]{246,247}
\subitem \leavevmode\hypertarget{idx:1285}{量子 Stanley--Reisner 理想},\enspace\IndexPageList[1285]{247,248,415}
\indexspace
\item \leavevmode\hypertarget{source:491}{}\hypertarget{idx:1325}{T 对偶},\enspace\IndexPageList[1325]{425,426}
\indexspace
\item \leavevmode\hypertarget{idx:1335}{Toda 格},\enspace\IndexPageList[1335]{251}
\indexspace
\item \leavevmode\hypertarget{idx:1381}{TQFT}，\textit{参见} \hyperlink{idx:1336}{拓扑量子场论}
\indexspace
\item \leavevmode\hypertarget{idx:1382}{TRR}，\textit{参见} \hyperlink{idx:1344}{拓扑递推关系}
\indexspace
\item \leavevmode\hypertarget{idx:1384}{$V$ 流形},\enspace\IndexPageList[1384]{431}，\textit{另见} \hyperlink{idx:0990}{轨形}
\indexspace
\item \leavevmode\hypertarget{idx:1388}{$\mathcal V_{\beta,k}$},\enspace\IndexPageList[1388]{387}，\textit{另见} \hyperlink{idx:1395}{$\mathcal V_{d,k}$}
\indexspace
\item \leavevmode\hypertarget{idx:1389}{$\mathcal V'_{\beta,k,i}$},\enspace\IndexPageList[1389]{387}，\textit{另见} \hyperlink{idx:1396}{$\mathcal V'_{d,k,i}$}
\indexspace
\item \leavevmode\hypertarget{idx:1390}{$\mathcal V_d$}
\subitem \leavevmode\hypertarget{idx:1391}{$\mathcal O_{\mathbb P^n}(\ell)$ 的 $\mathcal V_d$},\enspace\IndexPageList[1391]{360,363,367}
\subitem \leavevmode\hypertarget{idx:1392}{凸凹丛的 $\mathcal V_d$},\enspace\IndexPageList[1392]{372,374,381,398}
\subitem \leavevmode\hypertarget{idx:1393}{五次三维簇的 $\mathcal V_d$},\enspace\IndexPageList[1393]{204,207,224,311,322,325,355}
\subsubitem \leavevmode\hypertarget{idx:1394}{多重覆盖情形},\enspace\IndexPageList[1394]{375}
\indexspace
\item \leavevmode\hypertarget{idx:1395}{$\mathcal V_{d,k}$},\enspace\IndexPageList[1395]{380,381,383}
\indexspace
\item \leavevmode\hypertarget{idx:1396}{$\mathcal V'_{d,k,i}$},\enspace\IndexPageList[1396]{381,395,396}
\indexspace
\item \leavevmode\hypertarget{idx:1386}{Vafa--Intriligator 公式},\enspace\IndexPageList[1386]{250}
\indexspace
\item \leavevmode\hypertarget{idx:1398}{Verlinde 代数},\enspace\IndexPageList[1398]{250}
\indexspace
\item \leavevmode\hypertarget{idx:1401}{Virasoro 猜想},\enspace\IndexPageList[1401]{333,448}
\subitem \leavevmode\hypertarget{idx:1402}{对下列对象的推论}
\subsubitem \leavevmode\hypertarget{idx:1403}{引力后裔},\enspace\IndexPageList[1403]{334}
\subsubitem \leavevmode\hypertarget{idx:1404}{Gromov--Witten 不变量},\enspace\IndexPageList[1404]{333,334}
\indexspace
\item \leavevmode\hypertarget{idx:1399}{Virasoro 代数},\enspace\IndexPageList[1399]{333,334,446,447,448}
\subitem \leavevmode\hypertarget{idx:1400}{中心扩张},\enspace\IndexPageList[1400]{333,448}
\indexspace
\item \leavevmode\hypertarget{idx:1405}{Virasoro 约束},\enspace\IndexPageList[1405]{234}
\indexspace
\item \leavevmode\hypertarget{source:492}{}\hypertarget{idx:1430}{WDVV 方程}，\textit{参见} \hyperlink{idx:0226}{Dubrovin 联络}；\hyperlink{idx:0552}{Gromov--Witten 势}
\indexspace
\item \leavevmode\hypertarget{idx:1439}{Witten 猜想},\enspace\IndexPageList[1439]{333}
\indexspace
\item \leavevmode\hypertarget{idx:1443}{Yang--Mills 理论},\enspace\IndexPageList[1443]{440}
\indexspace
\item \leavevmode\hypertarget{idx:1444}{Yukawa 耦合},\enspace\IndexPageList[1444]{30,31,32,33,439,444}，\textit{另见} \hyperlink{idx:0183}{关联函数}
\subitem \leavevmode\hypertarget{idx:1445}{归一化 Yukawa 耦合},\enspace\IndexPageList[1445]{127,128,129,130,131,132,133,134,135,280,282,291,409}
\subsubitem \leavevmode\hypertarget{idx:1446}{定义},\enspace\IndexPageList[1446]{133}
\subsubitem \leavevmode\hypertarget{idx:1447}{用周期表示},\enspace\IndexPageList[1447]{412,422}
\subsubitem \leavevmode\hypertarget{idx:1448}{负号},\enspace\IndexPageList[1448]{31,44,130}
\subsubitem \leavevmode\hypertarget{idx:1449}{五次镜像的归一化 Yukawa 耦合}，\textit{参见} \hyperlink{idx:1179}{五次镜像，Yukawa 耦合}
\subsubitem \leavevmode\hypertarget{idx:1450}{与镜像映射的关系},\enspace\IndexPageList[1450]{132,133,134,135,174}
\subsubitem \leavevmode\hypertarget{idx:1451}{与 Gauss--Manin 联络的关系},\enspace\IndexPageList[1451]{128,131}
\subitem \leavevmode\hypertarget{idx:1452}{Calabi--Yau 三维簇的 Yukawa 耦合},\enspace\IndexPageList[1452]{125,126,127,128,129,130,131,132,133,134,135,279,282,350}
\subsubitem \leavevmode\hypertarget{idx:1453}{一维模},\enspace\IndexPageList[1453]{125,128,129,130}
\subsubitem \leavevmode\hypertarget{idx:1454}{二维模},\enspace\IndexPageList[1454]{125,126,127}
\subsubitem \leavevmode\hypertarget{idx:1455}{$r$ 维模},\enspace\IndexPageList[1455]{130,131,132}
\subitem \leavevmode\hypertarget{idx:1456}{环面 Yukawa 耦合},\enspace\IndexPageList[1456]{291,292,408}，\textit{另见} \hyperlink{idx:0190}{关联函数，B 模型关联函数，环面情形}
\subsubitem \leavevmode\hypertarget{idx:1457}{与归一化耦合的区别},\enspace\IndexPageList[1457]{291}
\indexspace
\item \leavevmode\hypertarget{idx:1458}{Zariski $p$ 形式},\enspace\IndexPageList[1458]{432}
\indexspace
\item \leavevmode\hypertarget{idx:0144}{凹丛},\enspace\IndexPageList[0144]{372}
\subitem \leavevmode\hypertarget{idx:0145}{多重覆盖情形},\enspace\IndexPageList[0145]{375}
\indexspace
\item \leavevmode\hypertarget{idx:0965}{八面体},\enspace\IndexPageList[0965]{74}
\indexspace
\item \leavevmode\hypertarget{idx:1244}{半丰沛},\enspace\IndexPageList[1244]{77}
\indexspace
\item \leavevmode\hypertarget{idx:1245}{半正},\enspace\IndexPageList[1245]{208,210,212,233}
\indexspace
\item \leavevmode\hypertarget{idx:0204}{保典范},\enspace\IndexPageList[0204]{77}
\indexspace
\item \leavevmode\hypertarget{idx:1087}{本原集合},\enspace\IndexPageList[1087]{61,73,246,247,248}
\subitem \leavevmode\hypertarget{idx:1088}{丰沛性判据},\enspace\IndexPageList[1088]{63}
\indexspace
\item \leavevmode\hypertarget{idx:1085}{本原上同调},\enspace\IndexPageList[1085]{106,109,276,287,384}
\subitem \leavevmode\hypertarget{idx:1086}{定义},\enspace\IndexPageList[1086]{97}
\indexspace
\item \leavevmode\hypertarget{idx:1089}{本原收缩},\enspace\IndexPageList[1089]{155,166,171}
\subitem \leavevmode\hypertarget{idx:1090}{I 型},\enspace\IndexPageList[1090]{156,157,166,167,170,249,250}
\subitem \leavevmode\hypertarget{idx:1091}{II 型},\enspace\IndexPageList[1091]{156,157,166}
\subitem \leavevmode\hypertarget{idx:1092}{III 型},\enspace\IndexPageList[1092]{156,157,166,169,250}
\indexspace
\item \leavevmode\hypertarget{idx:1334}{编时乘积},\enspace\IndexPageList[1334]{440}
\indexspace
\item \leavevmode\hypertarget{idx:1438}{变量的权},\enspace\IndexPageList[1438]{381,388,401,402}
\indexspace
\item \leavevmode\hypertarget{idx:1434}{表示的权},\enspace\IndexPageList[1434]{299,301,303,308,315,367,392,393}
\subitem \leavevmode\hypertarget{idx:1435}{障碍权},\enspace\IndexPageList[1435]{312,313,314,315,320}
\subitem \leavevmode\hypertarget{idx:1436}{法丛的权},\enspace\IndexPageList[1436]{308,312}
\subitem \leavevmode\hypertarget{idx:1437}{切权},\enspace\IndexPageList[1437]{301,304,308,310,313,320}
\indexspace
\item \leavevmode\hypertarget{idx:0081}{玻色}
\subitem \leavevmode\hypertarget{idx:0082}{玻色分量},\enspace\IndexPageList[0082]{450,452}
\subitem \leavevmode\hypertarget{idx:0083}{玻色场},\enspace\IndexPageList[0083]{442}
\subitem \leavevmode\hypertarget{idx:0084}{玻色变量},\enspace\IndexPageList[0084]{450,451}
\indexspace
\item \leavevmode\hypertarget{idx:0080}{玻色子},\enspace\IndexPageList[0080]{25,442,443,448}
\indexspace
\item \leavevmode\hypertarget{idx:0203}{产生算子},\enspace\IndexPageList[0203]{438,441}
\indexspace
\item \leavevmode\hypertarget{idx:0316}{场}
\subitem \leavevmode\hypertarget{idx:0317}{玻色场}，\textit{参见} \hyperlink{idx:0083}{玻色，玻色场}
\subitem \leavevmode\hypertarget{idx:0318}{经典场},\enspace\IndexPageList[0318]{434,437,439}
\subitem \leavevmode\hypertarget{idx:0319}{费米场}，\textit{参见} \hyperlink{idx:0312}{费米，费米场}
\subitem \leavevmode\hypertarget{idx:0320}{初级场},\enspace\IndexPageList[0320]{447}
\subitem \leavevmode\hypertarget{idx:0321}{量子场},\enspace\IndexPageList[0321]{437,439}
\subitem \leavevmode\hypertarget{idx:0322}{场论},\enspace\IndexPageList[0322]{434,437}，\textit{另见} \hyperlink{idx:0162}{共形场论}；\hyperlink{idx:1155}{量子场论}
\indexspace
\item \leavevmode\hypertarget{idx:0323}{场强},\enspace\IndexPageList[0323]{440,441}
\indexspace
\item \leavevmode\hypertarget{idx:0620}{超 Kähler},\enspace\IndexPageList[0620]{425}
\indexspace
\item \leavevmode\hypertarget{idx:1296}{超场},\enspace\IndexPageList[1296]{449,450,452}
\subitem \leavevmode\hypertarget{idx:1297}{手征超场}，\textit{参见} \hyperlink{idx:0121}{手征超场}
\indexspace
\item \leavevmode\hypertarget{idx:1309}{超对称场论},\enspace\IndexPageList[1309]{429,443,449}
\indexspace
\item \leavevmode\hypertarget{idx:1310}{超对称性},\enspace\IndexPageList[1310]{24,25,26,30,33,447,448,449}，\textit{另见} \hyperlink{idx:1254}{$\sigma$ 模型，超对称 $\sigma$ 模型}
\subitem \leavevmode\hypertarget{idx:1311}{$N=2$ 超对称代数},\enspace\IndexPageList[1311]{449}
\subitem \leavevmode\hypertarget{idx:1312}{超对称变换},\enspace\IndexPageList[1312]{25,443,445,449}
\indexspace
\item \leavevmode\hypertarget{idx:0280}{超额法丛},\enspace\IndexPageList[0280]{196,197}
\indexspace
\item \leavevmode\hypertarget{idx:0277}{超额维数},\enspace\IndexPageList[0277]{331}
\subitem \leavevmode\hypertarget{idx:0278}{稳定映射模空间的超额维数},\enspace\IndexPageList[0278]{205}
\subitem \leavevmode\hypertarget{idx:0279}{相交理论中的超额维数},\enspace\IndexPageList[0279]{205}
\indexspace
\item \leavevmode\hypertarget{idx:1292}{超共形}
\subitem \leavevmode\hypertarget{idx:1293}{$N=2$ 超共形代数},\enspace\IndexPageList[1293]{25}
\subitem \leavevmode\hypertarget{idx:1294}{$N=2$ 超共形表示},\enspace\IndexPageList[1294]{79}
\subitem \leavevmode\hypertarget{idx:1295}{超共形场论（SCFT）},\enspace\IndexPageList[1295]{25,27,29,30,32,33,35,40,76,79}
\indexspace
\item \leavevmode\hypertarget{idx:0615}{超几何方程},\enspace\IndexPageList[0615]{113,114,115,116,117,118,119,120,121,122,123,124,184,356}，\textit{另见} \hyperlink{idx:0055}{$\mathcal A$ 方程组}；\hyperlink{idx:0449}{GKZ 方程组}
\indexspace
\item \leavevmode\hypertarget{idx:0616}{超几何函数},\enspace\IndexPageList[0616]{118,187}
\subitem \leavevmode\hypertarget{idx:0617}{$p$ 进超几何函数},\enspace\IndexPageList[0617]{43}
\subitem \leavevmode\hypertarget{idx:0618}{广义超几何函数},\enspace\IndexPageList[0618]{43}
\subitem \leavevmode\hypertarget{idx:0619}{$HG$}，\textit{参见} \hyperlink{idx:0579}{$HG$}
\indexspace
\item \leavevmode\hypertarget{idx:1291}{超交换},\enspace\IndexPageList[1291]{242,244,252,253,256,257,263,457}
\indexspace
\item \leavevmode\hypertarget{idx:1299}{超流形},\enspace\IndexPageList[1299]{252,253,262,449}
\subitem \leavevmode\hypertarget{idx:1300}{超流形上的偏导数},\enspace\IndexPageList[1300]{252,327}
\subitem \leavevmode\hypertarget{idx:1301}{切丛},\enspace\IndexPageList[1301]{263}
\indexspace
\item \leavevmode\hypertarget{idx:1302}{超势},\enspace\IndexPageList[1302]{449,451,453}
\subitem \leavevmode\hypertarget{idx:1303}{五次超势},\enspace\IndexPageList[1303]{452}
\indexspace
\item \leavevmode\hypertarget{idx:1304}{超弦理论}
\subitem \leavevmode\hypertarget{idx:1305}{对偶性},\enspace\IndexPageList[1305]{427}
\subitem \leavevmode\hypertarget{idx:1306}{杂化超弦理论},\enspace\IndexPageList[1306]{25,26,427,443,444}
\subitem \leavevmode\hypertarget{idx:1307}{I 型超弦理论},\enspace\IndexPageList[1307]{25}
\subitem \leavevmode\hypertarget{idx:1308}{II 型超弦理论},\enspace\IndexPageList[1308]{25,168,427}
\indexspace
\item \leavevmode\hypertarget{idx:1313}{超向量空间},\enspace\IndexPageList[1313]{454,457}
\indexspace
\item \leavevmode\hypertarget{idx:1298}{超引力},\enspace\IndexPageList[1298]{33}
\indexspace
\item \leavevmode\hypertarget{idx:1242}{次级多胞形},\enspace\IndexPageList[1242]{143}
\indexspace
\item \leavevmode\hypertarget{idx:1239}{次级扇},\enspace\IndexPageList[1239]{66,140,143,144,158,159,165,166,171,178,179}，\textit{另见} \hyperlink{idx:0446}{GKZ 分解}
\subitem \leavevmode\hypertarget{idx:1240}{扩大次级扇},\enspace\IndexPageList[1240]{453}
\subitem \leavevmode\hypertarget{idx:1241}{例题},\enspace\IndexPageList[1241]{73,74,147,149,163}
\indexspace
\item \leavevmode\hypertarget{idx:0696}{大半径极限点},\enspace\IndexPageList[0696]{138,151,152,153,154,159,161,172,173,175,176,178,179,268,274,280,287,290,292}
\subitem \leavevmode\hypertarget{idx:0697}{等价性},\enspace\IndexPageList[0697]{154,175}
\subitem \leavevmode\hypertarget{idx:0698}{例题},\enspace\IndexPageList[0698]{164}
\subitem \leavevmode\hypertarget{idx:0699}{该点处的切空间},\enspace\IndexPageList[0699]{181}
\indexspace
\item \leavevmode\hypertarget{idx:0040}{代数叠},\enspace\IndexPageList[0040]{192}
\subitem \leavevmode\hypertarget{idx:0041}{代数叠上的丛}，\textit{参见} \hyperlink{idx:0261}{Euler 类，代数叠与粗模空间之间的比较}
\subitem \leavevmode\hypertarget{idx:0042}{Deligne--Mumford 代数叠},\enspace\IndexPageList[0042]{201}
\subitem \leavevmode\hypertarget{idx:0043}{代数叠上的相交理论},\enspace\IndexPageList[0043]{201,302}
\subitem \leavevmode\hypertarget{idx:0044}{商叠},\enspace\IndexPageList[0044]{201}
\subitem \leavevmode\hypertarget{idx:0045}{光滑代数叠},\enspace\IndexPageList[0045]{193,302}
\indexspace
\item \leavevmode\hypertarget{idx:0039}{代数空间},\enspace\IndexPageList[0039]{193}
\indexspace
\item \leavevmode\hypertarget{idx:0910}{单项式—除子镜像映射},\enspace\IndexPageList[0910]{80,83,181,291,294,295}
\indexspace
\item \leavevmode\hypertarget{idx:0909}{单项式序},\enspace\IndexPageList[0909]{248}
\indexspace
\item \leavevmode\hypertarget{idx:0901}{单值性变换},\enspace\IndexPageList[0901]{98,100,104,137,339}
\subitem \leavevmode\hypertarget{idx:0902}{对数},\enspace\IndexPageList[0902]{98,104,134,137,165,269,270,272,339}
\subitem \leavevmode\hypertarget{idx:0903}{A 模型联络的单值性变换}，\textit{参见} \hyperlink{idx:0011}{A 模型，联络，单值性}
\subitem \leavevmode\hypertarget{idx:0904}{拟幂幺单值性变换},\enspace\IndexPageList[0904]{98}
\subitem \leavevmode\hypertarget{idx:0905}{幂幺单值性变换},\enspace\IndexPageList[0905]{98,100,101,137,269,270}
\subsubitem \leavevmode\hypertarget{idx:0906}{最大幂幺}，\textit{参见} \hyperlink{idx:0765}{最大幂幺}
\indexspace
\item \leavevmode\hypertarget{idx:0900}{单值性定理},\enspace\IndexPageList[0900]{98}
\indexspace
\item \leavevmode\hypertarget{idx:0907}{单值性权滤过},\enspace\IndexPageList[0907]{99,103,104,127,128,138,172,173,177,184,273,291}
\subitem \leavevmode\hypertarget{idx:0908}{A 模型联络的单值性权滤过}，\textit{参见} \hyperlink{idx:0012}{A 模型，联络，单值性权滤过}
\indexspace
\item \leavevmode\hypertarget{idx:0243}{等变 Chern 类},\enspace\IndexPageList[0243]{299,303,389}
\indexspace
\item \leavevmode\hypertarget{idx:0250}{等变超平面类},\enspace\IndexPageList[0250]{51,124,358,369,389,391}
\subitem \leavevmode\hypertarget{idx:0251}{定义},\enspace\IndexPageList[0251]{300}
\subitem \leavevmode\hypertarget{idx:0252}{$N_d$ 的等变超平面类}，\textit{参见} \hyperlink{idx:0930}{$N_d$，射影空间，等变超平面类}
\indexspace
\item \leavevmode\hypertarget{idx:0253}{等变积分},\enspace\IndexPageList[0253]{303,358,390,391,396,399,400}
\subitem \leavevmode\hypertarget{idx:0254}{定义},\enspace\IndexPageList[0254]{302}
\indexspace
\item \leavevmode\hypertarget{idx:0244}{等变上同调},\enspace\IndexPageList[0244]{50,124}
\subitem \leavevmode\hypertarget{idx:0245}{定义},\enspace\IndexPageList[0245]{299}
\subitem \leavevmode\hypertarget{idx:0246}{参数 $\hbar$ 的说明},\enspace\IndexPageList[0246]{342}
\subitem \leavevmode\hypertarget{idx:0247}{局部化},\enspace\IndexPageList[0247]{300,359}，\textit{另见} \hyperlink{idx:0739}{局部化定理，等变上同调的局部化定理}
\subitem \leavevmode\hypertarget{idx:0248}{$\mathbb P^n$ 的等变上同调},\enspace\IndexPageList[0248]{300,301,358}
\subitem \leavevmode\hypertarget{idx:0249}{限制映射},\enspace\IndexPageList[0249]{366,367}
\indexspace
\item \leavevmode\hypertarget{idx:0242}{等变特征类},\enspace\IndexPageList[0242]{361}
\indexspace
\item \leavevmode\hypertarget{idx:0255}{等变向量丛},\enspace\IndexPageList[0255]{299,321}
\indexspace
\item \leavevmode\hypertarget{idx:0745}{低能极限},\enspace\IndexPageList[0745]{33}
\indexspace
\item \leavevmode\hypertarget{idx:1210}{递推关系}
\subitem \leavevmode\hypertarget{idx:1211}{微分方程的递推关系},\enspace\IndexPageList[1211]{118,185,186}
\subitem \leavevmode\hypertarget{idx:1212}{Gromov--Witten 不变量的递推关系}，\textit{参见} \hyperlink{idx:0523}{Gromov--Witten 不变量，计算}
\indexspace
\item \leavevmode\hypertarget{source:489}{}\hypertarget{idx:1209}{递推引理}，\textit{参见} \hyperlink{idx:0860}{镜像定理，Givental 镜像定理，递推引理}
\indexspace
\item \leavevmode\hypertarget{idx:0112}{典范延拓},\enspace\IndexPageList[0112]{98,177,270,272,283}
\subitem \leavevmode\hypertarget{idx:0113}{A 模型联络的典范延拓}，\textit{参见} \hyperlink{idx:0005}{A 模型，联络，典范延拓}
\indexspace
\item \leavevmode\hypertarget{idx:0238}{电磁势},\enspace\IndexPageList[0238]{440}
\indexspace
\item \leavevmode\hypertarget{idx:1283}{叠}，\textit{参见} \hyperlink{idx:0040}{代数叠}
\indexspace
\item \leavevmode\hypertarget{idx:0225}{对偶化层},\enspace\IndexPageList[0225]{34,69,430,432}
\indexspace
\item \leavevmode\hypertarget{idx:0224}{对偶锥},\enspace\IndexPageList[0224]{54}
\indexspace
\item \leavevmode\hypertarget{idx:1071}{多胞形},\enspace\IndexPageList[1071]{56}
\subitem \leavevmode\hypertarget{idx:1072}{刻面},\enspace\IndexPageList[1072]{56}
\subitem \leavevmode\hypertarget{idx:1073}{整多胞形},\enspace\IndexPageList[1073]{56}
\subitem \leavevmode\hypertarget{idx:1074}{Newton 多胞形},\enspace\IndexPageList[1074]{58,83}
\subitem \leavevmode\hypertarget{idx:1075}{法扇},\enspace\IndexPageList[1075]{57,74,76,158}
\subitem \leavevmode\hypertarget{idx:1076}{归一化体积},\enspace\IndexPageList[1076]{82}
\subitem \leavevmode\hypertarget{idx:1077}{极多胞形},\enspace\IndexPageList[1077]{57,69,74,76,143}
\subsubitem \leavevmode\hypertarget{idx:1078}{极对偶},\enspace\IndexPageList[1078]{58}
\subitem \leavevmode\hypertarget{idx:1079}{自反多胞形},\enspace\IndexPageList[1079]{83}，\textit{参见} \hyperlink{idx:1218}{自反多胞形}
\subitem \leavevmode\hypertarget{idx:1080}{多胞形环},\enspace\IndexPageList[1080]{57,59,60}
\subitem \leavevmode\hypertarget{idx:1081}{单纯多胞形},\enspace\IndexPageList[1081]{432}
\subitem \leavevmode\hypertarget{idx:1082}{对应的环面簇},\enspace\IndexPageList[1082]{57}
\indexspace
\item \leavevmode\hypertarget{idx:0915}{多重覆盖},\enspace\IndexPageList[0915]{225,230,232,311,366}，\textit{另见} \hyperlink{idx:0628}{瞬子数}
\subitem \leavevmode\hypertarget{idx:0916}{对 Gromov--Witten 不变量的贡献},\enspace\IndexPageList[0916]{225,227,311,312,313,314,315,316,317,375}
\subsubitem \leavevmode\hypertarget{idx:0917}{Lian--Liu--Yau 的证明},\enspace\IndexPageList[0917]{378}
\subsubitem \leavevmode\hypertarget{idx:0918}{Pandharipande 的证明},\enspace\IndexPageList[0918]{314,315,316}
\subsubitem \leavevmode\hypertarget{idx:0919}{$d=2$ 时的证明},\enspace\IndexPageList[0919]{312,313,314}
\subitem \leavevmode\hypertarget{idx:0920}{对虚拟基本类的贡献},\enspace\IndexPageList[0920]{225,228,317}
\subitem \leavevmode\hypertarget{idx:0921}{椭圆多重覆盖},\enspace\IndexPageList[0921]{232,237,323}
\subitem \leavevmode\hypertarget{idx:0922}{六结点五次曲线的多重覆盖},\enspace\IndexPageList[0922]{316,317}，\textit{另见} \hyperlink{idx:0218}{六结点五次曲线的二重覆盖}
\subitem \leavevmode\hypertarget{idx:0923}{非孤立曲线的多重覆盖},\enspace\IndexPageList[0923]{231}
\indexspace
\item \leavevmode\hypertarget{idx:0924}{多重镜像},\enspace\IndexPageList[0924]{68,83,144}
\indexspace
\item \leavevmode\hypertarget{idx:0937}{二次曲面网},\enspace\IndexPageList[0937]{389}
\indexspace
\item \leavevmode\hypertarget{idx:0947}{法锥},\enspace\IndexPageList[0947]{195,197,200,201,204,383}
\subitem \leavevmode\hypertarget{idx:0948}{内蕴法锥},\enspace\IndexPageList[0948]{201,202}
\subsubitem \leavevmode\hypertarget{idx:0949}{例题},\enspace\IndexPageList[0949]{202}
\subitem \leavevmode\hypertarget{idx:0950}{细化 Euler 类},\enspace\IndexPageList[0950]{196}
\subitem \leavevmode\hypertarget{idx:0951}{虚法锥},\enspace\IndexPageList[0951]{200}
\indexspace
\item \leavevmode\hypertarget{idx:0332}{翻越},\enspace\IndexPageList[0332]{66,156,170,248,249}
\subitem \leavevmode\hypertarget{idx:0333}{例题},\enspace\IndexPageList[0333]{164}
\subitem \leavevmode\hypertarget{idx:0334}{广义翻越},\enspace\IndexPageList[0334]{160}
\subitem \leavevmode\hypertarget{idx:0335}{平凡翻越},\enspace\IndexPageList[0335]{160}
\indexspace
\item \leavevmode\hypertarget{idx:0328}{翻转},\enspace\IndexPageList[0328]{160,162}
\subitem \leavevmode\hypertarget{idx:0329}{例题},\enspace\IndexPageList[0329]{164}
\subitem \leavevmode\hypertarget{idx:0330}{几何意义},\enspace\IndexPageList[0330]{160}，\textit{另见} \hyperlink{idx:0334}{翻越，广义翻越}
\subitem \leavevmode\hypertarget{idx:0331}{平凡翻转},\enspace\IndexPageList[0331]{160,161,164,165,179}
\indexspace
\item \leavevmode\hypertarget{idx:0050}{反典范}
\subitem \leavevmode\hypertarget{idx:0051}{反典范除子},\enspace\IndexPageList[0051]{69,71,72,74,77}
\subitem \leavevmode\hypertarget{idx:0052}{反典范线性系},\enspace\IndexPageList[0052]{78,79}
\indexspace
\item \leavevmode\hypertarget{idx:0053}{反对易子},\enspace\IndexPageList[0053]{445}
\indexspace
\item \leavevmode\hypertarget{idx:0114}{范畴商},\enspace\IndexPageList[0114]{60}
\indexspace
\item \leavevmode\hypertarget{idx:0945}{非等变极限},\enspace\IndexPageList[0945]{51,369,373,374,375,376,377,380,389,395,402,403}
\subitem \leavevmode\hypertarget{idx:0946}{定义},\enspace\IndexPageList[0946]{299}
\indexspace
\item \leavevmode\hypertarget{idx:0943}{非微扰态},\enspace\IndexPageList[0943]{425}
\indexspace
\item \leavevmode\hypertarget{idx:0942}{非线性 $\sigma$ 模型}，\textit{参见} \hyperlink{idx:1249}{$\sigma$ 模型}
\indexspace
\item \leavevmode\hypertarget{idx:0944}{非重整化定理},\enspace\IndexPageList[0944]{30,31,446}
\indexspace
\item \leavevmode\hypertarget{idx:0311}{费米}
\subitem \leavevmode\hypertarget{idx:0312}{费米场},\enspace\IndexPageList[0312]{442,443}，\textit{另见} \hyperlink{idx:1253}{$\sigma$ 模型，费米场}
\subitem \leavevmode\hypertarget{idx:0313}{费米变量},\enspace\IndexPageList[0313]{442,450,451}
\indexspace
\item \leavevmode\hypertarget{idx:0309}{费米子},\enspace\IndexPageList[0309]{25,442,443,444,448}
\subitem \leavevmode\hypertarget{idx:0310}{与超交换变量的关系},\enspace\IndexPageList[0310]{253,442}
\indexspace
\item \leavevmode\hypertarget{idx:0139}{复反射},\enspace\IndexPageList[0139]{431}
\indexspace
\item \leavevmode\hypertarget{idx:0140}{复化}
\subitem \leavevmode\hypertarget{idx:0141}{复化 Kähler 类},\enspace\IndexPageList[0141]{26,29,30,31,32,34,35,38,41,45,227,240,241,248,251,253,326,327,457}
\subitem \leavevmode\hypertarget{idx:0142}{复化 Kähler 模}，\textit{参见} \hyperlink{idx:0882}{模，Kähler 模}
\subitem \leavevmode\hypertarget{idx:0143}{复化 Kähler 空间},\enspace\IndexPageList[0143]{35,41,150,151,175,266,268,274}
\indexspace
\item \leavevmode\hypertarget{idx:0276}{赋值映射},\enspace\IndexPageList[0276]{195,202,212,228,233,245,321,324,355,380,383,388,398}
\indexspace
\item \leavevmode\hypertarget{idx:0171}{共轭动量},\enspace\IndexPageList[0171]{435,437}
\indexspace
\item \leavevmode\hypertarget{idx:0167}{共形不变性},\enspace\IndexPageList[0167]{443,447}
\indexspace
\item \leavevmode\hypertarget{idx:0162}{共形场论},\enspace\IndexPageList[0162]{446,447,448,449}
\subitem \leavevmode\hypertarget{idx:0163}{规范的选择},\enspace\IndexPageList[0163]{447}
\subitem \leavevmode\hypertarget{idx:0164}{规范群},\enspace\IndexPageList[0164]{447}
\subitem \leavevmode\hypertarget{idx:0165}{$\sigma$ 模型的共形场论},\enspace\IndexPageList[0165]{448}
\indexspace
\item \leavevmode\hypertarget{idx:0161}{共形代数},\enspace\IndexPageList[0161]{25,446,447}
\indexspace
\item \leavevmode\hypertarget{idx:0166}{共形群},\enspace\IndexPageList[0166]{447,448}
\indexspace
\item \leavevmode\hypertarget{idx:0183}{关联函数},\enspace\IndexPageList[0183]{24,29,79,175,181,264,283,326,387,439,444}
\subitem \leavevmode\hypertarget{idx:0184}{A 模型关联函数},\enspace\IndexPageList[0184]{30,31,32,34,130,226,282,289,291,408,410,412,445,446}，\textit{另见} \hyperlink{idx:1183}{五次三维簇，A 模型关联函数}
\subsubitem \leavevmode\hypertarget{idx:0185}{定义},\enspace\IndexPageList[0185]{244}
\subsubitem \leavevmode\hypertarget{idx:0186}{翻越的影响},\enspace\IndexPageList[0186]{250}
\subsubitem \leavevmode\hypertarget{idx:0187}{用瞬子数表示},\enspace\IndexPageList[0187]{249}
\subsubitem \leavevmode\hypertarget{idx:0188}{Calabi--Yau 三维簇的关联函数},\enspace\IndexPageList[0188]{227,249,278,279,350}
\subitem \leavevmode\hypertarget{idx:0189}{B 模型关联函数},\enspace\IndexPageList[0189]{31,32,133,279,282,288,291,408,410,446}，\textit{另见} \hyperlink{idx:0064}{B 模型，势函数}；\hyperlink{idx:1444}{Yukawa 耦合}；\hyperlink{idx:1164}{五次镜像}
\subsubitem \leavevmode\hypertarget{idx:0190}{环面情形},\enspace\IndexPageList[0190]{291}
\subitem \leavevmode\hypertarget{idx:0191}{亏格 $g$ 的关联函数},\enspace\IndexPageList[0191]{457}
\subitem \leavevmode\hypertarget{idx:0192}{$n$ 点关联函数},\enspace\IndexPageList[0192]{439,457}
\subitem \leavevmode\hypertarget{idx:0193}{Frobenius 代数的关联函数},\enspace\IndexPageList[0193]{455}
\subitem \leavevmode\hypertarget{idx:0194}{五次情形}，\textit{参见} \hyperlink{idx:1164}{五次镜像}；\hyperlink{idx:1181}{五次三维簇}
\subitem \leavevmode\hypertarget{idx:0195}{三点关联函数},\enspace\IndexPageList[0195]{29,30,31,32,33,244,249,250,278,279,445}
\subitem \leavevmode\hypertarget{idx:0196}{Voisin--Borcea 镜像构造的关联函数},\enspace\IndexPageList[0196]{92}
\indexspace
\item \leavevmode\hypertarget{idx:1383}{管状域},\enspace\IndexPageList[1383]{153}
\indexspace
\item \leavevmode\hypertarget{idx:0371}{广义 Berglund--H\"ubsch 转置法则},\enspace\IndexPageList[0371]{93}
\indexspace
\item \leavevmode\hypertarget{idx:0373}{广义 Fredholm 轨形丛},\enspace\IndexPageList[0373]{214}
\subitem \leavevmode\hypertarget{idx:0374}{定向 Euler 类},\enspace\IndexPageList[0374]{214}
\indexspace
\item \leavevmode\hypertarget{idx:0372}{广义翻越}，\textit{参见} \hyperlink{idx:0334}{翻越，广义翻越}
\indexspace
\item \leavevmode\hypertarget{idx:0953}{归一化 $3$ 形式},\enspace\IndexPageList[0953]{44,127,128,129,133,135,174,281,282,291,351,408}，\textit{另见} \hyperlink{idx:1445}{Yukawa 耦合，归一化 Yukawa 耦合}
\indexspace
\item \leavevmode\hypertarget{idx:0954}{归一化 $d$ 形式},\enspace\IndexPageList[0954]{177}
\indexspace
\item \leavevmode\hypertarget{idx:0350}{规范固定},\enspace\IndexPageList[0350]{440}
\indexspace
\item \leavevmode\hypertarget{idx:0353}{规范理论},\enspace\IndexPageList[0353]{440,451}
\indexspace
\item \leavevmode\hypertarget{idx:0351}{规范群},\enspace\IndexPageList[0351]{440,442,443,447,452}
\subitem \leavevmode\hypertarget{idx:0352}{电磁理论的规范群},\enspace\IndexPageList[0352]{441}
\indexspace
\item \leavevmode\hypertarget{idx:0354}{规范线性 $\sigma$ 模型},\enspace\IndexPageList[0354]{35,65,68,357,451,452,453}
\subitem \leavevmode\hypertarget{idx:0355}{作用量},\enspace\IndexPageList[0355]{451}
\subitem \leavevmode\hypertarget{idx:0356}{玻色部分},\enspace\IndexPageList[0356]{451}
\subitem \leavevmode\hypertarget{idx:0357}{关联函数},\enspace\IndexPageList[0357]{387}
\subitem \leavevmode\hypertarget{idx:0358}{例题},\enspace\IndexPageList[0358]{452}
\subitem \leavevmode\hypertarget{idx:0359}{规范群},\enspace\IndexPageList[0359]{451}
\subitem \leavevmode\hypertarget{idx:0360}{镜像对称},\enspace\IndexPageList[0360]{453}
\subitem \leavevmode\hypertarget{idx:0361}{相},\enspace\IndexPageList[0361]{166}，\textit{另见} \hyperlink{idx:1042}{相}
\subsubitem \leavevmode\hypertarget{idx:0362}{Calabi--Yau 相},\enspace\IndexPageList[0362]{452}
\subsubitem \leavevmode\hypertarget{idx:0363}{Landau--Ginzburg 相},\enspace\IndexPageList[0363]{453}
\subitem \leavevmode\hypertarget{idx:0364}{与环面几何的关系},\enspace\IndexPageList[0364]{452}
\indexspace
\item \leavevmode\hypertarget{idx:0349}{规范选择},\enspace\IndexPageList[0349]{441,442,447}
\indexspace
\item \leavevmode\hypertarget{idx:0990}{轨形},\enspace\IndexPageList[0990]{35,55,61,77,194,430}
\subitem \leavevmode\hypertarget{idx:0991}{Calabi--Yau 轨形}，\textit{参见} \hyperlink{idx:0095}{Calabi--Yau，Calabi--Yau 轨形}
\subitem \leavevmode\hypertarget{idx:0992}{定义},\enspace\IndexPageList[0992]{431}
\subitem \leavevmode\hypertarget{idx:0993}{微分同胚},\enspace\IndexPageList[0993]{65,68,433}
\subitem \leavevmode\hypertarget{idx:0994}{Gorenstein 轨形},\enspace\IndexPageList[0994]{430,431,433}
\subitem \leavevmode\hypertarget{idx:0995}{轨形上的 Hodge 理论},\enspace\IndexPageList[0995]{432}
\subitem \leavevmode\hypertarget{idx:0996}{无穷小形变},\enspace\IndexPageList[0996]{433}
\subitem \leavevmode\hypertarget{idx:0997}{轨形上的积分},\enspace\IndexPageList[0997]{363}
\subitem \leavevmode\hypertarget{idx:0998}{轨形上的 Kähler 形式},\enspace\IndexPageList[0998]{432}
\subitem \leavevmode\hypertarget{idx:0999}{辛轨形},\enspace\IndexPageList[0999]{65,433}
\subitem \leavevmode\hypertarget{idx:1000}{切层},\enspace\IndexPageList[1000]{433}
\indexspace
\item \leavevmode\hypertarget{idx:0380}{鬼场},\enspace\IndexPageList[0380]{447}
\indexspace
\item \leavevmode\hypertarget{idx:0078}{黑洞},\enspace\IndexPageList[0078]{168,427}
\subitem \leavevmode\hypertarget{idx:0079}{无质量黑洞},\enspace\IndexPageList[0079]{168}
\indexspace
\item \leavevmode\hypertarget{idx:1207}{互易引理},\enspace\IndexPageList[1207]{367}
\indexspace
\item \leavevmode\hypertarget{idx:1353}{环面簇},\enspace\IndexPageList[1353]{55,405,430}
\subitem \leavevmode\hypertarget{idx:1354}{仿射环面簇},\enspace\IndexPageList[1354]{55,152}
\subitem \leavevmode\hypertarget{idx:1355}{自同构群},\enspace\IndexPageList[1355]{70,71,81,84,85,139,140,141,148}
\subsubitem \leavevmode\hypertarget{idx:1356}{维数},\enspace\IndexPageList[1356]{71}
\subsubitem \leavevmode\hypertarget{idx:1357}{根},\enspace\IndexPageList[1357]{71,147,149}
\subitem \leavevmode\hypertarget{idx:1358}{Batyrev 量子环}，\textit{参见} \hyperlink{idx:0069}{Batyrev 量子环}
\subitem \leavevmode\hypertarget{idx:1359}{Chow 群},\enspace\IndexPageList[1359]{55,59}
\subitem \leavevmode\hypertarget{idx:1360}{上同调环},\enspace\IndexPageList[1360]{246,415}
\subitem \leavevmode\hypertarget{idx:1361}{完备环面簇},\enspace\IndexPageList[1361]{55}
\subitem \leavevmode\hypertarget{idx:1362}{环面簇上的除子},\enspace\IndexPageList[1362]{56}
\subsubitem \leavevmode\hypertarget{idx:1363}{丰沛除子},\enspace\IndexPageList[1363]{56,58,63}
\subsubitem \leavevmode\hypertarget{idx:1364}{由整体截面生成的除子},\enspace\IndexPageList[1364]{56}
\subsubitem \leavevmode\hypertarget{idx:1365}{整体截面},\enspace\IndexPageList[1365]{57}
\subsubitem \leavevmode\hypertarget{idx:1366}{半丰沛除子}，\textit{参见} \hyperlink{idx:1244}{半丰沛}
\subitem \leavevmode\hypertarget{idx:1367}{例题},\enspace\IndexPageList[1367]{72,73,74,75}
\subitem \leavevmode\hypertarget{idx:1368}{Fano 环面簇},\enspace\IndexPageList[1368]{69,70,72,74,76,412,413,414,415,416,418}，\textit{另见} \hyperlink{idx:1218}{自反多胞形}
\subitem \leavevmode\hypertarget{idx:1369}{齐次坐标环},\enspace\IndexPageList[1369]{59,70,71,72,73,108,115,140,148,453}
\subitem \leavevmode\hypertarget{idx:1370}{理想与簇的对应},\enspace\IndexPageList[1370]{61}
\subitem \leavevmode\hypertarget{idx:1371}{Kähler 锥},\enspace\IndexPageList[1371]{61,62,64,415,416,452}
\subitem \leavevmode\hypertarget{idx:1372}{Mori 锥},\enspace\IndexPageList[1372]{62,63,406,412,416}
\subitem \leavevmode\hypertarget{idx:1373}{Picard 群},\enspace\IndexPageList[1373]{55,247}
\subitem \leavevmode\hypertarget{idx:1374}{量子上同调}，\textit{参见} \hyperlink{idx:1141}{量子上同调，小量子上同调，光滑环面簇的小量子上同调}
\subitem \leavevmode\hypertarget{idx:1375}{单纯环面簇},\enspace\IndexPageList[1375]{55,60,61,64,70,71,431}
\subitem \leavevmode\hypertarget{idx:1376}{光滑环面簇},\enspace\IndexPageList[1376]{55,404}
\subitem \leavevmode\hypertarget{idx:1377}{用扇描述}，\textit{参见} \hyperlink{idx:0307}{扇，扇的环面簇}
\subitem \leavevmode\hypertarget{idx:1378}{用齐次坐标描述},\enspace\IndexPageList[1378]{60}
\subitem \leavevmode\hypertarget{idx:1379}{用多胞形描述}，\textit{参见} \hyperlink{idx:1082}{多胞形，对应的环面簇}
\subitem \leavevmode\hypertarget{idx:1380}{用辛约化描述},\enspace\IndexPageList[1380]{64,65,68,453}
\indexspace
\item \leavevmode\hypertarget{idx:1350}{环面镜像定理}，\textit{参见} \hyperlink{idx:0862}{镜像定理，Givental 镜像定理，环面镜像定理}
\indexspace
\item \leavevmode\hypertarget{idx:1347}{环面完全交},\enspace\IndexPageList[1347]{29,85}
\subitem \leavevmode\hypertarget{idx:1348}{Calabi--Yau 环面完全交},\enspace\IndexPageList[1348]{29,85,405,453}
\subitem \leavevmode\hypertarget{idx:1349}{nef 环面完全交},\enspace\IndexPageList[1349]{119,120,151,404,405,412}
\indexspace
\item \leavevmode\hypertarget{idx:0347}{基本类},\enspace\IndexPageList[0347]{195}
\indexspace
\item \leavevmode\hypertarget{idx:0763}{极大射影保典范部分解消},\enspace\IndexPageList[0763]{78}
\indexspace
\item \leavevmode\hypertarget{idx:0764}{极大射影细分}，\textit{参见} \hyperlink{idx:1096}{射影细分，极大射影细分}
\indexspace
\item \leavevmode\hypertarget{idx:0285}{极端跃迁},\enspace\IndexPageList[0285]{166,167,169,428,429}
\subitem \leavevmode\hypertarget{idx:0286}{与镜像对称的相容性},\enspace\IndexPageList[0286]{167,171}
\subitem \leavevmode\hypertarget{idx:0287}{对偶跃迁},\enspace\IndexPageList[0287]{170,171,181,182,428}
\subitem \leavevmode\hypertarget{idx:0288}{例题},\enspace\IndexPageList[0288]{169,170,171,181}
\subitem \leavevmode\hypertarget{idx:0289}{镜像跃迁}，\textit{参见} \hyperlink{idx:0287}{极端跃迁，对偶跃迁}
\subitem \leavevmode\hypertarget{idx:0290}{与超共形场论模的关系},\enspace\IndexPageList[0290]{167}
\subitem \leavevmode\hypertarget{idx:0291}{环面情形},\enspace\IndexPageList[0291]{168}
\subitem \leavevmode\hypertarget{idx:0292}{I 型},\enspace\IndexPageList[0292]{166,168}
\subitem \leavevmode\hypertarget{idx:0293}{II 型},\enspace\IndexPageList[0293]{166}
\subitem \leavevmode\hypertarget{idx:0294}{III 型},\enspace\IndexPageList[0294]{166,168}
\indexspace
\item \leavevmode\hypertarget{idx:0716}{极限权滤过}，\textit{参见} \hyperlink{idx:0907}{单值性权滤过}
\indexspace
\item \leavevmode\hypertarget{idx:0784}{极小模型},\enspace\IndexPageList[0784]{29}
\indexspace
\item \leavevmode\hypertarget{idx:0379}{几何工程理论},\enspace\IndexPageList[0379]{429}
\indexspace
\item \leavevmode\hypertarget{idx:0378}{几何商},\enspace\IndexPageList[0378]{60,140}
\indexspace
\item \leavevmode\hypertarget{idx:0046}{几乎复结构},\enspace\IndexPageList[0046]{228,231}
\subitem \leavevmode\hypertarget{idx:0047}{相容的几乎复结构},\enspace\IndexPageList[0047]{208,211}
\subitem \leavevmode\hypertarget{idx:0048}{受驯服的几乎复结构},\enspace\IndexPageList[0048]{208}
\indexspace
\item \leavevmode\hypertarget{idx:1431}{加权齐次多项式},\enspace\IndexPageList[1431]{449,450}
\indexspace
\item \leavevmode\hypertarget{idx:1432}{加权射影空间},\enspace\IndexPageList[1432]{29,72,73,76,93,94,163}
\subitem \leavevmode\hypertarget{idx:1433}{成为 Fano 空间的判据},\enspace\IndexPageList[1433]{70,73}
\indexspace
\item \leavevmode\hypertarget{idx:0638}{交同调},\enspace\IndexPageList[0638]{432}
\indexspace
\item \leavevmode\hypertarget{idx:0940}{结点形变},\enspace\IndexPageList[0940]{364}
\indexspace
\item \leavevmode\hypertarget{idx:0868}{镜像变换}，\textit{参见} \hyperlink{idx:0268}{Euler 数据，相关联的 Euler 数据，镜像变换}；\hyperlink{idx:0581}{$HG$，镜像变换}
\indexspace
\item \leavevmode\hypertarget{source:485}{}\hypertarget{idx:0786}{镜像猜想},\enspace\IndexPageList[0786]{351}
\subitem \leavevmode\hypertarget{idx:0787}{Calabi--Yau 三维簇的镜像猜想},\enspace\IndexPageList[0787]{281,284}
\subitem \leavevmode\hypertarget{idx:0788}{五次三维簇的镜像猜想},\enspace\IndexPageList[0788]{285,356}
\subitem \leavevmode\hypertarget{idx:0789}{高维镜像猜想},\enspace\IndexPageList[0789]{290}
\subitem \leavevmode\hypertarget{idx:0790}{与 Givental 镜像定理的关系},\enspace\IndexPageList[0790]{406}
\subitem \leavevmode\hypertarget{idx:0791}{环面镜像猜想},\enspace\IndexPageList[0791]{135,290,291,292,293,294,295,296,406}，\textit{另见} \hyperlink{idx:0862}{镜像定理，Givental 镜像定理，环面镜像定理}
\subsubitem \leavevmode\hypertarget{idx:0792}{例题},\enspace\IndexPageList[0792]{292,293}
\subsubitem \leavevmode\hypertarget{idx:0793}{三维簇的环面镜像猜想},\enspace\IndexPageList[0793]{290}
\subsubitem \leavevmode\hypertarget{idx:0794}{高维环面镜像猜想},\enspace\IndexPageList[0794]{296}
\subsubitem \leavevmode\hypertarget{idx:0795}{Hodge 理论环面镜像猜想},\enspace\IndexPageList[0795]{79,135,291,295,296,410}
\subsubitem \leavevmode\hypertarget{idx:0796}{数学环面镜像猜想},\enspace\IndexPageList[0796]{294}
\subsubitem \leavevmode\hypertarget{idx:0797}{表述},\enspace\IndexPageList[0797]{291,295}
\indexspace
\item \leavevmode\hypertarget{idx:0844}{镜像定理},\enspace\IndexPageList[0844]{29,296,297}
\subitem \leavevmode\hypertarget{idx:0845}{更广泛的理解},\enspace\IndexPageList[0845]{378,421,426}
\subitem \leavevmode\hypertarget{idx:0846}{Givental 镜像定理},\enspace\IndexPageList[0846]{379,421}
\subsubitem \leavevmode\hypertarget{idx:0847}{$\mathcal P$ 类},\enspace\IndexPageList[0847]{401,402,403}
\subsubitem \leavevmode\hypertarget{idx:0848}{计算亏格 $0$ 引力关联数},\enspace\IndexPageList[0848]{384}
\subsubitem \leavevmode\hypertarget{idx:0849}{坐标变换引理},\enspace\IndexPageList[0849]{402,403,404}
\subsubitem \leavevmode\hypertarget{idx:0850}{双重构造引理},\enspace\IndexPageList[0850]{398,401,403}
\subsubitem \leavevmode\hypertarget{idx:0851}{例题},\enspace\IndexPageList[0851]{382,387}
\subsubitem \leavevmode\hypertarget{idx:0852}{$\mathbb F_2$ 的镜像定理},\enspace\IndexPageList[0852]{417}
\subsubitem \leavevmode\hypertarget{idx:0853}{$\mathbb P^1$ 的镜像定理},\enspace\IndexPageList[0853]{347,380}
\subsubitem \leavevmode\hypertarget{idx:0854}{$\mathbb P^n$ 的镜像定理},\enspace\IndexPageList[0854]{348,380,389,394,395}
\subsubitem \leavevmode\hypertarget{idx:0855}{$\mathbb P(1,1,2,2,2)$ 中超曲面的镜像定理},\enspace\IndexPageList[0855]{406,407,408,409,411}
\subsubitem \leavevmode\hypertarget{idx:0856}{$\mathbb P^n$ 中 nef 完全交的镜像定理},\enspace\IndexPageList[0856]{379,381,382,396,402,403}
\subsubitem \leavevmode\hypertarget{idx:0857}{亏格 $1$ 版本},\enspace\IndexPageList[0857]{404}
\subsubitem \leavevmode\hypertarget{idx:0858}{如何求变量变换},\enspace\IndexPageList[0858]{382}
\subsubitem \leavevmode\hypertarget{idx:0859}{五次三维簇的镜像定理},\enspace\IndexPageList[0859]{49,384,385,386}
\subsubitem \leavevmode\hypertarget{idx:0860}{递推引理},\enspace\IndexPageList[0860]{392,394,395,397,398,401,403}
\subsubitem \leavevmode\hypertarget{idx:0861}{正则性引理},\enspace\IndexPageList[0861]{392,396,398,401}
\subsubitem \leavevmode\hypertarget{idx:0862}{环面镜像定理},\enspace\IndexPageList[0862]{404,405,410,412,414,428}，\textit{另见} \hyperlink{idx:1141}{量子上同调，小量子上同调，光滑环面簇的小量子上同调}
\subsubitem \leavevmode\hypertarget{idx:0863}{唯一性引理},\enspace\IndexPageList[0863]{401,403}
\subitem \leavevmode\hypertarget{idx:0864}{Lian--Liu--Yau 镜像定理},\enspace\IndexPageList[0864]{51,286,355,357,372,377,403,421}
\subsubitem \leavevmode\hypertarget{idx:0865}{对多重覆盖的应用},\enspace\IndexPageList[0865]{378}
\subsubitem \leavevmode\hypertarget{idx:0866}{环面镜像定理},\enspace\IndexPageList[0866]{421}
\subitem \leavevmode\hypertarget{idx:0867}{五次三维簇的镜像定理},\enspace\IndexPageList[0867]{47,48,49,50,51,52,286,355,357,375,377,385}
\indexspace
\item \leavevmode\hypertarget{idx:0613}{镜像定理中的齐次坐标变换}，\textit{参见} \hyperlink{idx:1438}{变量的权}
\indexspace
\item \leavevmode\hypertarget{idx:0826}{镜像对},\enspace\IndexPageList[0826]{28,29,32,76,79,425,444}，\textit{另见} \hyperlink{idx:0798}{镜像构造}
\subitem \leavevmode\hypertarget{idx:0827}{Hodge 理论镜像对}，\textit{参见} \hyperlink{idx:0754}{数学镜像对，Hodge 理论镜像对}
\subitem \leavevmode\hypertarget{idx:0828}{数学镜像对}，\textit{参见} \hyperlink{idx:0748}{数学镜像对}
\subitem \leavevmode\hypertarget{idx:0829}{物理定义},\enspace\IndexPageList[0829]{27}
\indexspace
\item \leavevmode\hypertarget{idx:0830}{镜像对称},\enspace\IndexPageList[0830]{24,27,28,29,32,33,39,58,69,76,105,130,134,138,139,151,152,154,165,167,170,175,273,279,284,289,351,446,449}，\textit{另见} \hyperlink{idx:0786}{镜像猜想}；\hyperlink{idx:0844}{镜像定理}
\subitem \leavevmode\hypertarget{idx:0831}{范畴镜像对称}，\textit{参见} \hyperlink{idx:0840}{镜像对称，同调镜像对称}
\subitem \leavevmode\hypertarget{idx:0832}{经典镜像对称},\enspace\IndexPageList[0832]{32,36,47,50,356,386,408}
\subitem \leavevmode\hypertarget{idx:0833}{最早的证据},\enspace\IndexPageList[0833]{83}
\subitem \leavevmode\hypertarget{idx:0834}{枚举方面},\enspace\IndexPageList[0834]{237}
\subitem \leavevmode\hypertarget{idx:0835}{K3 曲面的镜像对称},\enspace\IndexPageList[0835]{424}
\subitem \leavevmode\hypertarget{idx:0836}{规范线性 $\sigma$ 模型的镜像对称},\enspace\IndexPageList[0836]{453}
\subitem \leavevmode\hypertarget{idx:0837}{五次三维簇的镜像对称}，\textit{参见} \hyperlink{idx:1197}{五次三维簇，镜像对称}
\subitem \leavevmode\hypertarget{idx:0838}{刚性 Calabi--Yau 簇的镜像对称},\enspace\IndexPageList[0838]{88}
\subitem \leavevmode\hypertarget{idx:0839}{Hodge 理论镜像对称},\enspace\IndexPageList[0839]{48,177,385}，\textit{另见} \hyperlink{idx:0754}{数学镜像对，Hodge 理论镜像对}
\subitem \leavevmode\hypertarget{idx:0840}{同调镜像对称},\enspace\IndexPageList[0840]{34,426}
\subitem \leavevmode\hypertarget{idx:0841}{数学镜像对称},\enspace\IndexPageList[0841]{34,135,240,265,280,297,420,425}，\textit{另见} \hyperlink{idx:0748}{数学镜像对}；\hyperlink{idx:0786}{镜像猜想}
\subitem \leavevmode\hypertarget{idx:0842}{与 Givental 镜像定理的关系},\enspace\IndexPageList[0842]{406}
\subitem \leavevmode\hypertarget{idx:0843}{环面镜像对称},\enspace\IndexPageList[0843]{79}
\indexspace
\item \leavevmode\hypertarget{idx:0798}{镜像构造},\enspace\IndexPageList[0798]{284}
\subitem \leavevmode\hypertarget{idx:0799}{Batyrev 构造},\enspace\IndexPageList[0799]{76,79,83,93,95,143,146,147,148,158,162,166,167,169,171,178,284,290,291,294,295,404,406}
\subsubitem \leavevmode\hypertarget{idx:0800}{五次三维簇},\enspace\IndexPageList[0800]{84}
\subitem \leavevmode\hypertarget{idx:0801}{Batyrev--Borisov 构造},\enspace\IndexPageList[0801]{85,86,123,170,428}，\textit{另见} \hyperlink{idx:0935}{nef 分拆}
\subitem \leavevmode\hypertarget{idx:0802}{旗流形中的 Calabi--Yau 完全交的镜像构造},\enspace\IndexPageList[0802]{428}
\subitem \leavevmode\hypertarget{idx:0803}{Greene--Plesser 构造},\enspace\IndexPageList[0803]{76,79,112}
\subitem \leavevmode\hypertarget{idx:0804}{Voisin--Borcea 构造},\enspace\IndexPageList[0804]{88,91,92,95,424,425}
\subsubitem \leavevmode\hypertarget{idx:0805}{复模与 Kähler 模},\enspace\IndexPageList[0805]{91,92}
\subsubitem \leavevmode\hypertarget{idx:0806}{与 Batyrev 构造的关系},\enspace\IndexPageList[0806]{93}
\subsubitem \leavevmode\hypertarget{idx:0807}{Yukawa 耦合},\enspace\IndexPageList[0807]{92}
\indexspace
\item \leavevmode\hypertarget{idx:0809}{镜像映射},\enspace\IndexPageList[0809]{28,32,79,83,105,132,151,166,170,171,172,173,174,175,176,177,178,179,180,181,182,183,184,185,186,187,188,189,280,408}
\subitem \leavevmode\hypertarget{idx:0810}{定义},\enspace\IndexPageList[0810]{173,180}
\subsubitem \leavevmode\hypertarget{idx:0811}{需要作出的选择},\enspace\IndexPageList[0811]{173,175,177}
\subsubitem \leavevmode\hypertarget{idx:0812}{定义所需的数据},\enspace\IndexPageList[0812]{173}
\subsubitem \leavevmode\hypertarget{idx:0813}{定义的不确定性},\enspace\IndexPageList[0813]{173}
\subsubitem \leavevmode\hypertarget{idx:0814}{物理依据},\enspace\IndexPageList[0814]{181}
\subsubitem \leavevmode\hypertarget{idx:0815}{整性猜想的作用},\enspace\IndexPageList[0815]{173,174}
\subitem \leavevmode\hypertarget{idx:0816}{导数},\enspace\IndexPageList[0816]{175,176,177,181,289,290}，\textit{另见} \hyperlink{idx:0811}{镜像映射，定义，需要作出的选择}
\subitem \leavevmode\hypertarget{idx:0817}{五次三维簇的镜像映射}，\textit{参见} \hyperlink{idx:1195}{五次三维簇，镜像映射}
\subitem \leavevmode\hypertarget{idx:0818}{镜像坐标},\enspace\IndexPageList[0818]{133,135,175,282,288}
\subsubitem \leavevmode\hypertarget{idx:0819}{定义},\enspace\IndexPageList[0819]{172}
\subitem \leavevmode\hypertarget{idx:0820}{环面镜像映射},\enspace\IndexPageList[0820]{171,174,178,179,180,181,182,290,291,293,295}
\subsubitem \leavevmode\hypertarget{idx:0821}{定义},\enspace\IndexPageList[0821]{178,179}
\subsubitem \leavevmode\hypertarget{idx:0822}{五次三维簇的环面镜像映射},\enspace\IndexPageList[0822]{179,180}
\subsubitem \leavevmode\hypertarget{idx:0823}{不为恒等映射},\enspace\IndexPageList[0823]{171}
\subsubitem \leavevmode\hypertarget{idx:0824}{消除非唯一性},\enspace\IndexPageList[0824]{179}
\subitem \leavevmode\hypertarget{idx:0825}{通过超几何函数构造},\enspace\IndexPageList[0825]{182,183,184,185,186,187,188,189}
\indexspace
\item \leavevmode\hypertarget{idx:0808}{镜像坐标}，\textit{参见} \hyperlink{idx:0809}{镜像映射}
\indexspace
\item \leavevmode\hypertarget{idx:0723}{局部化},\enspace\IndexPageList[0723]{302,303,389,405}，\textit{另见} \hyperlink{idx:0734}{局部化定理}
\subitem \leavevmode\hypertarget{idx:0724}{叠的局部化},\enspace\IndexPageList[0724]{312}
\indexspace
\item \leavevmode\hypertarget{idx:0734}{局部化定理}
\subitem \leavevmode\hypertarget{idx:0735}{在枚举几何中的应用},\enspace\IndexPageList[0735]{305}
\subitem \leavevmode\hypertarget{idx:0736}{应用于 $\mathbb P^n$},\enspace\IndexPageList[0736]{301,302}
\subitem \leavevmode\hypertarget{idx:0737}{计算瞬子数 $n_4$},\enspace\IndexPageList[0737]{311}
\subitem \leavevmode\hypertarget{idx:0738}{计算五次三维簇上的直线},\enspace\IndexPageList[0738]{305}
\subitem \leavevmode\hypertarget{idx:0739}{等变上同调的局部化定理},\enspace\IndexPageList[0739]{300,342,361,362,370,371,372}
\subitem \leavevmode\hypertarget{idx:0740}{Chern 类多项式的局部化定理},\enspace\IndexPageList[0740]{304,314,315,316}
\subitem \leavevmode\hypertarget{idx:0741}{光滑叠的局部化定理},\enspace\IndexPageList[0741]{303,304,312,391,396,398,400}
\subitem \leavevmode\hypertarget{idx:0742}{等变积分的局部化定理},\enspace\IndexPageList[0742]{302,358,362,363,364,365,368,371,399}
\subitem \leavevmode\hypertarget{idx:0743}{证明多重覆盖公式},\enspace\IndexPageList[0743]{312}
\indexspace
\item \leavevmode\hypertarget{idx:0722}{局部算子},\enspace\IndexPageList[0722]{444,446}
\indexspace
\item \leavevmode\hypertarget{idx:0898}{矩映射},\enspace\IndexPageList[0898]{64,67,68,72,166,452}
\subitem \leavevmode\hypertarget{idx:0899}{例题},\enspace\IndexPageList[0899]{73}
\indexspace
\item \leavevmode\hypertarget{idx:0760}{矩阵积分},\enspace\IndexPageList[0760]{333}
\indexspace
\item \leavevmode\hypertarget{idx:0761}{矩阵模型},\enspace\IndexPageList[0761]{333}
\indexspace
\item \leavevmode\hypertarget{idx:0913}{可动锥},\enspace\IndexPageList[0913]{156,165}
\subitem \leavevmode\hypertarget{idx:0914}{反射可动锥},\enspace\IndexPageList[0914]{157,165}
\indexspace
\item \leavevmode\hypertarget{idx:0956}{可观测量},\enspace\IndexPageList[0956]{24,25,59,437}
\indexspace
\item \leavevmode\hypertarget{idx:1035}{裤形曲面},\enspace\IndexPageList[1035]{456,457}
\indexspace
\item \leavevmode\hypertarget{idx:0376}{亏格 $g$ 的引力 Gromov--Witten 势},\enspace\IndexPageList[0376]{326}
\subitem \leavevmode\hypertarget{idx:0377}{定义},\enspace\IndexPageList[0377]{327}
\indexspace
\item \leavevmode\hypertarget{idx:0375}{亏格 $g$ 耦合},\enspace\IndexPageList[0375]{326,335}
\indexspace
\item \leavevmode\hypertarget{idx:0172}{联络 $\nabla^c$},\enspace\IndexPageList[0172]{270}，\textit{另见} \hyperlink{idx:0003}{A 模型，联络}
\subitem \leavevmode\hypertarget{idx:0173}{平坦截面},\enspace\IndexPageList[0173]{270,271,272,277,339,341}
\indexspace
\item \leavevmode\hypertarget{idx:0174}{联络矩阵},\enspace\IndexPageList[0174]{48,50,99,102,129,131,269}
\indexspace
\item \leavevmode\hypertarget{idx:1109}{量子 $\mathcal D$ 模},\enspace\IndexPageList[1109]{428}
\indexspace
\item \leavevmode\hypertarget{idx:1155}{量子场论},\enspace\IndexPageList[1155]{436,437,438,439,440,441,442,445,453}
\subitem \leavevmode\hypertarget{idx:1156}{Hamilton 表述},\enspace\IndexPageList[1156]{439}
\subitem \leavevmode\hypertarget{idx:1157}{Lagrange 表述},\enspace\IndexPageList[1157]{439}
\subitem \leavevmode\hypertarget{idx:1158}{二维量子场论},\enspace\IndexPageList[1158]{24}
\indexspace
\item \leavevmode\hypertarget{idx:1159}{量子超平面截面原理},\enspace\IndexPageList[1159]{49,387,388,389,428}
\subitem \leavevmode\hypertarget{idx:1160}{$\mathbb P^n$ 的量子超平面截面原理},\enspace\IndexPageList[1160]{388}
\subitem \leavevmode\hypertarget{idx:1161}{凸簇的量子超平面截面原理},\enspace\IndexPageList[1161]{388}
\subitem \leavevmode\hypertarget{idx:1162}{与 Givental 镜像定理的关系},\enspace\IndexPageList[1162]{389}
\indexspace
\item \leavevmode\hypertarget{idx:1110}{量子反常},\enspace\IndexPageList[1110]{446,447,448}
\indexspace
\item \leavevmode\hypertarget{idx:1106}{量子化},\enspace\IndexPageList[1106]{33,436,438,441,442,447,448,453}
\subitem \leavevmode\hypertarget{idx:1107}{正则量子化},\enspace\IndexPageList[1107]{436,437}
\subitem \leavevmode\hypertarget{idx:1108}{径向量子化},\enspace\IndexPageList[1108]{447}
\indexspace
\item \leavevmode\hypertarget{idx:1111}{量子上同调},\enspace\IndexPageList[1111]{30,48,105,124,223,445,453,458}
\subitem \leavevmode\hypertarget{idx:1112}{大量子上同调},\enspace\IndexPageList[1112]{252,253,327,335,337}，\textit{另见} \hyperlink{idx:0552}{Gromov--Witten 势}
\subsubitem \leavevmode\hypertarget{idx:1113}{结合性},\enspace\IndexPageList[1113]{255,257,260}，\textit{另见} \hyperlink{idx:0568}{Gromov--Witten 势，WDVV 方程}
\subsubitem \leavevmode\hypertarget{idx:1114}{收敛性},\enspace\IndexPageList[1114]{253}
\subsubitem \leavevmode\hypertarget{idx:1115}{定义},\enspace\IndexPageList[1115]{254}
\subsubitem \leavevmode\hypertarget{idx:1116}{次数条件},\enspace\IndexPageList[1116]{258}
\subsubitem \leavevmode\hypertarget{idx:1117}{K\"unneth 公式},\enspace\IndexPageList[1117]{257}
\subsubitem \leavevmode\hypertarget{idx:1118}{$\mathbb P^2$ 的大量子上同调},\enspace\IndexPageList[1118]{260}
\subsubitem \leavevmode\hypertarget{idx:1119}{Calabi--Yau 三维簇的大量子上同调},\enspace\IndexPageList[1119]{262}
\subsubitem \leavevmode\hypertarget{idx:1120}{五次三维簇的大量子上同调},\enspace\IndexPageList[1120]{261}
\subsubitem \leavevmode\hypertarget{idx:1121}{与小量子乘积的关系},\enspace\IndexPageList[1121]{265,266}
\subsubitem \leavevmode\hypertarget{idx:1122}{限制},\enspace\IndexPageList[1122]{265,266,267,337,343,350}
\subsubitem \leavevmode\hypertarget{idx:1123}{环结构},\enspace\IndexPageList[1123]{256}，\textit{另见} \hyperlink{idx:0337}{Frobenius 代数}
\subsubitem \leavevmode\hypertarget{idx:1124}{用于定义 Givental 联络},\enspace\IndexPageList[1124]{334}
\subitem \leavevmode\hypertarget{idx:1125}{系数},\enspace\IndexPageList[1125]{251,252,253}
\subitem \leavevmode\hypertarget{idx:1126}{引力量子上同调}，\textit{参见} \hyperlink{idx:0481}{引力量子乘积}
\subitem \leavevmode\hypertarget{idx:1127}{修正量子上同调},\enspace\IndexPageList[1127]{386}
\subsubitem \leavevmode\hypertarget{idx:1128}{其中的关系},\enspace\IndexPageList[1128]{386}
\subitem \leavevmode\hypertarget{idx:1129}{小量子上同调},\enspace\IndexPageList[1129]{240,241,242,243,244,245,246,247,248,249,250,251,253,337,338,343,349,382,386,387}
\subsubitem \leavevmode\hypertarget{idx:1130}{结合性},\enspace\IndexPageList[1130]{242,251}
\subsubitem \leavevmode\hypertarget{idx:1131}{收敛性},\enspace\IndexPageList[1131]{241,250,251,267,268,423}
\subsubitem \leavevmode\hypertarget{idx:1132}{定义},\enspace\IndexPageList[1132]{241}
\subsubitem \leavevmode\hypertarget{idx:1133}{次数条件},\enspace\IndexPageList[1133]{244,245}
\subsubitem \leavevmode\hypertarget{idx:1134}{翻越的影响},\enspace\IndexPageList[1134]{249,250}
\subsubitem \leavevmode\hypertarget{idx:1135}{$(1+1)$ 维拓扑量子场论的小量子上同调},\enspace\IndexPageList[1135]{457}
\subsubitem \leavevmode\hypertarget{idx:1136}{$\mathbb F_2$ 的小量子上同调},\enspace\IndexPageList[1136]{416,417,418}
\subsubitem \leavevmode\hypertarget{idx:1137}{$\mathbb P^n$ 的小量子上同调},\enspace\IndexPageList[1137]{245,246,260,348}
\subsubitem \leavevmode\hypertarget{idx:1138}{Calabi--Yau 流形的小量子上同调},\enspace\IndexPageList[1138]{241,245,275}
\subsubitem \leavevmode\hypertarget{idx:1139}{Calabi--Yau 三维簇的小量子上同调},\enspace\IndexPageList[1139]{248,249,250,262,350,352}
\subsubitem \leavevmode\hypertarget{idx:1140}{Grassmann 簇的小量子上同调},\enspace\IndexPageList[1140]{250,428}
\subsubitem \leavevmode\hypertarget{idx:1141}{光滑环面簇的小量子上同调},\enspace\IndexPageList[1141]{246,413,415,416}，\textit{另见} \hyperlink{idx:0069}{Batyrev 量子环}
\subsubitem \leavevmode\hypertarget{idx:1142}{五次三维簇的小量子上同调},\enspace\IndexPageList[1142]{248,249,261}
\subsubitem \leavevmode\hypertarget{idx:1143}{其他例题},\enspace\IndexPageList[1143]{250,251}
\subsubitem \leavevmode\hypertarget{idx:1144}{性质},\enspace\IndexPageList[1144]{244}
\subsubitem \leavevmode\hypertarget{idx:1145}{与大量子乘积的关系},\enspace\IndexPageList[1145]{265,266}
\subsubitem \leavevmode\hypertarget{idx:1146}{其中的关系},\enspace\IndexPageList[1146]{344,345,348,387,410,413,415}
\subsubitem \leavevmode\hypertarget{idx:1147}{环结构},\enspace\IndexPageList[1147]{242}，\textit{另见} \hyperlink{idx:0337}{Frobenius 代数}
\subsubitem \leavevmode\hypertarget{idx:1148}{用于定义 A 模型联络},\enspace\IndexPageList[1148]{266,267}
\indexspace
\item \leavevmode\hypertarget{idx:1150}{量子微分方程},\enspace\IndexPageList[1150]{344,347,350,352,379,409,413,428}，\textit{另见} \hyperlink{idx:1151}{量子微分算子}
\indexspace
\item \leavevmode\hypertarget{idx:1151}{量子微分算子},\enspace\IndexPageList[1151]{344,345,352,415,418}
\subitem \leavevmode\hypertarget{idx:1152}{由量子上同调中的关系导出},\enspace\IndexPageList[1152]{345,347}
\subitem \leavevmode\hypertarget{idx:1153}{例题},\enspace\IndexPageList[1153]{347,348,350}
\subitem \leavevmode\hypertarget{idx:1154}{给出量子上同调中的关系},\enspace\IndexPageList[1154]{344}，\textit{另见} \hyperlink{idx:1146}{量子上同调，小量子上同调，其中的关系}
\indexspace
\item \leavevmode\hypertarget{idx:1149}{量子修正},\enspace\IndexPageList[1149]{35,41,171}
\indexspace
\item \leavevmode\hypertarget{idx:0205}{临界丛}，\textit{参见} \hyperlink{idx:0147}{凸凹丛，临界丛}
\indexspace
\item \leavevmode\hypertarget{idx:1231}{留数矩阵 $\operatorname{Res}(\nabla)$},\enspace\IndexPageList[1231]{100,102,269}
\indexspace
\item \leavevmode\hypertarget{idx:1230}{留数映射},\enspace\IndexPageList[1230]{106,109}
\indexspace
\item \leavevmode\hypertarget{idx:0218}{六结点五次曲线的二重覆盖},\enspace\IndexPageList[0218]{316,317,318,319,320}
\subitem \leavevmode\hypertarget{idx:0219}{对 Gromov--Witten 不变量及虚拟基本类的贡献},\enspace\IndexPageList[0219]{320}
\subitem \leavevmode\hypertarget{idx:0220}{经正规化分解的覆盖},\enspace\IndexPageList[0220]{318,319,320}
\subitem \leavevmode\hypertarget{idx:0221}{不经正规化分解的覆盖},\enspace\IndexPageList[0221]{318,319}
\subitem \leavevmode\hypertarget{idx:0222}{模},\enspace\IndexPageList[0222]{318}
\subitem \leavevmode\hypertarget{idx:0223}{切空间},\enspace\IndexPageList[0223]{319}
\indexspace
\item \leavevmode\hypertarget{idx:1037}{路径积分},\enspace\IndexPageList[1037]{439,440,441,445}，\textit{另见} \hyperlink{idx:0315}{Feynman 积分}
\indexspace
\item \leavevmode\hypertarget{idx:0241}{枚举几何},\enspace\IndexPageList[0241]{39,238}，\textit{另见} \hyperlink{idx:0534}{Gromov--Witten 不变量，枚举意义}
\indexspace
\item \leavevmode\hypertarget{idx:0938}{幂零轨道},\enspace\IndexPageList[0938]{277}
\indexspace
\item \leavevmode\hypertarget{idx:0939}{幂零轨道定理},\enspace\IndexPageList[0939]{99}
\indexspace
\item \leavevmode\hypertarget{idx:0870}{模},\enspace\IndexPageList[0870]{24,28}
\subitem \leavevmode\hypertarget{idx:0871}{复模},\enspace\IndexPageList[0871]{26,27,28,32,79,104,105,136,137,138,139,140,141,142,143,144,145,146,147,148,149,150,151,154,167,171,172,287,446}
\subsubitem \leavevmode\hypertarget{idx:0872}{一维复模},\enspace\IndexPageList[0872]{101,102,105,107,110,125,128,130,132,280}
\subsubitem \leavevmode\hypertarget{idx:0873}{紧化},\enspace\IndexPageList[0873]{137,138,140,142,143,144}
\subsubitem \leavevmode\hypertarget{idx:0874}{坐标},\enspace\IndexPageList[0874]{137,146,147,178,182,290}
\subsubitem \leavevmode\hypertarget{idx:0875}{维数},\enspace\IndexPageList[0875]{141,142}
\subsubitem \leavevmode\hypertarget{idx:0876}{例题},\enspace\IndexPageList[0876]{146,147,148,149,150,162,170}
\subsubitem \leavevmode\hypertarget{idx:0877}{最大幂幺边界点}，\textit{参见} \hyperlink{idx:0766}{最大幂幺，最大幂幺边界点}
\subsubitem \leavevmode\hypertarget{idx:0878}{Calabi--Yau 三维簇的复模},\enspace\IndexPageList[0878]{101,102,105,125,128,130,132,144,280,281,283}
\subsubitem \leavevmode\hypertarget{idx:0879}{Calabi--Yau 环面超曲面的复模},\enspace\IndexPageList[0879]{115,116,139,140,141,142,143,144,145,146,147,148,149,150}
\subsubitem \leavevmode\hypertarget{idx:0880}{多项式复模},\enspace\IndexPageList[0880]{84,139,140,141,142,147,148,149,161,170,406,450}
\subsubitem \leavevmode\hypertarget{idx:0881}{简化复模},\enspace\IndexPageList[0881]{118,140,141,142,143,144,145,146,147,148,149,150,158,161,165,178,179,290,294}
\subitem \leavevmode\hypertarget{idx:0882}{Kähler 模},\enspace\IndexPageList[0882]{26,27,28,34,35,79,105,139,150,151,152,153,154,155,156,157,158,159,160,161,162,163,164,165,166,167,168,169,170,171,172,175,241,268,287,446}
\subsubitem \leavevmode\hypertarget{idx:0883}{一维 Kähler 模},\enspace\IndexPageList[0883]{172}
\subsubitem \leavevmode\hypertarget{idx:0884}{边界},\enspace\IndexPageList[0884]{166}
\subsubitem \leavevmode\hypertarget{idx:0885}{紧化},\enspace\IndexPageList[0885]{155,158,159,165,166,167,175}
\subsubitem \leavevmode\hypertarget{idx:0886}{坐标},\enspace\IndexPageList[0886]{151,152,290}
\subsubitem \leavevmode\hypertarget{idx:0887}{扩大 Kähler 模},\enspace\IndexPageList[0887]{155,156,157,159,165}
\subsubitem \leavevmode\hypertarget{idx:0888}{例题},\enspace\IndexPageList[0888]{162,163,164,165,169}
\subsubitem \leavevmode\hypertarget{idx:0889}{全局 Kähler 模},\enspace\IndexPageList[0889]{165}
\subsubitem \leavevmode\hypertarget{idx:0890}{大半径极限点}，\textit{参见} \hyperlink{idx:0696}{大半径极限点}
\subsubitem \leavevmode\hypertarget{idx:0891}{Calabi--Yau 三维簇的 Kähler 模},\enspace\IndexPageList[0891]{155,156,157,165,280,281}
\subsubitem \leavevmode\hypertarget{idx:0892}{Calabi--Yau 环面超曲面的 Kähler 模},\enspace\IndexPageList[0892]{157,158,159,160,161,162,163,164,165,166}
\subsubitem \leavevmode\hypertarget{idx:0893}{部分紧化},\enspace\IndexPageList[0893]{151,152,153,154,155,157,161,164,165,173}
\subsubitem \leavevmode\hypertarget{idx:0894}{环面 Kähler 模},\enspace\IndexPageList[0894]{83,143,157,158,161,164,178,179,290,294}
\subitem \leavevmode\hypertarget{idx:0895}{超共形场论模},\enspace\IndexPageList[0895]{26,27,28,151,167,168,287,424}
\subsubitem \leavevmode\hypertarget{idx:0896}{例题},\enspace\IndexPageList[0896]{170}
\subitem \leavevmode\hypertarget{idx:0897}{稳定映射的模}，\textit{参见} \hyperlink{idx:1276}{稳定映射，模空间}
\indexspace
\item \leavevmode\hypertarget{idx:1249}{$\sigma$ 模型},\enspace\IndexPageList[1249]{26,27,29,30,32,33,34,38,191,325,357,442,443,444,445,446,451,452,453}
\subitem \leavevmode\hypertarget{idx:1250}{A 模型},\enspace\IndexPageList[1250]{444}
\subitem \leavevmode\hypertarget{idx:1251}{作用量}，\textit{参见} \hyperlink{idx:0034}{作用量，$\sigma$ 模型的作用量}
\subitem \leavevmode\hypertarget{idx:1252}{B 模型},\enspace\IndexPageList[1252]{444}
\subitem \leavevmode\hypertarget{idx:1253}{费米场},\enspace\IndexPageList[1253]{442,444,448}
\subitem \leavevmode\hypertarget{idx:1254}{超对称 $\sigma$ 模型},\enspace\IndexPageList[1254]{444,446}
\subitem \leavevmode\hypertarget{idx:1255}{拓扑 $\sigma$ 模型},\enspace\IndexPageList[1255]{457}
\subitem \leavevmode\hypertarget{idx:1256}{扭曲 $\sigma$ 模型},\enspace\IndexPageList[1256]{444,445}
\indexspace
\item \leavevmode\hypertarget{idx:0640}{内蕴法层},\enspace\IndexPageList[0640]{201,202}，\textit{另见} \hyperlink{idx:0948}{法锥，内蕴法锥}
\indexspace
\item \leavevmode\hypertarget{idx:0639}{内蕴法锥}，\textit{参见} \hyperlink{idx:0948}{法锥，内蕴法锥}
\indexspace
\item \leavevmode\hypertarget{idx:0240}{能动张量},\enspace\IndexPageList[0240]{447}
\indexspace
\item \leavevmode\hypertarget{idx:0239}{能动量守恒},\enspace\IndexPageList[0239]{447}
\indexspace
\item \leavevmode\hypertarget{idx:1163}{拟光滑},\enspace\IndexPageList[1163]{83,109,139,431}，\textit{另见} \hyperlink{idx:0990}{轨形}
\indexspace
\item \leavevmode\hypertarget{idx:0214}{判别式轨迹},\enspace\IndexPageList[0214]{127,145}，\textit{另见} \hyperlink{idx:0103}{Calabi--Yau，Calabi--Yau 环面超曲面，非退化情形}
\subitem \leavevmode\hypertarget{idx:0215}{例题},\enspace\IndexPageList[0215]{148,170,181,182}
\subitem \leavevmode\hypertarget{idx:0216}{主分量},\enspace\IndexPageList[0216]{145,148}
\indexspace
\item \leavevmode\hypertarget{idx:1036}{配分函数},\enspace\IndexPageList[1036]{333}
\indexspace
\item \leavevmode\hypertarget{idx:1257}{奇点},\enspace\IndexPageList[1257]{430,431,432,433}
\subitem \leavevmode\hypertarget{idx:1258}{$\mathbb Q$ 析因奇点},\enspace\IndexPageList[1258]{34,78,430}
\subitem \leavevmode\hypertarget{idx:1259}{典范奇点},\enspace\IndexPageList[1259]{34,35,430,431}
\subitem \leavevmode\hypertarget{idx:1260}{Cohen--Macaulay 奇点},\enspace\IndexPageList[1260]{430}
\subitem \leavevmode\hypertarget{idx:1261}{Gorenstein 奇点},\enspace\IndexPageList[1261]{34,35,78,139,430}
\subitem \leavevmode\hypertarget{idx:1262}{终端奇点},\enspace\IndexPageList[1262]{34,77,78,139,430,431,433}
\subitem \leavevmode\hypertarget{idx:1263}{环面型奇点},\enspace\IndexPageList[1263]{433}
\indexspace
\item \leavevmode\hypertarget{idx:0324}{旗流形},\enspace\IndexPageList[0324]{242,251,277}
\subitem \leavevmode\hypertarget{idx:0325}{旗流形中的完全交},\enspace\IndexPageList[0325]{387,389}
\subsubitem \leavevmode\hypertarget{idx:0326}{Calabi--Yau 完全交},\enspace\IndexPageList[0326]{428}
\indexspace
\item \leavevmode\hypertarget{idx:0577}{强 Lefschetz 定理},\enspace\IndexPageList[0577]{273,432}
\indexspace
\item \leavevmode\hypertarget{idx:1326}{切函子},\enspace\IndexPageList[1326]{199}
\indexspace
\item \leavevmode\hypertarget{idx:1327}{切权}，\textit{参见} \hyperlink{idx:1437}{表示的权，切权}
\indexspace
\item \leavevmode\hypertarget{idx:1328}{切—障碍复形},\enspace\IndexPageList[1328]{199,201,202,203,308}
\subitem \leavevmode\hypertarget{idx:1329}{用于虚拟基本类},\enspace\IndexPageList[1329]{201}
\subitem \leavevmode\hypertarget{idx:1330}{完美切—障碍复形},\enspace\IndexPageList[1330]{197,204}
\subsubitem \leavevmode\hypertarget{idx:1331}{定义},\enspace\IndexPageList[1331]{199}
\subsubitem \leavevmode\hypertarget{idx:1332}{例题},\enspace\IndexPageList[1332]{200}
\indexspace
\item \leavevmode\hypertarget{idx:0206}{曲率张量},\enspace\IndexPageList[0206]{443}
\indexspace
\item \leavevmode\hypertarget{idx:0612}{全纯辛},\enspace\IndexPageList[0612]{425}
\indexspace
\item \leavevmode\hypertarget{idx:0869}{缺失的镜像},\enspace\IndexPageList[0869]{95}
\indexspace
\item \leavevmode\hypertarget{idx:1333}{三点函数}，\textit{参见} \hyperlink{idx:0195}{关联函数，三点关联函数}
\indexspace
\item \leavevmode\hypertarget{idx:0297}{扇},\enspace\IndexPageList[0297]{54}，\textit{另见} \hyperlink{idx:1094}{射影细分}
\subitem \leavevmode\hypertarget{idx:0298}{完备扇},\enspace\IndexPageList[0298]{55}
\subitem \leavevmode\hypertarget{idx:0299}{扇的翻转}，\textit{参见} \hyperlink{idx:0328}{翻转}
\subitem \leavevmode\hypertarget{idx:0300}{Gröbner 扇}，\textit{参见} \hyperlink{idx:0453}{Gr\"obner 扇}
\subitem \leavevmode\hypertarget{idx:0301}{次级扇}，\textit{参见} \hyperlink{idx:1239}{次级扇}
\subitem \leavevmode\hypertarget{idx:0302}{单纯扇},\enspace\IndexPageList[0302]{55}
\subitem \leavevmode\hypertarget{idx:0303}{光滑扇},\enspace\IndexPageList[0303]{55}
\subitem \leavevmode\hypertarget{idx:0304}{支集},\enspace\IndexPageList[0304]{54}
\subitem \leavevmode\hypertarget{idx:0305}{支撑于线性回路上的扇}，\textit{参见} \hyperlink{idx:0720}{线性回路，由扇支撑}
\subitem \leavevmode\hypertarget{idx:0306}{对称性},\enspace\IndexPageList[0306]{71}
\subitem \leavevmode\hypertarget{idx:0307}{扇的环面簇},\enspace\IndexPageList[0307]{55,74}
\indexspace
\item \leavevmode\hypertarget{idx:1317}{扇的支集}，\textit{参见} \hyperlink{idx:0304}{扇，支集}
\indexspace
\item \leavevmode\hypertarget{idx:0126}{上同调场论},\enspace\IndexPageList[0126]{264}，\textit{另见} \hyperlink{idx:0235}{Dubrovin 形式体系}
\subitem \leavevmode\hypertarget{idx:0127}{关联函数},\enspace\IndexPageList[0127]{264}
\subitem \leavevmode\hypertarget{idx:0128}{高亏格情形},\enspace\IndexPageList[0128]{264}
\subitem \leavevmode\hypertarget{idx:0129}{Künneth 公式},\enspace\IndexPageList[0129]{264}
\indexspace
\item \leavevmode\hypertarget{idx:1351}{上同调的环面部分},\enspace\IndexPageList[1351]{157,166,179,181,290,384,405}
\subitem \leavevmode\hypertarget{idx:1352}{定义},\enspace\IndexPageList[1352]{80,294}
\indexspace
\item \leavevmode\hypertarget{idx:0967}{尚未解决的猜想},\enspace\IndexPageList[0967]{422,423}，\textit{另见} \hyperlink{idx:0786}{镜像猜想}
\subitem \leavevmode\hypertarget{idx:0968}{Calabi--Yau 三维簇}
\subsubitem \leavevmode\hypertarget{idx:0969}{用 $I_{\mathcal V}$ 表示 B 模型势函数},\enspace\IndexPageList[0969]{412,422}
\subsubitem \leavevmode\hypertarget{idx:0970}{用 $I_{\mathcal V}$ 表示 Yukawa 耦合},\enspace\IndexPageList[0970]{412,422}
\subitem \leavevmode\hypertarget{idx:0971}{复几何}
\subsubitem \leavevmode\hypertarget{idx:0972}{整性猜想},\enspace\IndexPageList[0972]{105,422}
\subsubitem \leavevmode\hypertarget{idx:0973}{简化模空间中的最大幂幺边界点},\enspace\IndexPageList[0973]{143,422}
\subitem \leavevmode\hypertarget{idx:0974}{量子上同调的收敛性},\enspace\IndexPageList[0974]{241,423}
\subitem \leavevmode\hypertarget{idx:0975}{枚举几何}
\subsubitem \leavevmode\hypertarget{idx:0976}{Clemens 猜想},\enspace\IndexPageList[0976]{225,423}
\subsubitem \leavevmode\hypertarget{idx:0977}{Gromov--Witten 不变量的枚举意义},\enspace\IndexPageList[0977]{219,423}
\subsubitem \leavevmode\hypertarget{idx:0978}{瞬子数的辛定义},\enspace\IndexPageList[0978]{231,423}
\subitem \leavevmode\hypertarget{idx:0979}{Kähler 几何}
\subsubitem \leavevmode\hypertarget{idx:0980}{锥猜想},\enspace\IndexPageList[0980]{153,422}
\subsubitem \leavevmode\hypertarget{idx:0981}{Calabi--Yau 环面超曲面的 Kähler 锥的环面部分},\enspace\IndexPageList[0981]{161,422}
\subitem \leavevmode\hypertarget{idx:0982}{Picard--Fuchs 方程}
\subsubitem \leavevmode\hypertarget{idx:0983}{$I_{\mathcal V}$ 是解},\enspace\IndexPageList[0983]{124,422}
\subsubitem \leavevmode\hypertarget{idx:0984}{$\mathcal A$ 方程组中缺失的方程},\enspace\IndexPageList[0984]{117,422}
\subitem \leavevmode\hypertarget{idx:0985}{$J_X$ 与 $J_{\mathcal V}$ 的关系},\enspace\IndexPageList[0985]{409,423}
\subitem \leavevmode\hypertarget{idx:0986}{虚拟基本类}
\subsubitem \leavevmode\hypertarget{idx:0987}{定义的等价性},\enspace\IndexPageList[0987]{199,422}
\subsubitem \leavevmode\hypertarget{idx:0988}{完全交的虚拟基本类},\enspace\IndexPageList[0988]{409,423}
\indexspace
\item \leavevmode\hypertarget{idx:1094}{射影细分},\enspace\IndexPageList[1094]{77,78,160}
\subitem \leavevmode\hypertarget{idx:1095}{定义},\enspace\IndexPageList[1095]{76}
\subitem \leavevmode\hypertarget{idx:1096}{极大射影细分},\enspace\IndexPageList[1096]{77,78,79,83,85,86,139,158,290,405}
\subsubitem \leavevmode\hypertarget{idx:1097}{定义},\enspace\IndexPageList[1097]{76}
\subitem \leavevmode\hypertarget{idx:1098}{简化射影细分},\enspace\IndexPageList[1098]{159,162,163,164,166,169}
\subsubitem \leavevmode\hypertarget{idx:1099}{定义},\enspace\IndexPageList[1099]{158}
\indexspace
\item \leavevmode\hypertarget{idx:1204}{升指标},\enspace\IndexPageList[1204]{241,262,440}
\indexspace
\item \leavevmode\hypertarget{idx:1440}{世界面},\enspace\IndexPageList[1440]{24,25,33,442,445,446,447,448,449,450}
\subitem \leavevmode\hypertarget{idx:1441}{非微扰修正},\enspace\IndexPageList[1441]{30}
\indexspace
\item \leavevmode\hypertarget{idx:1442}{世界体},\enspace\IndexPageList[1442]{425}
\indexspace
\item \leavevmode\hypertarget{idx:1083}{势函数}，\textit{参见} \hyperlink{idx:0064}{B 模型，势函数}；\hyperlink{idx:0552}{Gromov--Witten 势}
\indexspace
\item \leavevmode\hypertarget{idx:0121}{手征超场},\enspace\IndexPageList[0121]{449,450,451}
\indexspace
\item \leavevmode\hypertarget{idx:0120}{手征环},\enspace\IndexPageList[0120]{59,240,450}
\indexspace
\item \leavevmode\hypertarget{idx:0175}{守恒流},\enspace\IndexPageList[0175]{442}
\indexspace
\item \leavevmode\hypertarget{idx:0748}{数学镜像对},\enspace\IndexPageList[0748]{27,34,48,79,265,410,420}
\subitem \leavevmode\hypertarget{idx:0749}{全局数学镜像对},\enspace\IndexPageList[0749]{281}
\subitem \leavevmode\hypertarget{idx:0750}{高维数学镜像对},\enspace\IndexPageList[0750]{287,288,289,290}
\subsubitem \leavevmode\hypertarget{idx:0751}{定义},\enspace\IndexPageList[0751]{287}
\subsubitem \leavevmode\hypertarget{idx:0752}{四维簇的数学镜像对},\enspace\IndexPageList[0752]{288,289}
\subsubitem \leavevmode\hypertarget{idx:0753}{局限},\enspace\IndexPageList[0753]{288}
\subitem \leavevmode\hypertarget{idx:0754}{Hodge 理论镜像对},\enspace\IndexPageList[0754]{281,287}
\subitem \leavevmode\hypertarget{idx:0755}{Calabi--Yau 三维簇的数学镜像对},\enspace\IndexPageList[0755]{280,281,282,283,284,285}
\subsubitem \leavevmode\hypertarget{idx:0756}{定义},\enspace\IndexPageList[0756]{281}
\subsubitem \leavevmode\hypertarget{idx:0757}{用关联函数表述},\enspace\IndexPageList[0757]{282,285}
\subsubitem \leavevmode\hypertarget{idx:0758}{用势函数表述},\enspace\IndexPageList[0758]{283,285}
\subitem \leavevmode\hypertarget{idx:0759}{五次镜像与五次三维簇},\enspace\IndexPageList[0759]{285,286,356}
\indexspace
\item \leavevmode\hypertarget{idx:0077}{双变 Chow 群},\enspace\IndexPageList[0077]{203,205}
\indexspace
\item \leavevmode\hypertarget{idx:0076}{双有理锥猜想}，\textit{参见} \hyperlink{idx:0158}{锥猜想，双有理锥猜想}
\indexspace
\item \leavevmode\hypertarget{idx:0075}{双棱锥构造},\enspace\IndexPageList[0075]{93,95}
\indexspace
\item \leavevmode\hypertarget{idx:0622}{瞬子}
\subitem \leavevmode\hypertarget{idx:0623}{瞬子修正},\enspace\IndexPageList[0623]{32,446}
\subitem \leavevmode\hypertarget{idx:0624}{全纯瞬子},\enspace\IndexPageList[0624]{31,453}
\subitem \leavevmode\hypertarget{idx:0625}{朴素瞬子},\enspace\IndexPageList[0625]{453}
\subitem \leavevmode\hypertarget{idx:0626}{瞬子求和},\enspace\IndexPageList[0626]{445,446}
\subitem \leavevmode\hypertarget{idx:0627}{世界面瞬子},\enspace\IndexPageList[0627]{445}
\indexspace
\item \leavevmode\hypertarget{idx:0628}{瞬子数},\enspace\IndexPageList[0628]{31,226,227,228,229,230,231,232,233,249,331}
\subitem \leavevmode\hypertarget{idx:0629}{定义},\enspace\IndexPageList[0629]{226}
\subitem \leavevmode\hypertarget{idx:0630}{整除性质},\enspace\IndexPageList[0630]{47}
\subitem \leavevmode\hypertarget{idx:0631}{椭圆情形的类似概念},\enspace\IndexPageList[0631]{236,237}
\subitem \leavevmode\hypertarget{idx:0632}{枚举意义},\enspace\IndexPageList[0632]{228,229,230,231}
\subitem \leavevmode\hypertarget{idx:0633}{例题},\enspace\IndexPageList[0633]{228,229,230,231}，\textit{另见} \hyperlink{idx:1191}{五次三维簇，瞬子数}
\subitem \leavevmode\hypertarget{idx:0634}{五次三维簇的瞬子数}，\textit{参见} \hyperlink{idx:1191}{五次三维簇，瞬子数}
\subitem \leavevmode\hypertarget{idx:0635}{辛定义},\enspace\IndexPageList[0635]{231,232,233,423}
\indexspace
\item \leavevmode\hypertarget{idx:0989}{算子乘积展开},\enspace\IndexPageList[0989]{447,448}
\indexspace
\item \leavevmode\hypertarget{idx:1345}{拓扑 $\sigma$ 模型},\enspace\IndexPageList[1345]{190}
\subitem \leavevmode\hypertarget{idx:1346}{与引力耦合},\enspace\IndexPageList[1346]{190,210}
\indexspace
\item \leavevmode\hypertarget{idx:1344}{拓扑递推关系},\enspace\IndexPageList[1344]{334,335}
\indexspace
\item \leavevmode\hypertarget{idx:1336}{拓扑量子场论},\enspace\IndexPageList[1336]{190,191,240,439,445,453,454,455,456,457,458}
\subitem \leavevmode\hypertarget{idx:1337}{$(0+1)$ 维拓扑量子场论},\enspace\IndexPageList[1337]{454,455}
\subitem \leavevmode\hypertarget{idx:1338}{$(1+1)$ 维拓扑量子场论},\enspace\IndexPageList[1338]{455}
\subsubitem \leavevmode\hypertarget{idx:1339}{给出 Frobenius 代数},\enspace\IndexPageList[1339]{455,457}
\subitem \leavevmode\hypertarget{idx:1340}{$(2+1)$ 维拓扑量子场论},\enspace\IndexPageList[1340]{458}
\subitem \leavevmode\hypertarget{idx:1341}{$(d+1)$ 维拓扑量子场论},\enspace\IndexPageList[1341]{454}
\subitem \leavevmode\hypertarget{idx:1342}{定义},\enspace\IndexPageList[1342]{454}
\subitem \leavevmode\hypertarget{idx:1343}{与 A 模型的关系},\enspace\IndexPageList[1343]{457}
\indexspace
\item \leavevmode\hypertarget{idx:0176}{凸}
\subitem \leavevmode\hypertarget{idx:0177}{凸丛},\enspace\IndexPageList[0177]{372}，\textit{另见} \hyperlink{idx:0146}{凸凹丛}
\subitem \leavevmode\hypertarget{idx:0178}{凸线丛},\enspace\IndexPageList[0178]{387}
\subitem \leavevmode\hypertarget{idx:0179}{凸簇},\enspace\IndexPageList[0179]{321,381,388,389,404,405}
\subsubitem \leavevmode\hypertarget{idx:0180}{定义},\enspace\IndexPageList[0180]{203}
\indexspace
\item \leavevmode\hypertarget{idx:1248}{$\Sigma$ 凸}，\textit{参见} \hyperlink{idx:0202}{$\operatorname{cpl}(\Sigma)$，$\Sigma$ 凸性}
\indexspace
\item \leavevmode\hypertarget{idx:0146}{凸凹丛},\enspace\IndexPageList[0146]{360,372,373,375}
\subitem \leavevmode\hypertarget{idx:0147}{临界丛},\enspace\IndexPageList[0147]{373,375}
\subsubitem \leavevmode\hypertarget{idx:0148}{定义},\enspace\IndexPageList[0148]{374}
\indexspace
\item \leavevmode\hypertarget{idx:0181}{凸函数},\enspace\IndexPageList[0181]{62}，\textit{另见} \hyperlink{idx:0199}{$\operatorname{cpl}(\Sigma)$}
\subitem \leavevmode\hypertarget{idx:0182}{严格凸函数},\enspace\IndexPageList[0182]{62,65,66,68}
\indexspace
\item \leavevmode\hypertarget{idx:0327}{图 $\Gamma$ 的旗},\enspace\IndexPageList[0327]{308}，\textit{另见} \hyperlink{idx:0726}{$\overline M_{0,n}(\mathbb P^n,d)$ 中的局部化，不动点轨迹}
\indexspace
\item \leavevmode\hypertarget{idx:0131}{完全交},\enspace\IndexPageList[0131]{387}
\subitem \leavevmode\hypertarget{idx:0132}{$\mathbb P^n$ 中的完全交},\enspace\IndexPageList[0132]{238,387}
\subitem \leavevmode\hypertarget{idx:0133}{凸簇中的完全交},\enspace\IndexPageList[0133]{389}
\subsubitem \leavevmode\hypertarget{idx:0134}{nef 完全交},\enspace\IndexPageList[0134]{388}
\subitem \leavevmode\hypertarget{idx:0135}{旗流形中的完全交}，\textit{参见} \hyperlink{idx:0324}{旗流形}
\subitem \leavevmode\hypertarget{idx:0136}{Grassmann 簇中的完全交}，\textit{参见} \hyperlink{idx:0454}{Grassmann 簇}
\subitem \leavevmode\hypertarget{idx:0137}{环面完全交}，\textit{参见} \hyperlink{idx:1347}{环面完全交}
\subitem \leavevmode\hypertarget{idx:0138}{虚拟基本类}，\textit{参见} \hyperlink{idx:1419}{虚拟基本类，完全交的虚拟基本类}
\indexspace
\item \leavevmode\hypertarget{idx:1040}{微扰态},\enspace\IndexPageList[1040]{425}
\indexspace
\item \leavevmode\hypertarget{idx:1100}{伪流形},\enspace\IndexPageList[1100]{212}
\indexspace
\item \leavevmode\hypertarget{idx:1269}{稳定 $C^\ell$ 映射},\enspace\IndexPageList[1269]{211,214}
\subitem \leavevmode\hypertarget{idx:1270}{等价性},\enspace\IndexPageList[1270]{211}
\indexspace
\item \leavevmode\hypertarget{idx:1271}{稳定映射}
\subitem \leavevmode\hypertarget{idx:1272}{自同构},\enspace\IndexPageList[1272]{304,306,307,313,314,315,363,364,391,393,400}
\subitem \leavevmode\hypertarget{idx:1273}{定义},\enspace\IndexPageList[1273]{192}
\subitem \leavevmode\hypertarget{idx:1274}{稳定映射族},\enspace\IndexPageList[1274]{192}
\subitem \leavevmode\hypertarget{idx:1275}{图构造},\enspace\IndexPageList[1275]{342,357}
\subitem \leavevmode\hypertarget{idx:1276}{模空间},\enspace\IndexPageList[1276]{39,50,196}，\textit{另见} \hyperlink{idx:0725}{$\overline M_{0,n}(\mathbb P^n,d)$ 中的局部化}；\hyperlink{idx:0777}{$M_d$}
\subsubitem \leavevmode\hypertarget{idx:1277}{粗模空间},\enspace\IndexPageList[1277]{192,193,201,204,324}
\subsubitem \leavevmode\hypertarget{idx:1278}{维数},\enspace\IndexPageList[1278]{193}
\subsubitem \leavevmode\hypertarget{idx:1279}{预期维数},\enspace\IndexPageList[1279]{194,198,203,210,215,234,329,331}
\subsubitem \leavevmode\hypertarget{idx:1280}{模叠},\enspace\IndexPageList[1280]{193,198,203,204,324,325,380,381,382,387}
\subsubitem \leavevmode\hypertarget{idx:1281}{切空间},\enspace\IndexPageList[1281]{198,308,313}
\subitem \leavevmode\hypertarget{idx:1282}{万有稳定映射},\enspace\IndexPageList[1282]{202,324,325,328}
\indexspace
\item \leavevmode\hypertarget{idx:1164}{五次镜像},\enspace\IndexPageList[1164]{40,145,356}
\subitem \leavevmode\hypertarget{idx:1165}{$\mathcal A$ 方程组},\enspace\IndexPageList[1165]{116,117}
\subsubitem \leavevmode\hypertarget{idx:1166}{解},\enspace\IndexPageList[1166]{124}
\subitem \leavevmode\hypertarget{idx:1167}{B 模型关联函数},\enspace\IndexPageList[1167]{40,44,45,46,50,52}，\textit{另见} \hyperlink{idx:0189}{关联函数，B 模型关联函数}
\subitem \leavevmode\hypertarget{idx:1168}{B 模型势函数},\enspace\IndexPageList[1168]{286,356}
\subitem \leavevmode\hypertarget{idx:1169}{复模},\enspace\IndexPageList[1169]{41,285}
\subitem \leavevmode\hypertarget{idx:1170}{构造},\enspace\IndexPageList[1170]{40,84}
\subitem \leavevmode\hypertarget{idx:1171}{Frobenius 方法},\enspace\IndexPageList[1171]{185}
\subitem \leavevmode\hypertarget{idx:1172}{Gauss--Manin 联络},\enspace\IndexPageList[1172]{48}
\subitem \leavevmode\hypertarget{idx:1173}{齐次坐标环},\enspace\IndexPageList[1173]{85}
\subitem \leavevmode\hypertarget{idx:1174}{单值性},\enspace\IndexPageList[1174]{134,135,285}
\subitem \leavevmode\hypertarget{idx:1175}{Picard--Fuchs 方程},\enspace\IndexPageList[1175]{43,44,49,104,110,111,118,186,356}
\subsubitem \leavevmode\hypertarget{idx:1176}{解},\enspace\IndexPageList[1176]{46,124,133,180,183,185,187,356,376,385}
\subitem \leavevmode\hypertarget{idx:1177}{自反多胞形},\enspace\IndexPageList[1177]{84}
\subitem \leavevmode\hypertarget{idx:1178}{满足整性猜想},\enspace\IndexPageList[1178]{105,133,173,180}
\subitem \leavevmode\hypertarget{idx:1179}{Yukawa 耦合},\enspace\IndexPageList[1179]{44,180}
\subsubitem \leavevmode\hypertarget{idx:1180}{归一化 Yukawa 耦合},\enspace\IndexPageList[1180]{44,45,46,49,50,133,134,285,386}，\textit{另见} \hyperlink{idx:1167}{五次镜像，B 模型关联函数}
\indexspace
\item \leavevmode\hypertarget{idx:1181}{五次三维簇},\enspace\IndexPageList[1181]{38,39,40,41,42,43,44,45,46,47,48,49,50,51,52,375,452}
\subitem \leavevmode\hypertarget{idx:1182}{其上的 $2875$ 条直线},\enspace\IndexPageList[1182]{46,47,180,196,304,305}
\subitem \leavevmode\hypertarget{idx:1183}{A 模型关联函数},\enspace\IndexPageList[1183]{38,39,45,46,50,51,226,249,285,356,386}，\textit{另见} \hyperlink{idx:0184}{关联函数，A 模型关联函数}
\subitem \leavevmode\hypertarget{idx:1184}{大量子上同调},\enspace\IndexPageList[1184]{261}
\subitem \leavevmode\hypertarget{idx:1185}{Fermat 五次三维簇},\enspace\IndexPageList[1185]{166,197}
\subitem \leavevmode\hypertarget{idx:1186}{Gromov--Witten 不变量},\enspace\IndexPageList[1186]{196,197,224,225,226,316,355}
\subsubitem \leavevmode\hypertarget{idx:1187}{等变不变量},\enspace\IndexPageList[1187]{322}
\subsubitem \leavevmode\hypertarget{idx:1188}{公式},\enspace\IndexPageList[1188]{207,224,304,355,373}
\subsubitem \leavevmode\hypertarget{idx:1189}{通过局部化计算},\enspace\IndexPageList[1189]{311}
\subitem \leavevmode\hypertarget{idx:1190}{Gromov--Witten 势},\enspace\IndexPageList[1190]{51,260,261,356,373,385,411}
\subitem \leavevmode\hypertarget{idx:1191}{瞬子数},\enspace\IndexPageList[1191]{39,46,47,48,228,230,285,316,356,386}，\textit{另见} \hyperlink{idx:0628}{瞬子数}
\subsubitem \leavevmode\hypertarget{idx:1192}{$n_4$ 的计算},\enspace\IndexPageList[1192]{311}
\subsubitem \leavevmode\hypertarget{idx:1193}{$n_{10}$ 的细微之处},\enspace\IndexPageList[1193]{39,229,320}
\subitem \leavevmode\hypertarget{idx:1194}{Kähler 模},\enspace\IndexPageList[1194]{38,41,45,180}
\subitem \leavevmode\hypertarget{idx:1195}{镜像映射},\enspace\IndexPageList[1195]{41,42,43,45,46,49,51,133,179,180,185,286,356,376,385}
\subitem \leavevmode\hypertarget{idx:1196}{镜像}，\textit{参见} \hyperlink{idx:1164}{五次镜像}
\subitem \leavevmode\hypertarget{idx:1197}{镜像对称},\enspace\IndexPageList[1197]{38,41,43,45,46,285,286,311,356}，\textit{另见} \hyperlink{idx:0867}{镜像定理，五次三维簇的镜像定理}
\subitem \leavevmode\hypertarget{idx:1198}{镜像定理}，\textit{参见} \hyperlink{idx:0867}{镜像定理，五次三维簇的镜像定理}
\subitem \leavevmode\hypertarget{idx:1199}{其上的有理曲线},\enspace\IndexPageList[1199]{24,38,39,47,225,229,311,316}，\textit{另见} \hyperlink{idx:0124}{Clemens 猜想}
\subsubitem \leavevmode\hypertarget{idx:1200}{其上的六结点五次曲线},\enspace\IndexPageList[1200]{39,229,316,318,320}
\subsubitem \leavevmode\hypertarget{idx:1201}{法丛},\enspace\IndexPageList[1201]{317}
\subitem \leavevmode\hypertarget{idx:1202}{小量子上同调},\enspace\IndexPageList[1202]{248,249}
\subitem \leavevmode\hypertarget{idx:1203}{虚拟基本类},\enspace\IndexPageList[1203]{196,203,204,207}
\indexspace
\item \leavevmode\hypertarget{idx:1048}{物理学家的 Euler 数},\enspace\IndexPageList[1048]{83,425}，\textit{另见} \hyperlink{idx:0592}{Hodge 数，弦论 Hodge 数}
\indexspace
\item \leavevmode\hypertarget{idx:1265}{稀疏结式},\enspace\IndexPageList[1265]{145}
\indexspace
\item \leavevmode\hypertarget{idx:1287}{弦理论},\enspace\IndexPageList[1287]{24,25,30,442}，\textit{另见} \hyperlink{idx:1304}{超弦理论}
\subitem \leavevmode\hypertarget{idx:1288}{极限},\enspace\IndexPageList[1288]{429}
\subitem \leavevmode\hypertarget{idx:1289}{周期},\enspace\IndexPageList[1289]{429}
\indexspace
\item \leavevmode\hypertarget{idx:0721}{线性 $\sigma$ 模型}，\textit{参见} \hyperlink{idx:0354}{规范线性 $\sigma$ 模型}
\indexspace
\item \leavevmode\hypertarget{idx:0700}{线性关系格},\enspace\IndexPageList[0700]{62,64,67,159,161,163,178,182,247,413}
\subitem \leavevmode\hypertarget{idx:0701}{例题},\enspace\IndexPageList[0701]{73,116,118,147,149,183,187}
\subitem \leavevmode\hypertarget{idx:0702}{与 $\mathcal A$ 方程组的关系},\enspace\IndexPageList[0702]{114,120,122}
\indexspace
\item \leavevmode\hypertarget{idx:0717}{线性回路},\enspace\IndexPageList[0717]{159,161}
\subitem \leavevmode\hypertarget{idx:0718}{例题},\enspace\IndexPageList[0718]{159,163,164}
\subitem \leavevmode\hypertarget{idx:0719}{收缩 GKZ 分解},\enspace\IndexPageList[0719]{162}
\subitem \leavevmode\hypertarget{idx:0720}{由扇支撑},\enspace\IndexPageList[0720]{160,162}
\indexspace
\item \leavevmode\hypertarget{idx:1042}{相},\enspace\IndexPageList[1042]{69,166}，\textit{另见} \hyperlink{idx:0354}{规范线性 $\sigma$ 模型}
\indexspace
\item \leavevmode\hypertarget{idx:1041}{相变},\enspace\IndexPageList[1041]{452}
\indexspace
\item \leavevmode\hypertarget{idx:1397}{向量超场},\enspace\IndexPageList[1397]{451}
\indexspace
\item \leavevmode\hypertarget{idx:1387}{消失圈},\enspace\IndexPageList[1387]{42}
\indexspace
\item \leavevmode\hypertarget{idx:1264}{小子群},\enspace\IndexPageList[1264]{431}
\indexspace
\item \leavevmode\hypertarget{idx:1318}{辛}
\subitem \leavevmode\hypertarget{idx:1319}{辛作用},\enspace\IndexPageList[1319]{63}
\subitem \leavevmode\hypertarget{idx:1320}{辛流形},\enspace\IndexPageList[1320]{63,207,208,232,251}
\subsubitem \leavevmode\hypertarget{idx:1321}{半正辛流形}，\textit{参见} \hyperlink{idx:1245}{半正}
\subitem \leavevmode\hypertarget{idx:1322}{辛轨形},\enspace\IndexPageList[1322]{65,433}
\subitem \leavevmode\hypertarget{idx:1323}{辛约化},\enspace\IndexPageList[1323]{64,65,67,68,166,433}
\indexspace
\item \leavevmode\hypertarget{idx:1324}{辛商}，\textit{参见} \hyperlink{idx:1323}{辛，辛约化}
\indexspace
\item \leavevmode\hypertarget{idx:0208}{形变}
\subitem \leavevmode\hypertarget{idx:0209}{无穷小形变},\enspace\IndexPageList[0209]{80,176,198}
\subitem \leavevmode\hypertarget{idx:0210}{极小 Calabi--Yau 轨形的形变},\enspace\IndexPageList[0210]{433}
\subitem \leavevmode\hypertarget{idx:0211}{轨形的形变},\enspace\IndexPageList[0211]{433}
\subitem \leavevmode\hypertarget{idx:0212}{非多项式形变},\enspace\IndexPageList[0212]{81,170}
\subitem \leavevmode\hypertarget{idx:0213}{多项式形变},\enspace\IndexPageList[0213]{80,84}
\indexspace
\item \leavevmode\hypertarget{idx:1427}{虚法丛},\enspace\IndexPageList[1427]{305,323}
\indexspace
\item \leavevmode\hypertarget{idx:1428}{虚法锥}，\textit{参见} \hyperlink{idx:0951}{法锥，虚法锥}
\indexspace
\item \leavevmode\hypertarget{idx:1406}{虚拟基本类},\enspace\IndexPageList[1406]{39,195,197,198,199,200,201,202,203,204,205,206,214,230,311,317,322,325,328,329,331,336,381,383,404}
\subitem \leavevmode\hypertarget{idx:1407}{与遗忘映射的相容性},\enspace\IndexPageList[1407]{205,328}
\subitem \leavevmode\hypertarget{idx:1408}{用 Euler 类计算},\enspace\IndexPageList[1408]{205,231,234,312,331,332}
\subitem \leavevmode\hypertarget{idx:1409}{定义}，\textit{另见} \hyperlink{idx:0959}{障碍理论}；\hyperlink{idx:1328}{切—障碍复形}
\subsubitem \leavevmode\hypertarget{idx:1410}{Behrend--Fantechi 方法},\enspace\IndexPageList[1410]{201,202}
\subsubitem \leavevmode\hypertarget{idx:1411}{Behrend--Fantechi 定义},\enspace\IndexPageList[1411]{202}
\subsubitem \leavevmode\hypertarget{idx:1412}{定义的等价性},\enspace\IndexPageList[1412]{199,218,422}
\subsubitem \leavevmode\hypertarget{idx:1413}{Li--Tian 方法},\enspace\IndexPageList[1413]{199,200,201,204}
\subsubitem \leavevmode\hypertarget{idx:1414}{Li--Tian 定义},\enspace\IndexPageList[1414]{200,201}
\subitem \leavevmode\hypertarget{idx:1415}{次数等于预期维数},\enspace\IndexPageList[1415]{201,215,245,327}
\subitem \leavevmode\hypertarget{idx:1416}{等变虚拟基本类},\enspace\IndexPageList[1416]{322,323,405}
\subsubitem \leavevmode\hypertarget{idx:1417}{局部化公式},\enspace\IndexPageList[1417]{308,323,405}
\subitem \leavevmode\hypertarget{idx:1418}{例题},\enspace\IndexPageList[1418]{200,202}
\subitem \leavevmode\hypertarget{idx:1419}{完全交的虚拟基本类},\enspace\IndexPageList[1419]{409,423}
\subsubitem \leavevmode\hypertarget{idx:1420}{$\mathbb P^n$ 中完全交的虚拟基本类},\enspace\IndexPageList[1420]{383}
\subitem \leavevmode\hypertarget{idx:1421}{Calabi--Yau 三维簇的虚拟基本类},\enspace\IndexPageList[1421]{311,312}
\subitem \leavevmode\hypertarget{idx:1422}{Fermat 五次三维簇的虚拟基本类},\enspace\IndexPageList[1422]{197}
\subitem \leavevmode\hypertarget{idx:1423}{五次三维簇的虚拟基本类},\enspace\IndexPageList[1423]{196,203,204,383}
\subitem \leavevmode\hypertarget{idx:1424}{性质},\enspace\IndexPageList[1424]{203,204,205,206,218,327}
\subitem \leavevmode\hypertarget{idx:1425}{相对虚拟基本类},\enspace\IndexPageList[1425]{203}
\indexspace
\item \leavevmode\hypertarget{idx:1426}{虚邻域},\enspace\IndexPageList[1426]{211}
\indexspace
\item \leavevmode\hypertarget{idx:1429}{虚切空间},\enspace\IndexPageList[1429]{323}
\indexspace
\item \leavevmode\hypertarget{idx:1268}{旋量},\enspace\IndexPageList[1268]{449}
\indexspace
\item \leavevmode\hypertarget{idx:1266}{旋量丛},\enspace\IndexPageList[1266]{442}
\indexspace
\item \leavevmode\hypertarget{idx:1267}{旋量几何},\enspace\IndexPageList[1267]{449}
\indexspace
\item \leavevmode\hypertarget{idx:0049}{湮灭算子},\enspace\IndexPageList[0049]{438,441}
\indexspace
\item \leavevmode\hypertarget{idx:0479}{引力 Gromov--Witten 势}，\textit{参见} \hyperlink{idx:0376}{亏格 $g$ 的引力 Gromov--Witten 势}
\indexspace
\item \leavevmode\hypertarget{idx:0459}{引力关联数},\enspace\IndexPageList[0459]{48,219,234,324,325,326,327,328,329,330,331,332,333,334,343,405,406}，\textit{另见} \hyperlink{idx:0376}{亏格 $g$ 的引力 Gromov--Witten 势}；\hyperlink{idx:0375}{亏格 $g$ 耦合}
\subitem \leavevmode\hypertarget{idx:0460}{公理}，\textit{另见} \hyperlink{idx:0457}{引力类}
\subsubitem \leavevmode\hypertarget{idx:0461}{次数公理},\enspace\IndexPageList[0461]{327,333,334,340,341,348,383}
\subsubitem \leavevmode\hypertarget{idx:0462}{伸缩公理},\enspace\IndexPageList[0462]{329,330,331,333,341}
\subsubitem \leavevmode\hypertarget{idx:0463}{除子公理},\enspace\IndexPageList[0463]{328,329,330,331,333,337,341,380}
\subsubitem \leavevmode\hypertarget{idx:0464}{基本类公理},\enspace\IndexPageList[0464]{328,329,330,331,341,347,348}
\subsubitem \leavevmode\hypertarget{idx:0465}{分裂公理},\enspace\IndexPageList[0465]{328,329,330,336}
\subitem \leavevmode\hypertarget{idx:0466}{系数},\enspace\IndexPageList[0466]{326}
\subitem \leavevmode\hypertarget{idx:0467}{定义},\enspace\IndexPageList[0467]{325}
\subitem \leavevmode\hypertarget{idx:0468}{用 Gromov--Witten 不变量表示亏格 $0$ 引力关联数},\enspace\IndexPageList[0468]{329}
\subitem \leavevmode\hypertarget{idx:0469}{$\mathbb P^1$ 的亏格 $0$ 引力关联数},\enspace\IndexPageList[0469]{329,380}
\subsubitem \leavevmode\hypertarget{idx:0470}{递推关系},\enspace\IndexPageList[0470]{330,331}
\subitem \leavevmode\hypertarget{idx:0471}{Calabi--Yau 三维簇的亏格 $0$ 引力关联数},\enspace\IndexPageList[0471]{331}
\subitem \leavevmode\hypertarget{idx:0472}{亏格 $1$ 引力关联数 $\langle\tau_1\rangle_{1,0}$},\enspace\IndexPageList[0472]{331,332}
\subsubitem \leavevmode\hypertarget{idx:0473}{与 Euler 示性数的关系},\enspace\IndexPageList[0473]{333}
\subitem \leavevmode\hypertarget{idx:0474}{由拓扑递推关系得到的递推式},\enspace\IndexPageList[0474]{334}，\textit{另见} \hyperlink{idx:1344}{拓扑递推关系}
\subitem \leavevmode\hypertarget{idx:0475}{与 Gromov--Witten 不变量的关系},\enspace\IndexPageList[0475]{326}
\subitem \leavevmode\hypertarget{idx:0476}{符号记法},\enspace\IndexPageList[0476]{338,379}
\indexspace
\item \leavevmode\hypertarget{idx:0477}{引力后裔},\enspace\IndexPageList[0477]{325,334}
\subitem \leavevmode\hypertarget{idx:0478}{与引力关联数的关系},\enspace\IndexPageList[0478]{326}
\indexspace
\item \leavevmode\hypertarget{idx:0457}{引力类},\enspace\IndexPageList[0457]{329,332}，\textit{另见} \hyperlink{idx:0459}{引力关联数}
\subitem \leavevmode\hypertarget{idx:0458}{定义},\enspace\IndexPageList[0458]{328}
\indexspace
\item \leavevmode\hypertarget{idx:0481}{引力量子乘积},\enspace\IndexPageList[0481]{257,327}
\subitem \leavevmode\hypertarget{idx:0482}{系数},\enspace\IndexPageList[0482]{327}
\indexspace
\item \leavevmode\hypertarget{idx:0480}{引力势},\enspace\IndexPageList[0480]{333}
\indexspace
\item \leavevmode\hypertarget{idx:1205}{有理多面锥},\enspace\IndexPageList[1205]{54}
\subitem \leavevmode\hypertarget{idx:1206}{强凸有理多面锥},\enspace\IndexPageList[1206]{54}
\indexspace
\item \leavevmode\hypertarget{idx:1234}{右行}
\subitem \leavevmode\hypertarget{idx:1235}{右行费米子},\enspace\IndexPageList[1235]{25,444}
\subitem \leavevmode\hypertarget{idx:1236}{右行扇区},\enspace\IndexPageList[1236]{448}
\indexspace
\item \leavevmode\hypertarget{idx:0197}{余切复形},\enspace\IndexPageList[0197]{201}
\indexspace
\item \leavevmode\hypertarget{idx:0198}{余切线 $\mathcal L_i$},\enspace\IndexPageList[0198]{325,392}
\indexspace
\item \leavevmode\hypertarget{idx:0281}{预期维数}
\subitem \leavevmode\hypertarget{idx:0282}{Kuranishi 理论中的预期维数},\enspace\IndexPageList[0282]{197}
\subitem \leavevmode\hypertarget{idx:0283}{稳定映射模空间的预期维数}，\textit{参见} \hyperlink{idx:1276}{稳定映射，模空间}
\subitem \leavevmode\hypertarget{idx:0284}{零点轨迹的预期维数},\enspace\IndexPageList[0284]{196}
\indexspace
\item \leavevmode\hypertarget{idx:1084}{预稳定映射},\enspace\IndexPageList[1084]{202}
\indexspace
\item \leavevmode\hypertarget{idx:0451}{粘合引理},\enspace\IndexPageList[0451]{361,365,403}
\indexspace
\item \leavevmode\hypertarget{idx:0217}{占优猜想},\enspace\IndexPageList[0217]{140,178}
\indexspace
\item \leavevmode\hypertarget{idx:0957}{障碍丛},\enspace\IndexPageList[0957]{231,312,313,315}，\textit{另见} \hyperlink{idx:1435}{表示的权，障碍权}
\indexspace
\item \leavevmode\hypertarget{idx:0958}{障碍空间},\enspace\IndexPageList[0958]{198}
\indexspace
\item \leavevmode\hypertarget{idx:0959}{障碍理论},\enspace\IndexPageList[0959]{201,323}
\subitem \leavevmode\hypertarget{idx:0960}{完美障碍理论},\enspace\IndexPageList[0960]{197,203,322}
\subsubitem \leavevmode\hypertarget{idx:0961}{定义},\enspace\IndexPageList[0961]{202}
\subsubitem \leavevmode\hypertarget{idx:0962}{等变完美障碍理论},\enspace\IndexPageList[0962]{322}
\subsubitem \leavevmode\hypertarget{idx:0963}{例题},\enspace\IndexPageList[0963]{202}
\subsubitem \leavevmode\hypertarget{idx:0964}{用于虚拟基本类},\enspace\IndexPageList[0964]{203}
\indexspace
\item \leavevmode\hypertarget{idx:1385}{真空态},\enspace\IndexPageList[1385]{438,440,448,450}
\indexspace
\item \leavevmode\hypertarget{idx:0636}{整性猜想},\enspace\IndexPageList[0636]{128,129,130,132,133,135,174,178,180,273,281,282,285,289,291,293,408,409,422}
\subitem \leavevmode\hypertarget{idx:0637}{表述},\enspace\IndexPageList[0637]{105,173}
\indexspace
\item \leavevmode\hypertarget{idx:0952}{正规化正合列},\enspace\IndexPageList[0952]{313,315}
\indexspace
\item \leavevmode\hypertarget{idx:1225}{正则奇点},\enspace\IndexPageList[1225]{43,98,100,269,270}
\indexspace
\item \leavevmode\hypertarget{idx:1226}{正则三角剖分},\enspace\IndexPageList[1226]{66,67}
\indexspace
\item \leavevmode\hypertarget{idx:1227}{正则性引理}，\textit{参见} \hyperlink{idx:0861}{镜像定理，Givental 镜像定理，正则性引理}
\indexspace
\item \leavevmode\hypertarget{idx:1314}{支撑函数},\enspace\IndexPageList[1314]{55}
\subitem \leavevmode\hypertarget{idx:1315}{上凸支撑函数},\enspace\IndexPageList[1315]{56}
\subsubitem \leavevmode\hypertarget{idx:1316}{严格上凸支撑函数},\enspace\IndexPageList[1316]{56,58}
\indexspace
\item \leavevmode\hypertarget{idx:0621}{指标方程},\enspace\IndexPageList[0621]{100,102,103,186}
\indexspace
\item \leavevmode\hypertarget{idx:0119}{中心荷},\enspace\IndexPageList[0119]{333,448}
\indexspace
\item \leavevmode\hypertarget{idx:1208}{重构定理},\enspace\IndexPageList[1208]{219,220,223}
\indexspace
\item \leavevmode\hypertarget{idx:1229}{重整化},\enspace\IndexPageList[1229]{25,30,31,446}，\textit{另见} \hyperlink{idx:0944}{非重整化定理}
\indexspace
\item \leavevmode\hypertarget{idx:1038}{周期},\enspace\IndexPageList[1038]{42,43,50,98,115,117,118,182,356}，\textit{另见} \hyperlink{idx:1059}{Picard--Fuchs 方程，解}
\indexspace
\item \leavevmode\hypertarget{idx:1039}{周期域},\enspace\IndexPageList[1039]{277}
\indexspace
\item \leavevmode\hypertarget{idx:1093}{主 $\mathcal A$ 行列式},\enspace\IndexPageList[1093]{145}
\indexspace
\item \leavevmode\hypertarget{idx:1286}{驻相法},\enspace\IndexPageList[1286]{440,445}
\indexspace
\item \leavevmode\hypertarget{idx:0149}{锥}
\subitem \leavevmode\hypertarget{idx:0150}{Calabi--Yau 锥}，\textit{参见} \hyperlink{idx:0110}{Calabi--Yau 锥}
\subitem \leavevmode\hypertarget{idx:0151}{对偶锥}，\textit{参见} \hyperlink{idx:0224}{对偶锥}
\subitem \leavevmode\hypertarget{idx:0152}{锥的面}，\textit{参见} \hyperlink{idx:0296}{锥的面}
\subitem \leavevmode\hypertarget{idx:0153}{Gorenstein 锥},\enspace\IndexPageList[0153]{87}
\subsubitem \leavevmode\hypertarget{idx:0154}{自反 Gorenstein 锥}，\textit{参见} \hyperlink{idx:1213}{自反 Gorenstein 锥}
\subitem \leavevmode\hypertarget{idx:0155}{齐次自伴锥},\enspace\IndexPageList[0155]{154}
\subitem \leavevmode\hypertarget{idx:0156}{有理多面锥}，\textit{参见} \hyperlink{idx:1205}{有理多面锥}
\indexspace
\item \leavevmode\hypertarget{idx:0157}{锥猜想},\enspace\IndexPageList[0157]{153,157,173,175,268,422}
\subitem \leavevmode\hypertarget{idx:0158}{双有理锥猜想},\enspace\IndexPageList[0158]{156}
\subitem \leavevmode\hypertarget{idx:0159}{推论},\enspace\IndexPageList[0159]{153}
\subitem \leavevmode\hypertarget{idx:0160}{已获证明的情形},\enspace\IndexPageList[0160]{153}
\indexspace
\item \leavevmode\hypertarget{idx:0296}{锥的面},\enspace\IndexPageList[0296]{54}
\indexspace
\item \leavevmode\hypertarget{idx:0168}{锥形奇点},\enspace\IndexPageList[0168]{145,166}
\subitem \leavevmode\hypertarget{idx:0169}{锥形奇点跃迁},\enspace\IndexPageList[0169]{166,427,428,429}，\textit{另见} \hyperlink{idx:0292}{极端跃迁，I 型}
\subsubitem \leavevmode\hypertarget{idx:0170}{例题},\enspace\IndexPageList[0170]{167}
\indexspace
\item \leavevmode\hypertarget{idx:1290}{子轨形},\enspace\IndexPageList[1290]{78,432}
\indexspace
\item \leavevmode\hypertarget{idx:1213}{自反 Gorenstein 锥},\enspace\IndexPageList[1213]{85,87}
\subitem \leavevmode\hypertarget{idx:1214}{对偶性},\enspace\IndexPageList[1214]{87}
\subitem \leavevmode\hypertarget{idx:1215}{指数},\enspace\IndexPageList[1215]{87}
\subitem \leavevmode\hypertarget{idx:1216}{与 nef 分拆的关系},\enspace\IndexPageList[1216]{87}
\subitem \leavevmode\hypertarget{idx:1217}{与自反多胞形的关系},\enspace\IndexPageList[1217]{87}
\indexspace
\item \leavevmode\hypertarget{idx:1224}{自反层},\enspace\IndexPageList[1224]{60}
\indexspace
\item \leavevmode\hypertarget{idx:1218}{自反多胞形},\enspace\IndexPageList[1218]{29,58,70,76,77,78,86,119,139,141,157,158,167,168,178,290,405,406,428}
\subitem \leavevmode\hypertarget{idx:1219}{Batyrev 对偶},\enspace\IndexPageList[1219]{69,79}
\subitem \leavevmode\hypertarget{idx:1220}{分类},\enspace\IndexPageList[1220]{70}
\subitem \leavevmode\hypertarget{idx:1221}{定义},\enspace\IndexPageList[1221]{69}
\subitem \leavevmode\hypertarget{idx:1222}{例题},\enspace\IndexPageList[1222]{72,74,94,111,147,148,162,169}
\subitem \leavevmode\hypertarget{idx:1223}{五次镜像的自反多胞形},\enspace\IndexPageList[1223]{84}
\indexspace
\item \leavevmode\hypertarget{idx:0130}{组合对偶性},\enspace\IndexPageList[0130]{65,69}
\indexspace
\item \leavevmode\hypertarget{idx:0762}{最大退化},\enspace\IndexPageList[0762]{144}，\textit{另见} \hyperlink{idx:0766}{最大幂幺，最大幂幺边界点}
\indexspace
\item \leavevmode\hypertarget{idx:0765}{最大幂幺}
\subitem \leavevmode\hypertarget{idx:0766}{最大幂幺边界点},\enspace\IndexPageList[0766]{105,116,125,127,128,132,133,136,137,138,142,144,151,154,161,172,174,175,176,177,178,179,182,183,184,274,280,289,290,291,292}
\subsubitem \leavevmode\hypertarget{idx:0767}{定义},\enspace\IndexPageList[0767]{104}
\subsubitem \leavevmode\hypertarget{idx:0768}{等价性},\enspace\IndexPageList[0768]{138,143}
\subsubitem \leavevmode\hypertarget{idx:0769}{例题},\enspace\IndexPageList[0769]{112,113,148,150,165}
\subsubitem \leavevmode\hypertarget{idx:0770}{存在性},\enspace\IndexPageList[0770]{139,143,422}
\subsubitem \leavevmode\hypertarget{idx:0771}{该点处的切空间},\enspace\IndexPageList[0771]{181}
\subitem \leavevmode\hypertarget{idx:0772}{最大幂幺单值性},\enspace\IndexPageList[0772]{41,42,101,102,103,104,105,132,136,150,176,186,287}
\subsubitem \leavevmode\hypertarget{idx:0773}{一维模空间},\enspace\IndexPageList[0773]{101,102,103,104}
\subsubitem \leavevmode\hypertarget{idx:0774}{$r$ 维模空间},\enspace\IndexPageList[0774]{104,105}
\indexspace
\item \leavevmode\hypertarget{idx:0711}{左行}
\subitem \leavevmode\hypertarget{idx:0712}{左行费米子},\enspace\IndexPageList[0712]{25,444}
\subitem \leavevmode\hypertarget{idx:0713}{左行扇区},\enspace\IndexPageList[0713]{448}
\indexspace
\item \leavevmode\hypertarget{idx:0033}{作用量},\enspace\IndexPageList[0033]{435,439}
\subitem \leavevmode\hypertarget{idx:0034}{$\sigma$ 模型的作用量},\enspace\IndexPageList[0034]{443,444}
\subsubitem \leavevmode\hypertarget{idx:0035}{玻色部分},\enspace\IndexPageList[0035]{443}
\subsubitem \leavevmode\hypertarget{idx:0036}{费米部分},\enspace\IndexPageList[0036]{443}
\end{theindex}
\endgroup
